% modulo fr''uher eunf''uhren, ABC, Elementarteilersatz,
% Euklidischer Algorithmus, Moduln, Polynome in beliebiger Variablenmenge
% (wegen Beweis der Existenz eines algebraischen Abschluss).
% Kreisteilungspolynome, Endliche K''orper, Wedderburn
\documentclass[12pt,a4paper,twoside]{article}
\usepackage[cmtip,arrow]{xy}
\usepackage{amssymb,amsthm,amsmath,ngerman,makeidx,color,pb-diagram,pb-xy,verbatim}
\usepackage[dvips]{graphicx}
%\usepackage{graphicx}
\makeindex
\input{latams}
\newcommand{\ei}[1]{\emph{#1}\index{#1}}
\DeclareMathOperator{\ord}{ord}
\DeclareMathOperator{\ggT}{ggT}
\DeclareMathOperator{\grad}{grad}
\DeclareMathOperator{\Fix}{Fix}
\begin{document}
\title{Algebra I+II}
\author{Vorlesungen Wintersemester 2004/05 und Sommersemester 2005\\
Peter M\"uller}
\maketitle
\tableofcontents
\section{Einf"uhrung}
Die Algebra hat ihre Wurzeln in der Geometrie und Zahlentheorie.
\paragraph{Klassische Beispiele}
\begin{itemize}
\item Kann man mit Zirkel und Lineal beliebige Winkel dritteln, oder
  W\"urfel verdoppeln?
\item Kann man die Nullstellen von Polynomen stets durch Wurzeln und die
  vier Grundrechenarten ausdr\"ucken?
\item Kann man f\"ur Funktionen wie $e^{-x^2}$ explizite Stammfunktionen
  angeben?
\item Wie findet man rationale oder ganzzahlige L\"osungen von Systemen
  von Polynomen in mehreren Ver\"anderlichen, z.B.\ (f\"ur festes $n\in\NNN$)
  $X^n+Y^n=Z^n$ mit $X,Y,Z\in\ZZZ$ (Fermat-Problem)?
\end{itemize}
\paragraph{Moderne Beispiele}
\begin{itemize}
\item Wie \"ubertr\"agt man digitale Daten, bei deren \"Ubertragung Fehler
  auftreten k\"onnen, durch Einbau von m\"oglichst wenig Redundanz, aber mit
  m\"oglichst guter Fehlererkennungs- und -korrekturm\"oglichkeit? Dies
  f\"uhrt zur Kodierungstheorie.
\item Wie \"ubertr\"agt man verschl\"usselte Daten \"offentlich, so dass nur
  der vorgesehene Empf\"anger sie entschl\"usseln kann? Das geht sogar, ohne
  dass Sender und Empf\"anger vorher geheime Schl\"ussel austauschen m\"ussen,
  selbst die d\"urfen \"offentlich mitgeteilt werden (public key
  cryptography)! Dies f\"uhrt zur Kryptographie.
\end{itemize}
In der Algebra haben sich wichtige Begriffe herauskristallisiert, n\"amlich
Gruppen, Ringe und K\"orper. Diesen Begriffen ist die Vorlesung Algebra I
gewidmet.

Die K\"orpertheorie wir in der Algebra II fortgesetzt. Ein Schwerpunkt wird
die Galoistheorie sein, eine reizvolle Kombination aus Gruppen- und
K\"orpertheorie.
\section{Gruppen}
\subsection{Definitionen, Beispiele}
\begin{Definition}
  Eine Menge $G$ mit einer zweistelligen Verkn\"upfung $G\times G\to G$,
  $(x,y)\mapsto x\cdot y$, hei\ss t \ei{Halbgruppe}, wenn das
  \ei{Assoziativgesetz} gilt, d.h.\ $(x\cdot y)\cdot z=x\cdot(y\cdot z)$ f\"ur
  alle $x,y,z\in G$. H\"aufig schreibt man $xy$ statt $x\cdot y$.
\end{Definition}
\paragraph{Beispiele von Halbgruppen}
\begin{itemize}
\item $(\NNN,+)$, $(\NNN,\cdot)$, $(\ZZZ,+)$, $(\ZZZ,\cdot)$.
\item $\text{Abb}(X,X)=$ Menge der Abbildungen von $X$ auf sich. Das Produkt
  $f\cdot g$ f\"ur $f,g\in\text{Abb}(X,X)$ ist als die Komposition definiert,
  also $(f\cdot g)(x)=f(g(x))$ f\"ur $x\in X$. In der Gruppentheorie werden
  wir h\"aufig $x^f$ statt $f(x)$ schreiben.
\item Sei $G$ die Potenzmenge einer Menge. Dann sind $(G,\cap)$ und $(G,\cup)$
  Halbgruppen.
\end{itemize}
\begin{Definition}
  Eine Halbgruppe hei\ss t \ei{abelsch} oder \ei{kommutativ}, wenn $xy=yx$
  f\"ur alle $x,y\in G$ gilt.
\end{Definition}
\begin{Definition}
  Die Halbgruppe $G$ ist ein \ei{Monoid}, wenn es ein \ei{Einselement} $e$
  (oder auch \ei{neutrales Element} genannt) gibt, d.h.\ es gilt $ex=xe=x$
  f\"ur alle $x\in G$. Man schreibt dann $x^0=e$.
\end{Definition}
\begin{Beispiel*}
  Die Halbgruppe $(\NNN,+)$ ist kein Monoid, $(\NNN,\cdot)$ aber schon.
\end{Beispiel*}
Wir kommen nun zum zentralen Begriff dieses Kapitels.
\begin{Definition}
  Ein Monoid $G$ hei\ss t \ei{Gruppe}, wenn es f\"ur alle $x\in G$ ein $y\in
  G$ gibt mit $xy=yx=e$.
\end{Definition}
\begin{Bemerkung*}
  In der Definition ist $y$ eindeutig. Denn aus $y'x=xy'=e$ folgt
  $y'=y'e=y'(xy)=(y'x)y=ey=y$.

Man schreibt $x^{-1}$ f\"ur $y$, und $x^{-n}:=(x^{-1})^n$ f\"ur
$n\in\NNN$. Mit diesen Festsetzungen gilt dann $x^mx^n=x^{m+n}$ und
$(x^m)^n=x^{mn}$ f\"ur alle $m,n\in\ZZZ$.
\end{Bemerkung*}
\paragraph{Einfache Folgerungen f\"ur eine Gruppe $G$}
\begin{itemize}
\item Seien $a,b\in G$. Dann sind die Gleichungen $ax=b$ und $ya=b$ eindeutig
  durch $x=a^{-1}b$ und $y=ba^{-1}$ l\"osbar. Im allgemeinen gilt $x\ne y$, so
  dass eine Bruchschreibweise $\frac{b}{a}$ nur f\"ur abelsche Gruppen
  sinnvoll ist.
\item Es gelten K\"urzungsregeln: Aus $ax=ay$ (oder $xa=ya$) folgt $x=y$.
\item $(ab)^{-1}=b^{-1}a^{-1}$.
\end{itemize}
\paragraph{Beispiele von Gruppen}
\begin{itemize}
\item $(V,+)$ mit Vektorraum $V$.
\item $(K\setminus\{0\},\cdot)$ mit K\"orper $K$.
\item $S_X=$ Teilmenge der Bijektionen aus $\text{Abb}(X,X)$, wieder mit
  Komposition als Verkn\"upfung. F\"ur $X=\{1,2,\dots,n\}$ schreibt man $S_n$
  statt $S_X$. Man nennt $S_n$ die \ei{symmetrische Gruppe} vom Grad $n$.
\item $\GL_n(K)=\{A\in M_n(K)|\;\det A\ne 0\}$, $K$ K\"orper.
\item $\SL_n(\ZZZ)=\{A\in M_n(\ZZZ)|\;\det A=1\}$.
\end{itemize}
\begin{Bemerkung*}
  H\"aufig schreibt man $1$ f\"ur das neutrale Element $e$. Ist die Gruppe
  abelsch, so verwendet man oft eine additive Notation, also $a+b$ und $-a$
  statt $a\cdot b$ und $a^{-1}$. Nat\"urlich verwendet man dennoch z.B.\ f\"ur
  multiplikative Gruppen von K\"orpern die multiplikative Notation.
\end{Bemerkung*}
\subsection{Untergruppen und zyklische Gruppen}
\begin{Definition}
  Eine nichtleere Teilmenge $U$ einer Gruppe $G$ hei\ss t \ei{Untergruppe} von
  $G$, wenn aus $a,b\in U$ schon $a\cdot b\in U$ und $a^{-1}\in U$ folgt. (In
  diesem Fall ist $(U,\cdot)$ eine Gruppe.) Man schreibt $U\le G$.
\end{Definition}
Es folgt sehr einfach, dass Untergruppen von Untergruppen wieder Untergruppen
sind. Auch Schnitte von Untergruppen sind wieder Untergruppen. (Was ist mit
Vereinigungen?)
\paragraph{Beispiele}
\begin{itemize}
\item $G=\ZZZ$, $U=2\ZZZ=\{2m|\;m\in\ZZZ\}$
\item $G=S_5$, $U=\{g\in G|\;g(5)=5\}=S_4$
\item $G=\GL_n(K)$, $U=\SL_n(K)$
\end{itemize}
\begin{Definition}
  Sei $M\subseteq G$ eine Teilmenge der Gruppe $G$. Die von $M$ \ei{erzeugte
  Untergruppe} $\gen{M}$ ist die kleinste $M$ enthaltende Untergruppe von $G$.
\end{Definition}
Wichtige Beispiele von (Unter)gruppen sind solche, die von einem Element
erzeugt werden.
\begin{Definition}
  Eine \ei{zyklische Gruppe} ist eine von einem Element erzeugte Gruppe.
\end{Definition}
\paragraph{Beispiele zyklischer Gruppen}
\begin{itemize}
\item $G=\ZZZ=\gen{1}$
\item $G=\{e^{\frac{2\pi ik}{n}}|\;k=0,1,\dots,n-1\}=\gen{e^{\frac{2\pi
        i}{n}}}$ f\"ur $n\in\NNN$
\item $G=\ZZZ/n\ZZZ=\gen{1+n\ZZZ}$ f\"ur $n\in\NNN$
\end{itemize}
\begin{Definition}
Ist $G$ eine Gruppe, dann hei\ss t $\abs{G}\in\NNN\cup\{\infty\}$ die
\ei{Ordnung} von $G$.

F\"ur $a\in G$ hei\ss t $\ord(a)=\abs{\gen{a}}\in\NNN\cup\{\infty\}$ die
\ei{Ordnung} von $a$.
\end{Definition}
\begin{Lemma}
  Ist $\ord(a)=n\in\NNN$, dann ist $n$ die kleinste nat\"urliche Zahl $m$ mit
  $a^m=e$. Ferner gilt $\gen{a}=\{e=a^0,a^1,\dots,a^{n-1}\}$, und $a^i=a^j$
  ist \"aquivalent zu $i\equiv j\pmod{n}$.
\end{Lemma}
\begin{proof}[Beweis]
Zun\"achst m\"ussen wir sehen, dass es \"uberhaupt eine nat\"urliche Zahl $m$
  gibt $a^m=e$. Das folgende Argument unter Verwendung des Schubfachprinzips
  wird uns in vielen Varianten noch \"ofters begegnen: Unter den $n+1$
  Elementen $\{a^0,a^1,\dots,a^n\}$ muss es zwei gleiche geben, da die Menge
  eine Teilmenge einer $n$-Elementigen Menge ist. Sei etwa $a^i=a^j$ mit $0\le
  i<j\le n$. Dann gilt $j-i\in\NNN$ und $a^{j-i}=e$.
  
  Sei nun $m\in\NNN$ minimal mit $a^m=e$. Dann gilt
  $\gen{a}=\{a^0,a^1,\dots,a^{m-1}\}$. Ferner gilt $a^i\ne a^j$ f\"ur $0\le
  i<j\le m-1$, denn andernfalls folgte $a^{j-i}=e$, im Widerspruch zur
  Minimalit\"at von $m$. Es folgt $m=n$ und der Rest der Behauptung.
\end{proof}
\begin{Definition}
  Ein Element der Ordnung $2$ nennt man \ei{Involution}.
\end{Definition}
Wir werden sp\"ater sehen, dass die Struktur der Untergruppen selbst in
endlichen Gruppen sehr kompliziert sein kann. \"Ubersichtlich ist das hingegen
in zyklischen Gruppen.
\begin{Satz}
  Die Gruppe $G$ sei zyklisch. Dann ist jede Untergruppe von $G$ zyklisch.
  Hat $G$ endliche Ordnung $n$, dann hat $G$ f\"ur jeden Teiler $r$ von $n$
  genau eine Untergruppe der Ordnung $r$. (Ferner ist die Ordnung jeder
  Untergruppe ein Teiler von $n$.)
\end{Satz}
\begin{proof}[Beweis]
  Sei $U$ eine Untergruppe von $G=\gen{a}$. W\"ahle $s\in\NNN$ minimal mit
  $a^s\in U$. Wir behaupten $U=\gen{a^s}$. Dazu sei $m\in\ZZZ$ mit $a^m\in
  U$. Schreibe $m=us+v$ mit $0\le v\le s-1$ (Division mit Rest). Wegen
  $(a^s)^ua^v=a^m\in U$ und $a^s\in U$ folgt $a^v\in U$, also $v=0$ wegen der
  Minimalit\"at von $s$. Aber $a^m=a^{us}\in\gen{a^s}$. Hieraus folgt der
  erste Teil der Behauptung.

Sei nun $n\in\NNN$ die Ordnung von $a$. Wegen $a^n=e\in U$ folgt wie eben,
dass $s$ ein Teiler von $n$ ist, also $n=rs$ mit $r\in\NNN$. Offenbar gilt
$\gen{a^s}=\{a^0,a^s,a^{2s},\dots,a^{(r-1)s}\}$, also
$\abs{U}=\ord(a^s)=r$. Umgekehrt ist f\"ur jeden Teiler $r$ von $n$ die Menge
$\gen{a^s}=\{a^0,a^s,a^{2s},\dots,a^{(r-1)s}\}$ mit $s=n/r$ eine Untergruppe
der Ordnung $r$.
\end{proof}
\subsection{Nebenklassen}\index{Nebenklassen}
Sei $U$ eine Untergruppe der Gruppe $G$. F\"ur $g\in G$ nennt man die Menge
$Ug:=\{ug|\;u\in U\}$ eine \ei{Rechtsnebenklasse}. Analog ist $gU$ eine
\ei{Linksnebenklasse}. Die Menge der Rechts- bzw.\ Linksnebenklassen
bezeichnen wir mit $U\backslash G$ bzw.\ $G/U$.

In diesem Abschnitt verstehen wir unter Nebenklassen, wenn nichts anderes
gesagt wird, Rechtsnebenklassen. Mit entsprechenden Modifikationen gelten die
folgenden Aussagen auch f\"ur Linksnebenklassen.

\begin{Lemma}
  F\"ur zwei Nebenklassen $Ug$ und $Uh$ gilt entweder $Ug=Uh$, oder $Ug\cap
  Uh=\emptyset$.
\end{Lemma}
\begin{proof}[Beweis]
  Sei $x\in Ug\cap Uh$, also $x=u_1g=u_2h$ f\"ur $u_1,u_2\in U$. Es folgt
  $g=u_3h$ mit $u_3=u_1^{-1}u_2\in U$, also $Ug=Uu_3h=Uh$.
\end{proof}
Da die Inversenabbildung $g\mapsto g^{-1}$ eine Bijektion von $G$ ist, ist
auch die Abbildung $Ug\mapsto (Ug)^{-1}=g^{-1}U^{-1}=g^{-1}U$ eine Bijektion
zwischen $U\backslash G$ und $G/U$. Insbesondere gilt
$\abs{G/U}=\abs{U\backslash G}$. Man nennt $[G:U]:=\abs{U\backslash G}$ den
\ei{Index} von $U$ in $G$.
\begin{Definition}
  Aus obigem Lemma gewinnt man eine disjunkte Zerlegung $G=\biguplus_i Ug_i$
  f\"ur geeignete $g_i\in G$, die \ei{Nebenklassenzerlegung}. Die Menge der
  $g_i$ nennt man \ei{Rechtstransversale} oder \ei{Vertretersystem} von
  $U\backslash G$.
\end{Definition}
\begin{Bemerkung*}
  Eine Rechtstransversale ist im allgemeinen keine Linkstransversale. Dennoch
  kann man z.B.\ f\"ur endliche Gruppen $G$ zeigen, dass es zu jeder
  Untergruppe $U$ eine simultane Rechts- und Linkstransversale gibt.
\end{Bemerkung*}
F\"ur alle $g\in G$ gilt $\abs{Ug}=\abs{U}$; zusammen mit der
Nebenklassenzerlegung folgt der fundamentale Satz von Lagrange:
\begin{Satz}[Lagrange]
  Sei $U$ eine Untergruppe der Gruppe $G$. Dann gilt $\abs{G}=[G:U]\abs{U}$.
  Insbesondere ist (falls $\abs{G}<\infty$) $\abs{U}$ ein Teiler von
  $\abs{G}$.
\end{Satz}
\begin{Bemerkung*}
  Die Umkehrung des Satzes von Lagrange gilt im allgemeinen nicht, es muss
  nicht f\"ur jeden Teiler $r$ von $\abs{G}$ eine Untergruppe der Ordnung $r$
  existieren. Ist hingegen $r$ eine Primpotenz, so werden wir sp\"ater sehen,
  dass es stets Untergruppen der Ordnung $r$ gibt.
\end{Bemerkung*}
Der Satz von Lagrange hat einige bemerkenswerte Folgen:
\begin{Korollar}
  Sei $G$ eine endliche Gruppe. F\"ur alle $g\in G$ ist $\ord(g)$ ein Teiler
  von $\abs{G}$.
\end{Korollar}
\begin{proof}[Beweis]
  Klar, da $\gen{g}$ eine Untergruppe der Ordnung $\ord(g)$ ist.
\end{proof}
\begin{Korollar}[Kleiner Satz von Fermat]
  $a\in\ZZZ$ sei nicht durch die Primzahl $p$ teilbar. Dann ist $a^{p-1}-1$
  durch $p$ teilbar.
\end{Korollar}
\begin{proof}[Beweis]
  Die Gruppe der invertierbaren Elemente im Monoid $(\ZZZ/p\ZZZ,\cdot)$ hat
  Ordnung $p-1$.
\end{proof}
\begin{Korollar}
  Gruppen von Primzahlordnung sind zyklisch.
\end{Korollar}
\begin{proof}[Beweis]
  Sei $e\ne g\in G$. Dann ist $\ord(g)>1$ ein Teiler der Primzahl $\abs{G}$,
  also $\ord(g)=\abs{G}$ und daher $G=\gen{g}$.
\end{proof}
\subsection{Normalteiler und Faktorgruppen}
Sind $A$ und $B$ Teilmengen einer Halbgruppe $G$, dann schreiben wir
$AB:=\{ab|\;a\in A,b\in B\}$.

Auf den Begriff des Normalteilers st\"o\ss t man, wenn man versucht, auf dem
Quotientenraum $U\backslash G$ auf sinnvolle Weise ein Produkt einzuf\"uhren. Es sollte
$(Ug)\cdot(Uh)=Ugh$ f\"ur alle $g,h\in G$ gelten. Speziell f\"ur $h=1$ folgt
$gU\subseteq (Ug)U=Ug$, und daraus (durch Inversenbildung und Ersetzen von $g$
durch $g^{-1}$) $gU=Ug$. Diese spezielle Eigenschaft von Untergruppen
kennzeichnet Normalteiler.
\begin{Definition}
  Eine Untergruppe $N\le G$ hei\ss t \ei{Normalteiler} von $G$, wenn $Ng=gN$
  f\"ur alle $g\in G$ gilt. Man sagt auch, $N$ ist \ei{normal} in $G$, und
  schreibt $N\trianglelefteq G$.
\end{Definition}
Zur \"Uberpr\"ufung der Normalteilereigenschaft ist folgende Aussage oft
n\"utzlich.
\begin{Lemma}\label{L:normalkonj}
  Die Untergruppe $N$ von $G$ ist genau dann ein Normalteiler, wenn
  $g^{-1}Ng\subseteq N$ gilt f\"ur alle $g\in G$.
\end{Lemma}
\begin{proof}[Beweis]
  Aus $g^{-1}Ng\subseteq N$ folgt $Ng\subseteq gN$. Mit $g^{-1}$ statt $g$
  gilt $Ng^{-1}\subseteq g^{-1}N$, also $gN=(Ng^{-1})^{-1}\subseteq
  (g^{-1}N)^{-1}=Ng$, und die eine Richtung der Behauptung folgt. Die andere
  Richtung ist sowieso klar.
\end{proof}
\paragraph{Beispiele}
\begin{itemize}
\item Untergruppen abelscher Gruppen sind normal.
\item Sei $U\le G$ mit $[G:U]=2$. F\"ur $g\in G$ gilt entweder $g\in U$, und
  daher $Ug=U=gU$, oder $g\notin U$, und dann ist $Ug=gU$, da $G$ die disjunkte
  Vereinigung von $U$ und $Ug$ bzw.\ von $U$ und $gU$ ist. Daher sind
  Untergruppen vom Index $2$ stets Normalteiler.
\item Es gilt $\SL_n(K)\trianglelefteq\GL_n(K)$. Dazu seien $A\in\SL_n(K)$ und
  $B\in\GL_n(K)$ beliebig. Wegen $\det(B^{-1}AB)=\det(B)^{-1}\det(A)\det(B)=1$
  gilt $B^{-1}AB\in\SL_n(K)$, und die Behauptung folgt aus Lemma
  \ref{L:normalkonj}.
\item Die Untergruppe $U$ der oberen Dreiecksmatrizen in $\GL_n(K)$ ist f\"ur
  $n\ge2$ nicht normal in $\GL_n(K)$. Dazu sei $D$ die Antidiagonalmatrix mit
  Eintr\"agen $1$. Dann besteht $D^{-1}UD$ aus den unteren Dreiecksmatrizen.
\end{itemize}
Der folgende Satz bildet aus einer Gruppe $G$ und einem Normalteiler $N$ eine
kleinere Gruppe $G/N$. Dieses Prinzip ist wichtig, um \"uber eine eventuell
komplizierte Gruppe $G$ Informationen aus den kleineren Gruppen $N$ und $G/N$
gewinnen zu k\"onnen.
\begin{Satz}
  Sei $N\trianglelefteq G$. Dann gilt $(Ng)(Nh)=Ngh$ f\"ur alle $g,h\in
  G$. Mit diesem Produkt wird $G/N$ zu einer Gruppe, der \ei{Faktorgruppe} von
  $G$ modulo $N$.
\end{Satz}
\begin{proof}[Beweis]
  Durch Verwendung von $Nx=xN$ und der Assoziativit\"at folgt $NgNh=NNgh=Ngh$.
  Wir m\"ussen zeigen, dass wir ein wohldefiniertes Produkt auf $G/N$
  bekommen. Dazu sei $Ng=Ng'$ und $Nh=Nh'$. Dies liefert $g'=ug$ und $h'=vh$
  mit $u,v\in N$. Wegen $gvg^{-1}=w\in N$ und $uw\in N$ gilt
  $g'h'=ugvh=u(gvg^{-1})gh=uwgh$ und daher $Ng'h'=Ngh$.

  Assoziativit\"at des Produkt auf $G/N$ ist klar, das neutrale Element ist
  $N$, und das inverse Element von $Ng$ ist $Ng^{-1}$.
\end{proof}
\begin{Definition}
  Eine Gruppe $G>\{e\}$ hei\ss t \ei{einfach}, wenn sie au\ss er $G$ und
  $\{e\}$ keine Normalteiler hat.
\end{Definition}
Einfache Gruppen spielen eine \"ahnliche Rolle wie die Primzahlen f\"ur die
nat\"urlichen Zahlen. Nach dem Satz von Lagrange ist z.B.\ eine Gruppe von
Primzahlordnung einfach. Weniger triviale Beispiele werden uns in den
folgenden Abschnitten begegnen.
\subsection{Symmetrische und alternierende Gruppen}
Sei $n\in\NNN$ und $S_n$ die Gruppe der Permutationen von $\{1,2,\dots,n\}$.
Wir verwenden die Exponentialnotation $a^\phi$ f\"ur das Bild von $a$ unter
$\phi$. Das hat den Vorteil der bequemen Notation wie z.B.\ 
$a^{\phi\psi}=(a^\phi)^\psi$ f\"ur $\phi,\psi\in S_n$. Man kann Permutationen
$\phi\in S_n$ durch eine Wertetabelle $\phi=
\begin{pmatrix}
  1 & 2 & \cdots & n\\
1^\phi & 2^\phi & \cdots & n^\phi
\end{pmatrix}
$ angeben. Transparenter, und f\"ur praktische wie auch theoretische Zwecke
geeigneter, ist die Zykelnotation. Dabei hei\ss t $\psi\in S_n$ ein
\ei{$m$-Zykel}, wenn es eine Folge $a_1,a_2,\dots,a_m$ verschiedener Zahlen
aus $\{1,2,\dots,n\}$ gibt mit $a_i^\phi=a_{i+1}$ f\"ur $1\le i\le m-1$,
$a_m^\phi=a_1$ und $a^\phi=a$ f\"ur alle $a\notin\{a_1,a_2,\dots,a_m\}$. Die
\ei{Zykelnotation} f\"ur $\psi$ ist $\psi=(a_1,a_2,\dots,a_m)$, wobei man auch
h\"aufig $\psi=(a_1\, a_2\,\ldots\, a_m)$ schreibt. In den folgenden Beweisen
erspart man sich Fallunterscheidungen, wenn man $a_{m+1}=a_1, a_{m+2}=a_2,
\dots$ setzt.

Ist $\phi\in S_n$ eine Permutation, so nennt man die Menge der von $\phi$
bewegten Elemente den Tr\"ager von $\phi$; er besteht also aus den $a$ mit
$a^\phi\ne a$.

Jedes Element $\phi$ aus $S_n$ ist ein Produkt von Zykeln $\psi_i$ mit
disjunkten Tr\"agern. Diese Zykel kommutieren paarweise, also
$\psi_i\psi_j=\psi_j\psi_i$ f\"ur alle $i,j$. Das Produkt diser Zykel in
Zykelnotation ist die \ei{Zykelnotation} f\"ur $\phi$. Kommt in $\phi$ ein
trivialer Zykel $(a)$ vor, so l\"asst man den Faktor $(a)$ h\"aufig weg.
\paragraph{Beispiele}
\begin{itemize}
\item $n=3$, $\phi=\begin{pmatrix}
  1 & 2 & 3\\
3 & 2 & 1
\end{pmatrix}
=(1 3)(2)=(1 3)$.
\item $n=5$, $\phi=\begin{pmatrix}
  1 & 2 & 3 & 4 & 5\\
3 & 5 & 1 & 2 & 4
\end{pmatrix}
=(1 3)(2 5 4)$, $\phi^{-1}=(1 3)(2 4 5)$, $\phi^2=(1)(3)(2 4 5)=(2 4 5)$.
\item $n=5$, $\phi=(1 2 3 4)$, $\psi=(2 3)(4 5)$, $\phi\psi=(1 3 5 4)$.
\end{itemize}
Ein wichtiger Begriff der Gruppentheorie ist die \ei{Konjugiertheit}, den wir
hier am Beispiel der symmetrischen Gruppe n\"aher kennenlernen wollen. Die
Konjugiertheit ist eine Verallgemeinerung des Begriffs der \"Ahnlichkeit f\"ur
die Gruppe der invertierbaren Matrizen $\GL_n(K)$.
\begin{Definition}
  Sei $G$ eine Gruppe. Die Elemente $a,b\in G$ hei\ss en \ei{konjugiert}, wenn
  es ein $g\in G$ gibt mit $b=g^{-1}ag$. Statt $g^{-1}ag$ schreibt man auch
  $a^g$, was die eing\"angigen Beziehungen $a^{gh}=(a^g)^h$ und
  $(ab)^g=a^gb^g$ liefert.
\end{Definition}
Die Konjugiertheit schreibt sich in Zykelnotation sehr \"ubersichtlich:
\begin{Satz}
  Sei $\phi=(a_1\,a_2\dots)(b_1\,b_2\dots)\dots\in S_n$ und $\psi\in S_n$.
  Dann gilt
  $\phi^\psi=(a_1^\psi\,a_2^\psi\dots)(b_1^\psi\,b_2^\psi\dots)\dots$.
\end{Satz}
\begin{proof}[Beweis]
  Wir m\"ussen sehen, wohin z.B.\ das Element $a_i^\psi$ von $\phi^\psi$
  abgebildet wird. Wir rechnen
  \begin{align*}
    (a_i^\psi)^{\phi^\psi} &= (a_i^\psi)^{\psi^{-1}\phi\psi}\\
&= a_i^{\psi\psi^{-1}\phi\psi} \\
&= a_i^{\phi\psi} \\
&= (a_i^\phi)^\psi \\
&= a_{i+1}^\psi
  \end{align*}
und die Behauptung folgt.
\end{proof}
Die Zykell\"angen von $\alpha$ sind die L\"angen der in der Zykelschreibweise
von $\alpha$ auftretenden Zykel. Obiger Satz zeigt, dass sich die Zykell\"angen
unter Konjugation nicht \"andern. Es gilt auch die Umkehrung:
\begin{Satz}\label{S:zykkon}
  Seien $\alpha,\beta\in S_n$ Permutationen mit den gleichen Zykell\"angen.
  Dann sind $\alpha$ und $\beta$ in $S_n$ konjugiert.
\end{Satz}
\begin{proof}[Beweis]
  Seien $m_1,m_2,\dots$ die Zykell\"angen. Schreibe
  \begin{align*}
    \alpha &=
    (a_{1,1}\,a_{1,2}\dots a_{1,m_1})(a_{2,1}\,a_{2,2}\dots a_{2,m_2})\dots\\
    \beta &=
    (b_{1,1}\,b_{1,2}\dots b_{1,m_1})(b_{2,1}\,b_{2,2}\dots b_{2,m_2})\dots,\\
\end{align*}
wobei wir auch die Zykel der L\"angen $1$ mit aufschreiben. Jedes Element aus
$\{1,2,\dots,n\}$ kommt daher genau einmal sowohl unter den $a_{i,j}$, als
auch unter den $b_{i,j}$, vor. Sei $\psi\in S_n$ definiert durch
$b_{i,j}=a_{i,j}^\psi$. Aus der folgenden Rechnung ergibt sich
$\beta=\alpha^\psi$:
\begin{align*}
  b_{i,j}^{\beta} &= b_{i,j+1}\\
&= a_{i,j+1}^\psi\\
&= a_{i,j}^{\alpha\psi}\\
&= a_{i,j}^{\psi\psi^{-1}\alpha\psi}\\
&= b_{i,j}^{\psi^{-1}\alpha\psi}
\end{align*}
\end{proof}
Eine wichtige Klasse von Permutationen sind die $2$-Zykel. Man nennt sie auch
\ei{Transpositionen}. Die folgende Aussage zeigt, dass die symmetrischen
Gruppen von Transpositionen erzeugt werden.
\begin{Satz}
  Jede Permutation $\phi\in S_n$ ist ein Produkt von
  Transpositionen, und jeder Zykel der L\"ange $m$ ist ein Produkt von $m-1$
  Transpositionen.
\end{Satz}
\begin{proof}[Beweis]
  Da jede Permutation ein Produkt von Zykeln ist, gen\"ugt es, die Aussage
  f\"ur einen Zykel zu beweisen. Die Behauptung folgt induktiv wegen
  $(a_1\,a_2\dots a_m)=(a_2\,a_3\dots a_m)(a_1\,a_2)$.
\end{proof}
Sind $m_1,m_2,\dots,m_r$ die Zykell\"angen eines Elements $\phi\in S_n$, so
nennt man $\ell(\phi):=\sum_{i=1}^r(m_i-1)$ die \ei{L\"ange} von
$\phi$. Obiger Satz zeigt, dass sich $\phi$ als Produkt von $\ell(\phi)$
Transpositionen schreiben l\"asst.
\begin{Satz}
  F\"ur $\phi,\psi\in S_n$ gilt
  $\ell(\phi\psi)\equiv\ell(\phi)+\ell(\psi)\pmod{2}$.
\end{Satz}
\begin{proof}[Beweis]
  Schreibt man $\phi$ und $\psi$ jeweils als Produkt von Transpositionen, so
  sieht man, dass die Behauptung aus folgender Aussage folgt: Sei
  $\alpha=\tau_1\tau_2\dots\tau_m$ ein Produkt von $m$ Transpositionen. Dann
  gilt $\ell(\alpha)\equiv m\pmod{2}$. Das beweisen wir durch Induktion \"uber
  $m$. Die Behauptung ist trivial f\"ur $m=1$. Wir nehmen nun an, dass sie
  richtig ist f\"ur $m$. Sei $\tau$ eine Transposition. Die Richtigkeit der
  Aussage f\"ur $m+1$ folgt, wenn wir zeigen, dass
  $\ell(\alpha\tau)=\ell(\alpha)+\delta$ gilt mit $\delta\in\{-1,1\}$. Wir
  schreiben $\tau=(u\,v)$, und unterscheiden zwei F\"alle:
\begin{itemize}
\item[(i)] $u$ und $v$ sind in einem Zykel $\beta=(a_1\,a_2\cdots a_r)$ von
  $\alpha$ enthalten, also etwa $a_i=u$ und $a_j=v$ mit $i<j$. Zur Bestimmung
  der Zykeldarstellung von $\alpha\tau$ muss man lediglich $\beta\tau$
  bestimmen. Man erh\"alt $\beta\tau=(a_1\,a_2\cdots
  a_r)(a_i\,a_j)=(a_1\,a_2\dots a_{i-1}\,a_j\,a_{j+1}\dots
  a_r)(a_i\,a_{i+1}\dots a_{j-1})$, also
  $\ell(\beta\tau)=(r-j+i-1)+(j-i-1)=r=\ell(\beta)-1$, und daher
  $\ell(\alpha\tau)=\ell(\alpha)-1$
\item[(ii)] $u$ und $v$ liegen in verschiedenen Zykeln $(a_1\dots a_r)$ und
  $(b_1\dots b_s)$ von $\alpha$, wobei wir $a_1=u$ und $b_1=v$ annehmen
  d\"urfen. Aus $(a_1\dots a_r)(b_1\dots b_s)(a_1\,b_1)=(a_1\,a_2\dots
  a_r\,b_1\,b_2\dots b_s)$ folgt $\ell(\alpha\tau)=\ell(\alpha)+1$.
\end{itemize}
\end{proof}
\begin{Definition}
  Permutationen $\phi$ mit $\ell(\phi)\equiv0\pmod{2}$ hei\ss en gerade, und
  solche mit $\ell(\phi)\equiv1\pmod{2}$ hei\ss en ungerade.
\end{Definition}
\begin{Satz}
  Sei $1<n\in\NNN$. Die Menge der geraden Permutationen aus $S_n$ bildet eine
  Gruppe $A_n$ mit $[S_n:A_n]=2$. Man nennt $A_n$ die \ei{alternierende
  Gruppe} vom Grad $n$.
\end{Satz}
\begin{proof}[Beweis]
  Aus obigem Satz folgt die multiplikative Abgeschlossenheit von $A_n$, somit
  ist $A_n$ eine Untergruppe von $S_n$. Sei $\tau\in S_n$ eine Transposition,
  und $\phi\in S_n$, aber $\phi\notin A_n$. Da $\tau$ und $\phi$ ungerade
  Permutationen sind, ist nach dem obigen Satz $\phi\tau$ gerade, also
  $\phi\tau\in A_n$, und daher $\phi\in A_n\tau$. Somit ist $S_n=A_n\cup
  A_n\tau$ eine Nebenklassenzerlegung, und die Behauptung folgt.
\end{proof}
Unser Ziel ist es zu zeigen, dass die alternierenden Gruppen $A_n$ f\"ur
$n\ge5$ einfach sind. Hierf\"ur ben\"otigen wir noch einige Vorbereitungen:
\begin{Satz}\label{S:Anerz}
  Jede gerade Permutation ist ein Produkt von $3$-Zykeln. Insbesondere wird
  $A_n$ von den $3$-Zykeln aus $S_n$ erzeugt.
\end{Satz}
\begin{proof}[Beweis]
  Da jede gerade Permutation das Produkt einer geraden Anzahl von
  Transpositionen ist, gen\"ugt es zu zeigen, dass das Produkt von zwei
  Transpositionen $\sigma=(a\,b)$ und $\tau=(c\,d)$ ein Produkt von $3$-Zykeln
  ist. Wir unterscheiden zwei F\"alle: (i) $\sigma$ und $\tau$ bewegen einen
  gemeinsamen Punkt. Sei also etwa $a=c$ (und $b\ne d$, da es sonst nichts zu
  beweisen gibt). Dann gilt $\sigma\tau=(a\,b)(a\,d)=(a\,b\,d)$. (ii) $\sigma$
  und $\tau$ bewegen keinen gemeinsamen Punkt. In diesem Fall gilt
  $\sigma\tau=(a\,b)(c\,d)=(a\,b)(b\,c)(b\,c)(c\,d)=(a\,c\,b)(b\,d\,c)$.
\end{proof}
\begin{Lemma}
  Sei $n\ge5$. Dann sind alle $3$-Zykel aus $A_n$ konjugiert.
\end{Lemma}
\begin{proof}[Beweis]
  Seien $\alpha$ und $\beta$ zwei $3$-Zykel. Nach Satz \ref{S:zykkon} gibt es
  $\gamma\in S_n$ mit $\beta=\alpha^\gamma$. Liegt $\gamma$ in $A_n$, dann
  sind wir fertig. Das sei also nicht der Fall. Wegen $n\ge5$ gibt es eine
  Transposition $\tau$, die zwei Fixpunkte von $\beta$ vertauscht. Somit gilt
  $\beta=\alpha^\gamma=\alpha^{\gamma\tau}$, und die Behauptung folgt wegen
  $\gamma\tau\in A_n$.
\end{proof}
\begin{Satz}
  F\"ur $n\ge5$ ist die Gruppe $A_n$ einfach.
\end{Satz}
\begin{proof}
  Sei $N>1$ ein Normalteiler von $A_n$. Unser Ziel ist es zu zeigen, dass $N$
  einen $3$-Zykel $\rho$ enth\"alt. F\"ur alle $\gamma\in A_n$ folgt dann
  n\"amlich $\rho^\gamma\in N^\gamma=N$, und wegen obigem Lemma enth\"alt $N$
  dann alle $3$-Zykel aus $A_n$. Diese aber erzeugen nach Satz \ref{S:Anerz}
  die Gruppe $A_n$, und somit gilt $N=A_n$.
  
  Es bleibt die Existenz eines $3$-Zykels $\rho\in N$ zu zeigen. Dazu starten
  wir mit einem beliebigen Element $1\ne\phi\in N$. Ist $\sigma$ ein
  $3$-Zykel, dann liegt das Produkt der $3$-Zykel $\sigma^{-1}$ und
  $\phi^{-1}\sigma\phi$ in $N$, da $\sigma^{-1}\phi^{-1}\sigma$ und $\phi$ in
  $N$ liegen. Unser Ziel ist es, $\sigma$ geschickt zu w\"ahlen, so dass
  $\rho$ ein $3$-Zykel wird. Das gelingt nicht immer in einem Schritt, aber da
  das Produkt von zwei $3$-Zykeln h\"ochstens $6$ Elemente bewegt, hat $\rho$
  schon mal mindestens $n-6$ Fixpunkte.
  
  Ist $\phi$ ein $3$-Zykel, dann sind wir sowieso fertig. Da $\phi$ gerade
  ist, kann $\phi$ keine Transposition sein.  Somit bleibt der Fall zu
  betrachten, dass $\phi$ mindestens $4$ Punkte bewegt.
  
  Hierzu unterscheiden wir drei F\"alle:

(i) $\phi$ hat eine Zykell\"ange $\ge4$. Dann k\"onnen wir
$\phi:=(a\,b\,c\,d\dots)\dots$ schreiben. Wir setzen
$\sigma:=(a\,b\,c)$. Wegen
$\phi^{-1}\sigma\phi=(a^\phi\,b^\phi\,c^\phi)=(b\,c\,d)$ gilt
$\rho=\sigma^{-1}\phi^{-1}\sigma\phi=(a\,c\,b)(b\,c\,d)=(a\,d\,b)$ ist
$\rho\in N$ ein $3$-Zykel.

(ii) $\phi$ enth\"alt einen $3$-Zykel $(a\,b\,c)$. Da $\phi$ kein $3$-Zykel
ist, wird mindestens ein weiteres Element $d$ bewegt, also
$\phi=(a\,b\,c)(d\,e\dots)\dots$. Wir setzen $\sigma=(a\,b\,d)$. Es folgt
$\phi^{-1}\sigma\phi=(a\,b\,d)^\phi=(b\,c\,e)$ und
$\rho=(a\,d\,b)(b\,c\,e)=(a\,d\,c\,e\,b)$. Somit sind wir im Fall (i), und
k\"onnen so wie dort einen $3$-Zykel konstruieren.

(iii) $\phi$ hat nur Zykell\"angen $2$. Dann gilt $\phi=(a\,b)(c\,d)\dots$,
und wir setzen $\sigma=(a\,c\,e)$, wobei $e$ von $a,b,c,d$ verschieden ist.
(Das geht wegen $n\ge5$!) Es gilt
$\phi^{-1}\sigma\phi=(a\,c\,e)^\phi=(b\,d\,f)$ mit $f=e^\phi$. Ist $e=f$, dann
gilt $\rho=(a\,e\,c)(b\,d\,e)=(a\,b\,d\,e\,c)$, und wir fahren fort wie im
Fall (i). Ist hingegen $e\ne f$, dann ist $\rho=(a\,e\,c)(b\,d\,f)$ das
Produkt zweier disjunkter $3$-Zykel, und wir fahren fort wie im Fall (ii).

Wie wir sehen, gelangen wir stets nach endlich vielen Schritten zu einem
$3$-Zykel aus $N$. Die Behauptung folgt.
\end{proof}
\subsection{Homomorphismen}
Strukturerhaltende Abbildungen spielen in der gesamten Mathematik eine
wichtige Rolle. In der Algebra hei\ss en sie meist Homomorphismen. Eine
Abbildung $\Phi:X\to Y$ zwischen zwei Mengen schreiben wir entweder in der
gew\"ohnlichen oder der Exponentialnotation, d.h.\ das Bild von $x$ unter
$\Phi$ ist $\Phi(x)$ oder $x^\Phi$.
\begin{Definition}
  Eine Abbildung $\Phi:G\to H$ von der Gruppe $G$ in die Gruppe $H$ hei\ss t
  ein \ei{Homomorphismus}, wenn $(xy)^\Phi=x^\Phi y^\Phi$ gilt f\"ur alle
  $x,y\in G$.
\end{Definition}
\paragraph{Einfache Folgerungen}
Die folgenden Aussagen f\"ur einen Homomorphismus $\Phi:G\to H$ erh\"alt man
direkt aus den Definitionen
\begin{itemize}
\item Ein Homomorphismus bildet das neutrale Element auf das neutrale Element
  ab, das folgt aus der K\"urzungsregel und $e^\Phi=(ee)^\Phi=e^\Phi e^\Phi$.
\item Es gilt $(x^{-1})^\Phi=(x^\Phi)^{-1}$, man schreibt daf\"ur auch
  manchmal $x^{-\Phi}$.
\item Die Komposition von Homomorphismen ist ein Homomorphismus.
\item Das Bild $U^\Phi=\{u^\Phi|\,u\in U\}$ einer Untergruppe $U\le G$ ist eine
  Untergruppe von $H$.
\item Das Urbild $V^{\Phi^{-1}}$ einer Untergruppe $V\le H$ ist eine
  Untergruppe von $G$. Ist dabei $V$ normal in $H$, dann ist $V^{\Phi^{-1}}$
  normal in $G$. (Sind auch Bilder von Normalteilern normal?)
\item Sind $x,y$ in $G$ konjugiert, so sind $x^\Phi,y^\Phi$ in $H$ konjugiert.
\item $\Phi$ ist bereits durch seine Werte auf einem Erzeugendensystem von $G$
  festgelegt.
\end{itemize}
Homomorphismen $\Phi:G\to H$ mit speziellen Eigenschaften haben noch weitere
Bezeichnungen:

%  \centering
  \begin{tabular}{ll}
\ei{Monomorphismus:} & $\Phi$ ist injektiv\\
\ei{Epimorphismus:} & $\Phi$ ist surjektiv\\
\ei{Isomorphismus:} & $\Phi$ ist bijektiv\\
\ei{Endomorphismus:} & $G=H$\\
\ei{Automorphismus:} & $G=H$ und $\Phi$ ist bijektiv
  \end{tabular}

Die Umkehrabbildung eines Isomorphismus $\Phi$ bezeichnet man mit
$\Phi^{-1}$. Aus
$xy=(x^{\Phi^{-1}})^\Phi(y^{\Phi^{-1}})^\Phi=
(x^{\Phi^{-1}}y^{\Phi^{-1}})^\Phi$, folgt, nach Anwenden von $\Phi^{-1}$, dass
auch $\Phi^{-1}$ ein Homomorphismus und damit ein Isomorphismus
ist. Insbesondere ist die Menge der Automorphismen einer Gruppe $G$ selber
eine Gruppe, man bezeichnet sie mit $\Aut(G)$. Besteht zwischen den Gruppen
$G$ und $H$ ein Isomophismus, dann nennt man $G$ und $H$ \ei{isomorph}, und
schreibt $G\cong H$.
\paragraph{Beispiele von Homomorphismen}
\begin{itemize}
\item $G=(\CCC,+)$, $H=(\CCC\setminus\{0\},\cdot)$, $x^\Phi:=e^x$.
\item $G=\GL_n(K)$, $H=(K\setminus\{0\},\cdot)$, $x^\Phi:=\det(x)$.
\item $G=S_n$, $H=(\{-1,1\},\cdot)$, $x^\Phi:=(-1)^{l(x)}$.
\item $G=H$ abelsch, $n\in\ZZZ$, $x^\Phi:=x^n$.
\item $N\trianglelefteq G$, $H=G/N$, $x^\Phi:=Nx$.
\end{itemize}
\begin{Definition}
  Der \ei{Kern} eines Homomorphismus $\Phi:G\to H$ ist die Menge der $g\in G$
  mit $g^\Phi=e$, man schreibt auch $\text{Kern}(\Phi)$ f\"ur diese Menge. Als
  Urbild des Normalteilers $\{e\}$ von $H$ ist $\text{Kern}(\Phi)$ ein
  Normalteiler von $G$. Mit $\text{Bild}(\Phi)$ bezeichnen wir das Bild
  $G^\Phi$ von $G$ unter $\Phi$.
\end{Definition}
\begin{Bemerkung*}
  Man rechnet sofort nach, dass ein Homomorphismus $\Phi$ genau dann injektiv
  ist, wenn $\text{Kern}(\Phi)=\{e\}$ gilt.
\end{Bemerkung*}
Analog zu einem Satz der linearen Algebra erhalten wir den folgenden wichtigen
Homomorphiesatz:
\begin{Satz}
  Sei $\Phi:G\to H$ ein Homomorphismus. Dann ist die Abbildung
  $\Psi:G/\text{Kern}(\Phi)\to\text{Bild}(\Phi)$, $\text{Kern}(\Phi)x\mapsto
  x^\Phi$ wohldefiniert, und liefert einen Isomorphismus
  $G/\text{Kern}(\Phi)\cong\text{Bild}(\Phi)$.
\end{Satz}
\begin{proof}[Beweis]
Die Wohldefiniertheit und Surjektivit\"at folgen direkt aus der
Definition. Setze $N:=\text{Kern}(\Phi)$. Wegen $NxNy=Nxy$ gilt
$(NxNy)^\Psi=(Nxy)^\Psi=(xy)^\Phi=x^\Phi y^\Phi=(Nx)^\Psi(Ny)^\Psi$. Somit ist
$\Psi$ ein Homomorphismus. Ferner ist $\Psi$ injektiv, denn aus
$e_H=(Nx)^\Psi=x^\Phi$ folgt $x=e_G$. ($e_G$ und $e_H$ bezeichnen die
neutralen Elemente von $G$ bzw.\ $H$.)  
\end{proof}
Es war wohl schon fr\"uher klar, dass es ``im Prinzip'' nur die zyklischen
$\ZZZ$ und $\ZZZ/n\ZZZ$ gibt. Mit den jetzigen Begriffen k\"onnen wir das
pr\"azise formulieren und beweisen:
\begin{Korollar}
  Bis auf Isomorphie gibt es nur die folgenden zyklischen Gruppen: $(\ZZZ,+)$
  und $(\ZZZ/n\ZZZ,+)$ f\"ur ein $n\in\NNN$.
\end{Korollar}
\begin{proof}[Beweis]
  Sei $g$ ein Erzeuger einer zyklischen Gruppe. Die Abbildung $\ZZZ\to G$,
  $m\mapsto g^m$ ist ein Epimorphismus. Sei $N$ der Kern. Dann gilt nach
  obigem Satz $G\cong\ZZZ/N$. Ist $N=\{0\}$, dann gilt nat\"urlich
  $\ZZZ/N\cong\ZZZ$. Sei nun $N\ne\{0\}$, und $n$ die kleinste nat\"urliche
  Zahl in $N$. Dann gilt $N=n\ZZZ$, und die Behauptung folgt.
\end{proof}
Eine weitere wichtige Erkenntnis ist der Satz, dass jede Gruppe isomorph ist
zu einer Untergruppe einer symmetrischen Gruppe.
\begin{Satz}[Cayley]
  Jede Gruppe $G$ ist isomorph zu einer Untergruppe von $\text{Sym}(G)$.
\end{Satz}
\begin{proof}[Beweis]
  F\"ur $g\in G$ sei $g^\Phi$ die Permutation der Elemente von $G$, welche
  $x\in G$ auf $xg$ abbildet. Aus
  $x^{(gh)^\Phi}=xgh=(x^{g^\Phi})^{h^\Phi}=x^{g^\Phi h^\Phi}$ folgt, dass
  $\Phi:G\to\text{Sym}(G)$ ein Homomorphismus ist. Der Kern von $\Phi$ besteht
  offenbar nur aus dem neutralen Element, somit ist $G$ isomorph zu $G^\Phi$.
\end{proof}
Wie in der Linearen Algebra folgen aus dem Homomorphiesatz einige
Isomorphies\"atze. Beachte, dass wenn $N$ ein Normalteiler und $U$ eine
Untergruppe einer Gruppe $G$ sind, dann ist $UN$ eine Untergruppe von $G$.
\begin{Satz}
  \begin{itemize}
  \item[(a)] Sei $G$ eine Gruppe mit Untergruppe $U$ und Normalteiler $N$.
    Dann gilt $U/U\cap N\cong UN/N$.
  \item[(b)] Seien $G$ eine Gruppe mit Normalteilern $N\subseteq M$. Dann gilt
  $(G/N)/(M/N)\cong G/M$.
  \end{itemize}
\end{Satz}
\begin{proof}[Beweis]
  (a) Sei $\Phi$ der nat\"urliche Homomophismus $G\to G/N$. Das Bild von $U$
  besteht aus den Nebenklassen $uN$ f\"ur $u\in U$, ist also $UN/N$. Der Kern
  der Einschr\"ankung von $\Phi$ auf $U$ ist $U\cap N$, die Behauptung folgt
  nun aus dem Homomorphiesatz.

(b) Die Abbildung $G/N\to G/M$, $gN\mapsto gM$ ist ein wohldefinierter
Epimorphismus mit Kern $M/N$, die Behauptung folgt wiederum aus dem
Homomorphiesatz.
\end{proof}
\subsection{Gruppenoperationen}
In diesem Abschnitt wollen wir Operationen von Gruppen auf Mengen
betrachten. Einige grundlegende Aussagen werden sp\"ater als Hilfsmittel zum
Beweis interner Aussagen \"uber endliche Gruppen dienen, die nichts mit
Gruppenoperationen zu tun haben.
\begin{Definition}
  Eine \ei{Operation} einer Gruppe $G$ auf einer Menge $M$ ist eine Abbildung
  $M\times G\to M$, $(m,g)\mapsto m^g$, f\"ur die $m^e=m$ und $(m^g)^h=m^{gh}$
  gilt f\"ur alle $m\in M$, $g\in G$.
\end{Definition}
\begin{Bemerkung*}
  $G$ operiere auf $M$ wie oben. Sei $\Phi:G\to\text{Sym}(M)$ die Abbildung,
  die $g\in G$ auf die Permutation $M\to M$, $m\mapsto m^g$ abbildet. Dann ist
  $\Phi$ ein Homomorphismus. Umgekehrt liefert jeder Homomorphismus
  $G\to\text{Sym}(G)$ eine Operation von $G$ auf $M$.
\end{Bemerkung*}
Operiert die Gruppe $G$ auf der Menge $M$, so erh\"alt man auf nat\"urliche
Weise eine \"Aquivalenzrelation auf $M$: Zwei Elemente $m_1,m_2\in M$ sind
\"aquivalent genau dann, wenn es ein Element $g\in G$ gibt mit
$m_2=m_1^g$. (Man verifiziere die drei Eigenschaften Reflexivit\"at, Symmetrie
und Transitivit\"at.) Die \"Aquivalenzklassen nennt man \ei{Bahnen}. Die Bahn
durch $m$ besteht offenbar aus den Elementen $m^g$, $g\in G$, und wird daher
mit $m^G$ bezeichnet.

Eine Operation von $G$ auf $M$ hei\ss t \ei{transitiv}, wenn $M$ nur aus einer
Bahn besteht. In diesem Fall wird $M$ gelegentlich auch \ei{homogener Raum}
genannt.

\paragraph{Beispiele}
\begin{itemize}
\item Sei $U$ eine Untergruppe der Gruppe $G$. Dann operiert $G$ auf dem
  Nebenklassenraum $U\backslash G$: Hierbei schickt $g$ die Nebenklasse $Ux$
  auf die Nebenklasse $Uxg$. Diese Operation ist offensichtlich transitiv.
  Gleich werden wir sehen, dass jede transitive Operation auf einer Menge $M$,
  bis auf Umbenennung der Elemente von $M$, eine solche Operation auf einem
  Nebenklassenraum ist.
  
  Diese Operation liefert einen Homomorphismus $G\to\text{Sym}(U\backslash
  G)$. Es ist nat\"urlich interessant, den Kern $N$ zu bestimmen. Dabei
  besteht $N$ aus genau den Elementen $g\in G$, die jede Nebenklasse $Ux$
  festlassen, also $Uxg=Ux$ f\"ur alle $x\in G$ erf\"ullen. Aber $Uxg=Ux$ ist
  \"aquivalent zu $xgx^{-1}\in U$, und das ist \"aquivalent zu $g\in U^x$.
  Daher besteht $N$ aus dem Schnitt der Konjugierten von $U$. Gilt
  $[G:U]=n<\infty$, dann ist nach dem Homomorphiesatz $G/N$ isomorph zu einer
  Untergruppe der symmetrischen Gruppe auf $n$ Elementen, insbesondere ist
  $[G:N]$ endlich und ein Teiler von $n!$.
\item Die Gruppe $G$ operiert auf sich selbst durch Konjugation. Dabei schickt
  $g$ das Element $x$ auf $g^{-1}xg=x^g$. F\"ur $\abs{G}>1$ ist diese
  Operation nicht transitiv, da $\{e\}$ eine Bahn ist.
\item Die lineare Gruppe $\GL_n(K)$ operiert auf dem Vektorraum $K^n$, aber
  sie operiert zum Beispiel auch auf der Menge der $1$-dimensionalen
  Unterr\"aume von $K^n$.
\item Sei $G$ die Gruppe der gleichsinnigen Symmetrien eines W\"urfels. Diese
  Gruppe hat Ordnung $24$, denn eine Seitenfl\"ache l\"asst sich auf $6$
  m\"ogliche Fl\"achen abbilden, und danach hat man noch $4$ m\"ogliche
  Drehungen dieser Fl\"ache. Diese Gruppe hat verschiedene transitive
  Operationen auf Objekten, die zum W\"urfel geh\"oren: $G$ operiert transitiv
  auf den $8$ Ecken, transitiv auf den $6$ Fl\"achen, transitiv auf den $12$
  Kanten, aber auch transitiv auf den $4$ Raumdiagonalen.
\end{itemize}

Operiert $G$ auf $M$, so nennt man $G_m:=\{g\in G|\;m^g=m\}$ den
\ei{Stabilisator} von $m\in M$. Andere gebr\"auchliche Bezeichnungen sind
\ei{Punktstabilisator}, \ei{Standgruppe} oder \ei{Isotropiegruppe}. Man
rechnet sofort nach, dass $G_m$ eine Untergruppe von $G$ ist.

\begin{Definition}
  Die Gruppe $G$ operiere auf den Mengen $M$ und $N$. Eine Abbildung
  $\phi:M\to N$ nennen wir $G$-\ei{\"aquivariant}, wenn
  $(m^g)^\phi=(m^\phi)^g$ gilt f\"ur alle $m\in M$, $g\in G$.
\end{Definition}
\begin{Satz}
  Die Gruppe $G$ operiere transitiv auf der Menge $M$. Sei $U$ der
  Stabilisator eines Elements $m\in M$. Dann wird durch $\phi:U\backslash G\to
  M$, $Ux\mapsto m^x$ eine $G$-\"aquivariante Bijektion definiert.
\end{Satz}
\begin{proof}[Beweis]
  Zun\"achst m\"ussen wir sehen, dass $\phi$ wohldefiniert ist. Dazu sei etwa
  $Ux=Uy$. Dann gilt $y=ux$ mit $u\in U$, und wegen $m^u=m$ folgt
  $m^x=m^{ux}=m^y$.
  
  $\phi$ ist offenbar surjektiv. Ferner ist $\phi$ injektiv, denn aus
  $m^x=m^y$ folgt $m^{yx^{-1}}=m$, also $yx^{-1}\in U$ und dann $Ux=Uy$.
  
  Die $G$-\"Aquivarianz folgt aus
  $((Ux)^\phi)^g=(m^x)^g=m^{xg}=(Uxg)^\phi=((Ux)^g)^\phi$.
\end{proof}
Eine wichtige Konsequenz ist
\begin{Korollar}\label{K:Bahnlaenge}
  Die Gruppe $G$ operiere auf $M$. Die L\"ange der Bahn durch $m$ berechnet
  sich durch $\abs{m^G}=[G:G_m]$.
\end{Korollar}
\begin{proof}[Beweis]
  $G$ operiert transitiv auf der Bahn $m^G$. Nach obigem Satz besteht eine
  Bijektion zwischen $m^G$ und $\text{St}_G(m)\backslash G$.
\end{proof}
Folgende einfache Aussage wird h\"aufig ohne Kommentar verwendet:
\begin{Lemma}
  $G$ operiere auf $M$. Die Elemente $u,v\in M$ seien in einer gemeinsamen
  Bahn. Dann sind die Stabilisatoren $G_u$ und $G_v$ in $G$ konjugiert.
  Genauer: Ist $v=u^g$, dann gilt $G_v=g^{-1}G_ug=G_u^g$.
\end{Lemma}
\begin{proof}[Beweis]
  Sei $v=u^g$. Dann ist $h\in G_v$ $\Longleftrightarrow$ $v^h=v$
  $\Longleftrightarrow$ $u^{gh}=u^g$ $\Longleftrightarrow$ $u^{ghg^{-1}}=u$
  $\Longleftrightarrow$ $ghg^{-1}\in G_u$ $\Longleftrightarrow$ $h\in
  g^{-1}G_ug$.
\end{proof}
An einigen Beispielen wollen wir die Anwendbarkeit der neuen Resultate und
Konzepte verdeutlichen. Vorher ben\"otigen wir noch einige Begriffe.
\begin{Definition}
  Die Gruppe $G$ operiere durch Konjugation auf sich selbst. Den Stabilisator
  eines Elements $x$ unter dieser Operation nennt man den \ei{Zentralisator}
  von $x$ in $G$, und schreibt daf\"ur $C_G(x)$. Die Menge $C_G(x)$ besteht
  also aus allen $g\in G$ mit $gx=xg$.
  
  Die Bahn $x^G=\{g^{-1}xg|\;g\in G\}$ nennt man \ei{Konjugationsklasse} von
  $x$ in $G$.
  
  Ist $X$ eine Teilmenge von $G$, so nennt man $C_G(X):=\{g\in
  G|\;x^g=x\text{ f\"ur alle $x\in X$}\}$ den Zentralisator von $X$ in $G$.
  Offenbar ist $C_G(X)$ der Schnitt der Zentralisatoren $C_G(x)$ f\"ur alle
  $x\in X$.
  
  Ist speziell $X=G$, so nennt man $Z(G):=C_G(G)$ das \ei{Zentrum} von $G$.
  
  $G$ operiert auch durch Konjugation auf der Menge der Teilmengen von $G$.
  Ist $X$ eine Teilmenge, dann nennt man den Stabilisator von $X$ den
  \ei{Normalisator} von $X$ in $G$, und schreibt daf\"ur $N_G(X)$. Es gilt
  also $N_G(X)=\{g\in G|\;X^g=X\}$.
  
  Da Stabilisatoren Untergruppen sind, folgt sofort, dass $N_G(X)$, $C_G(x)$,
  $C_G(X)$ und $Z(G)$ Untergruppen von $G$ sind, was man aber auch direkt
  nachrechnen kann. F\"ur $x\in X$ folgt aus der Definition sofort $Z(G)\le
  C_G(X)\le C_G(x)$ und $C_G(X)\le N_G(X)$.
\end{Definition}
\begin{Satz}[Bahnengleichung]
Die Gruppe $G$ operiere auf der endlichen Menge $M$. Seien $m_1,m_2,\dots,m_r$
Repr\"asentanten der Bahnen. Dann gilt
\[
\abs{M}=\sum_{i=1}^r[G:G_{m_i}].
\]
\end{Satz}
\begin{proof}[Beweis]
  $M$ ist die disjunkte Vereinigung der Bahnen durch die $m_i$, die Behauptung
  folgt dann aus Satz \ref{K:Bahnlaenge}.
\end{proof}
\begin{Satz}[Klassengleichung]
  Sind $x_1,x_2,\dots,x_r$ die Repr\"asentanten der Konjugationsklassen der
  endlichen Gruppe $G$, dann gilt
\[
\abs{G}=\sum_{i=1}^r[G:C_G(x_i)].
\]
\end{Satz}
\begin{proof}[Beweis]
  Dies ist obiger Satz, angewandt auf die Operation von $G$ auf sich durch
  Konjugation.
\end{proof}
\begin{Korollar}\label{K:Zentrum_p-Gruppe}
  Sei $p$ eine Primzahl, $n\in\NNN$, und $G$ eine Gruppe der Ordnung
  $p^n$. Dann gilt $\abs{Z(G)}>1$.
\end{Korollar}
\begin{proof}[Beweis]
  Seien $1=x_1,x_2,\dots,x_r$ die Repr\"asentanten der Konjugationsklassen von
  $G$. Dann gilt $p^n=\sum_{i=1}^r[G:C_G(x_i)]$. Nach dem Satz von Lagrange
  sind die Indizes $[G:C_G(x_i)]$ Potenzen von $p$. Die linke Seite der
  Klassengleichung ist durch $p$ teilbar, aber $[G:C_G(x_1)]=1$ ist es
  nicht. Daher muss es einen Index $i>1$ geben, so dass auch $[G:C_G(x_i)]$
  nicht durch $p$ teilbar ist. Als Potenz von $p$ muss dann $[G:C_G(x_i)]=1$
  gelten, also $G=C_G(x_i)$ und somit $1\ne x_i\in Z(G)$.
\end{proof}
Die folgende Formel wird manchmal die Burnsidesche Bahnenformel genannt,
obwohl sie schon fr\"uher bei Cauchy und Frobenius auftaucht.
\begin{Satz}[Cauchy-Frobenius Bahnenformel]
  Die endliche Gruppe $G$ operiere auf der endlichen Menge $M$. F\"ur $g\in G$
  sei $\chi(g)$ die Anzahl der Fixpunkte $m\in M$ mit $m^g=m$. Dann ist
\[
\frac{1}{G}\sum_{g\in G}\chi(g)
\]
die Anzahl der Bahnen von $G$ auf $M$.
\end{Satz}
\begin{proof}[Beweis]
  Es gen\"ugt die Aussage f\"ur transitive Operationen zu beweisen, da die
  Anzahl der Fixpunkte von $g$ die Summe der Anzahlen der Fixpunkte auf den
  einzelnen Bahnen ist. Sei $S$ die Menge der Paare $(m,g)\in M\times G$ mit
  $m^g=m$. Wir z\"ahlen $S$ einmal \"uber die Elemente $g$, und einmal \"uber
  die Elemente $m$ ab. Die erste Abz\"ahlung liefert $\abs{S}=\sum_{g\in
    G}\chi(m)$, und die zweite Abz\"ahlung ergibt $\abs{S}=\sum_{m\in
    M}\abs{G_m}$. Wegen der Transitivit\"at gilt $\abs{M}=[G:G_m]$ f\"ur alle
  $m\in M$, also $\abs{G_m}=\abs{G}/\abs{M}$, und die Behauptung folgt.
\end{proof}
Wir notieren zwei Folgerungen.
\begin{Korollar}
  Die Gruppe $G$ operiere transitiv auf der endlichen Menge $M$. Es gelte
  $\abs{M}>1$. Dann enth\"alt $G$ ein fixpunktfreies Element.
\end{Korollar}
\begin{proof}[Beweis]
  Wir d\"urfen $G$ ersetzen mit seinem homomorphen Bild in $\text{Sym}(M)$,
  insbesondere ist dann $G$ endlich. Nach obigem Satz haben die Elemente von
  $G$ durchschnittlich einen Fixpunkt. Aber $e$ hat $\abs{M}>1$ Fixpunkte,
  daher muss es zum Ausgleich ein Element mit weniger als einem Fixpunkt
  geben.
\end{proof}
\begin{Korollar}
  Die echte Untergruppe $U$ von $G$ habe endlichen Index. Dann ist $G$ nicht
  die Vereinigung der Konjugierten von $U$.
\end{Korollar}
\begin{proof}[Beweis]
  Betrachte die Operation von $G$ auf $U\backslash G$. Nach obigem Korollar
  hat $G$ ein fixpunktfreies Element $g$. Das hei\ss t $Uxg\ne Ux$ f\"ur alle
  $x\in G$, und somit $g\notin U^x$ f\"ur alle $x\in G$.
\end{proof}
\begin{Bemerkung*}
  Diese Aussage wird falsch, wenn man die Voraussetzung des endlichen Index
  wegl\"asst. In der Theorie der sogenannten algebraischen Gruppen spielen
  gerade gewisse Untergruppen (die Borelgruppen) eine wichtige Rolle, deren
  Konjugierte im zusammenh\"angenden Fall die Gruppe ausf\"ullen. Ein
  Spezialfall davon ist z.B.\ die bekannte Aussage aus der Linearen Algebra,
  dass sich jede Matrix aus $\GL_n(\CCC)$ auf obere Dreiecksgestalt
  transformieren l\"asst.
\end{Bemerkung*}
\subsection{Produkte}
Sei $I$ eine Indexmenge, und f\"ur alle $i\in I$ sei $G_i$ eine Gruppe. Die
Menge der Tupel $(g_i)_{i\in I}$ mit $g_i\in G_i$ wird mit der
komponentenweise Multiplikation $(g_i)_{i\in I}(h_i)_{i\in I}:=(g_ih_i)_{i\in
  I}$ zu einer Gruppe. Man nennt diese Gruppe das \emph{direkte
  Produkt}\index{direktes Produkt} der Gruppen $G_i$, und schreibt daf\"ur
$\prod_{i\in I}G_i$. Ist $I$ die endliche Menge $\{1,2,\dots,n\}$, dann
schreibt man daf\"ur auch $G_1\times G_2\times\dots\times G_n$, und stellt die
Elemente durch $n$-Tupel $(g_1,g_2,\dots,g_n)$ dar.

Das \emph{eingeschr\"ankte direkte Produkt}\index{eingeschr\"anktes direkte
  Produkt} ist der Normalteiler bestehend aus all den Tupeln $g_i\in G_i$, in
  denen f\"ur alle bis auf endlich viele Indizes $i$ das Element $g_i$ das
  neutrale Element ist. Man nennt das manchmal auch \ei{direkte
  Summe}. Gebr\"auchliche Schreibweisen sind $\bigoplus_{i\in I}G_i$ und
  $\coprod_{i\in I}G_i$. Ist die Indexmenge $I$ endlich, dann stimmen
  nat\"urlich direkte Summe und direktes Produkt \"uberein.

  \begin{Lemma}
    Seien $A$ und $B$ Normalteiler der Gruppe $G$ mit $A\cap B=\{e\}$ und
    $G=AB$. Dann gilt $G\cong A\times B$.
  \end{Lemma}
  \begin{proof}[Beweis]
    Sei $a\in A$, $b\in B$. Wegen $a^{-1}b^{-1}a\in B$ und $b\in B$ gilt
    $a^{-1}b^{-1}ab\in B$. Wegen $a^{-1}\in A$ und $b^{-1}ab\in A$ gilt aber
    auch $a^{-1}b^{-1}ab\in A$, also $a^{-1}b^{-1}ab\in A\cap B=\{e\}$ und
    somit $ab=ba$. Folglich ist die Abbildung $A\times B\to G$, $(a,b)\mapsto
    ab$ ein Homomorphismus. Gem\"a\ss\ Voraussetzung ist dieser Homomorphismus
    surjektiv. Er ist aber auch injektiv, denn aus $ab=e$ folgt $a,b\in A\cap
    B$, also $a=b=e$.
  \end{proof}
  \begin{Bemerkung*}
    Das Produkt $AB$ im Lemma nennt man auch das \emph{interne direkte
    Produkt}\index{internes direktes Produkt} der Normalteiler $A$ und
    $B$. Im Gegensatz dazu bezeichnet man $A\times B$ manchmal als das
    \emph{externe direkte Produkt}\index{externes direktes Produkt} von $A$
    und $B$.
  \end{Bemerkung*}
Eine Verallgemeinerung des obigen Lemmas auf mehrere Faktoren ist
\begin{Satz}
  Sei $I$ eine Indexmenge, und $G_i$, $i\in I$ seien Untergruppen einer Gruppe
  $G$. Auf $I$ w\"ahle man eine Totalordnung $<$. Die folgenden Aussagen sind
  \"aquivalent.
  \begin{itemize}
\item[(a)] $G_i\trianglelefteq G$ f\"ur alle $i$, die $G_i$ erzeugen $G$, und
  $G_i\cap \hat G_i=\{e\}$, wo $\hat G_i$ das Erzeugnis der Gruppen $G_j$ mit
  $j\in I\setminus\{i\}$ ist.
\item[(b)] $G_i\trianglelefteq G$ f\"ur alle $i$, und f\"ur jedes $g\in G$
  gibt es, bis auf Faktoren $e$, genau eine Darstellung $g=g_{i_1}g_{i_2}\dots
  g_{i_n}$ mit $g_i\in G_i$ und $i_1<i_2<\dots i_n$.
\item[(c)] F\"ur alle $i\ne j$, $i,j\in I$ gilt $g_ig_j=g_jg_i$ f\"ur $g_i\in
  G_i$, $g_j\in G_j$, und f\"ur jedes $g\in G$ gibt es, bis auf Faktoren $e$,
  genau eine Darstellung $g=g_{i_1}g_{i_2}\dots g_{i_n}$ mit $g_i\in G_i$ und
  $i_1<i_2<\dots i_n$.
  \end{itemize}
  Gilt eine der Aussagen, dann ist $G$ ist isomorph zur direkten Summe der
  $G_i$.
\end{Satz}
\begin{proof}[Beweis]
Aus $G_i\trianglelefteq G$ folgt, wie im Beweis des obigen Lemmas, dass
  $a^{-1}b^{-1}ab\in G_i\cap G_j$ liegt f\"ur $a\in G_i$, $b\in G_j$.
  
  ``(a) $\Longrightarrow$ (b)'': F\"ur $i\ne j$ gilt $G_j\le \hat G_i$, also
  $G_i\cap G_j=\{e\}$. Daher vertauschen die Elemente aus $G_i$ mit denen aus
  $G_j$. Da die $G_i$ ganz $G$ erzeugen, hat jedes Element $g\in G$ eine
  Darstellung $g=g_{i_1}g_{i_2}\dots g_{i_n}$ mit $g_i\in G_i$ und
  $i_1<i_2<\dots i_n$. Sei eine weitere solche Darstellung f\"ur $g$ gegeben.
  Durch Erg\"anzung der Indexmengen durch Einf\"ugen der Faktoren $e$ d\"urfen
  wir $g=h_{i_1}h_{i_2}\dots h_{i_n}$ mit $h_i\in G_i$ annehmen. Es folgt
  $h_{i_1}^{-1}g_{i_1}=h_{i_2}h_{i_3}\dots h_{i_n}g_{i_n}^{-1}\dots
  g_{i_3}^{-1}g_{i_2}^{-1}\in G_{i_1}\cap \hat G_{i_1}$ und daraus
  $g_{i_1}=h_{i_1}$. Fortsetzen des Verfahrens liefert nacheinander
  $g_{i_2}=h_{i_2}$, \dots, $g_{i_n}=h_{i_n}$. Die Eindeutigkeitsaussage
  folgt.
  
  ``(b) $\Longrightarrow$ (c)'': Dies folgt aus dem Beginn des Beweises.
  
  ``(c) $\Longrightarrow$ (a)'': Aus der Vertauschbarkeitsvoraussetzung
  folgt $G_i\trianglelefteq G$ f\"ur alle $i$. Wegen der eindeutigen
  Elementdarstellung aus den $G_i$ folgt $G_i\cap \hat G_i=\{e\}$.
  
  Es gelte nun (a), (b) und (c). Dann ist die nat\"urliche Abbildung
  $\bigoplus_{i\in I}G_i\to G$, die durch multiplikative Fortsetzung der
  Einbettungsabbildungen $G_i\to G$ definiert ist, ein Epimorphismus. Sei
  $(g_{i_1},g_{i_2},\dots,g_{i_n})$ im Kern, also $g_{i_1}g_{i_2}\dots
  g_{i_n}=e$. Wegen der eindeutigen Darstellbarkeit von $e$ folgt
  $g_{i_1}=g_{i_2}=\dots=g_{i_n}=e$, und die Abbildung ist somit auch
  injektiv. Die Behauptung folgt.
\end{proof}
\paragraph{Beispiele von Produkten}
\begin{itemize}
\item Seien $m$ und $n$ teilerfremde nat\"urliche Zahlen, und $G$ eine
  zyklische Gruppe der Ordnung $mn$. Dann hat $G$ Untergruppen $A$ der Ordnung
  $m$ und $B$ der Ordnung $n$. Nach dem Satz von Lagrange gilt $A\cap
  B=\{e\}$. Nach obigem Satz gilt $AB\cong A\times B$. Die rechte Seite hat
  Ordnung $mn$, daher auch die linke Seite und es folgt $G=AB\cong A\times
  B$.
\item Sei $\CCC^\star$ die multiplikative Gruppe der komplexen Zahlen, und
  $S^1$ die Untergruppe der komplexen Zahlen vom Betrag $1$. Sei $P$ die
  Gruppe der positiven reellen Zahlen. Dann gilt $\CCC^\star\cong S^1\times
  P$.
\item Sei $V$ ein Vektorraum \"uber einem K\"orper $K$, und $\{e_i|i\in I\}$
  eine (eventuell unendliche) Basis. Dann ist $(V,+)$ die interne direkte
  Summe der Unterr\"aume $Ke_i$, und isomorph zur externen direkten Summe
  $\bigoplus_{i\in I}K$.
\end{itemize}
In obigem Beispiel hatten wir ein etwas indirektes Argument zur Bestimmung der
Ordnung von $AB$. Die Gr\"o\ss e von Produkten $AB$ l\"asst sich ganz
allgemein bestimmen. Man beachte, dass $AB$ im allgemeinen keine Gruppe sein
muss.
\begin{Satz}
  Seien $A$ und $B$ Untergruppen der endlichen Gruppe $G$. Dann gilt
  $\abs{AB}=\frac{\abs{A}\abs{B}}{\abs{A\cap B}}$.
\end{Satz}
\begin{proof}[Beweis]
  1.\ Beweis: Sei $\{b_1,b_2,\dots,b_n\}$ ein Vertretersystem der
  Rechtsnebenklassen $(A\cap B)\backslash B$. Dann sind die Nebenklassen
  $Ab_i$ disjunkt, denn aus $Ab_i=Ab_j$ folgt $b_ib_j^{-1}\in A\cap B$, also
  $(A\cap B)b_i=(A\cap B)b_j$. Ferner hat jedes Element aus $AB$ die Form
  $ab_i$ mit $a\in A$: Dazu sei $a'\in A$, $b\in B$. Schreibe $b=a''b_i$ f\"ur
  geeignetes $a''\in A\cap B$. Dann gilt $a'b=a'a''b_i$, und die Behauptung
  folgt mit $a=a'a''$.
    
  2.\ Beweis: Betrachte die Operation von $G$ auf $A\backslash G$. Die Bahn
  der Untergruppe $B$ durch $A$ besteht aus den Nebenklassen $Ab$, $b\in B$.
  Die Vereinigung dieser Nebenklassen ist $AB$, und jede Nebenklasse besteht
  aus $\abs{A}$ Elementen. Daher hat die Bahn die L\"ange $\abs{AB}/\abs{A}$.
  Der Stabilisator in $B$ von $A$ ist $A\cap B$. Nach Korollar
  \ref{K:Bahnlaenge} wird die Bahnl\"ange durch $[B:A\cap B]$ gegeben, und die
  Behauptung folgt.
\end{proof}
Wir kommen nun zum wichtigen Begriff des \emph{semidirekten
  Produkts}\index{semidirektes Produkt}, einer Verallgemeinerung des direkten
Produkts zweier Faktoren. Hierzu sei $U$ eine Untergruppe der Gruppe $G$, und
$N$ ein Normalteiler von $G$, so dass $U\cap N=\{e\}$ und $G=UN$ gelten. In
dieser Situation sagt man, $G$ sei das \emph{semidirekte Produkt} des
Normalteilers $N$ mit der Untergruppe $H$. Ferner nennt man $H$ ein
\ei{Komplement} von $N$ in $G$. Im folgenden konstruieren wir aus dieser
Situation eine Verallgemeinerung des externen direkten Produkts. Dazu
beobachten wir, dass jedes Element aus $G$ eine eindeutige Darstellung der
Form $un$ mit $u\in U$, $n\in N$ hat. Das Ziel ist es, das Produkt zweier
solcher Elemente $u_1n_1$ und $u_2n_2$ wieder in dieser Form darzustellen. Wir
rechnen
\[
u_1n_1u_2n_2=u_1u_2u_2^{-1}n_1u_2n_2=u_1u_2n_1^{u_2}n_2.
\]
Man beachte, dass $u_1u_2\in U$ und $n_1^{u_2}n_2\in N$ gilt. Um also solche
Produkte zu berechnen, muss man Produkte innerhalb $U$ und $N$ berechnen, und
die Konjugationswirkung von $H$ aud $N$ kennen. Das f\"uhrt direkt auf den
Begriff des (externen) semidirekten Produkts.
\begin{Satz}
  Seien $N$ und $U$ Gruppen, und $\phi:U\to\Aut(N)$ ein Homomorphismus. Auf
  der Menge $G$ der Paare $(u,n)$, $u\in U$, $n\in N$ definiere man ein
  Produkt durch $(u_1,n_1)(u_2,n_2)=(u_1u_2,n_1^{u_2^\phi}n_2)$. Mit
  diesem Produkt ist $G$ eine Gruppe. $G$ hat die zu $U$ und $N$ isomorphen
  Untergruppen $\tilde U=\{(u,e)|u\in U\}$ und $\tilde N=\{(e,n)|n\in
  N\}$. $G$ ist ein semidirektes Produkt des Normalteilers $\tilde N$ mit
  $\tilde U$. Es gilt $(e,n)^{(u,e)}=(e,n^{u^\phi})$.
\end{Satz}
\begin{proof}[Beweis]
  Die Aussagen ergeben sich alle durch direktes Nachrechnen.
\end{proof}
\begin{Bemerkung*}
  Dir erste Teil des Satzes liefert eine wichtige Konstruktionsmethode f\"ur
  Gruppen. Gelegentlich schreibt man f\"ur $G$ auch $U\ltimes_\phi N$. Ist
  $\phi$ der triviale Homomorphismus, dann ist $U\ltimes_\phi N$ einfach das
  direkte Produkt $U\times N$.
\end{Bemerkung*}
\paragraph{Beispiele semidirekter Gruppen}
\begin{itemize}
\item Sei $S_n$ die symmetrische Gruppe auf $n$ Punkten, $A_n$ die
  alternierende Gruppe, und $C$ die von einer Transposition erzeugte
  Untergruppe von $S_n$. Dann gilt $S_n=C\ltimes A_n$, aber nicht $S_n=C\times
  A_n$.
\item Sei $C_n$ eine zyklische Gruppe der Ordnung, $C=\gen{c}$ eine Gruppe der
  Ordnung $2$, und $\phi:C\to\Aut(C_n)$ definiert durch
  $g^{c^\phi}:=g^{-1}$. Die Gruppe $G:=C\ltimes_\phi C_n$ ist die
  \ei{Diedergruppe} $D_n$ der Ordnung $2n$. F\"ur $n\ge3$ kann man sich diese
  Gruppe geometrisch veranschaulichen, sie besteht aus den
  Kongruenzabbildungen eines regul\"aren $n$-Ecks. Die Gruppe $C_n$ besteht
  dabei aus den gleichsinnigen Kongruenzen, und $c$ ist eine Spiegelung.
\item Es gilt $\GL_n(K)=U\ltimes \SL_n(K)$, wo $U$ die Untergruppe der
  Diagonalmatrizen ist, die au\ss er links oben nur $1$er enth\"alt.
\end{itemize}
\subsection{Endliche abelsche Gruppen}
Das Ziel dieses Abschnitts ist die Bestimmung der endlichen abelschen Gruppen.
Wir beginnen mit einem einfachen
\begin{Lemma}
  Seien $a,b$ Elemente einer Gruppe mit teilerfremden Ordnungen. Ferner gelte
  $ab=ba$. Dann gilt $\ord(ab)=\ord(a)\ord(b)$.
\end{Lemma}
\begin{proof}[Beweis]
  Sei $r=\ord(a)$, $s=\ord(b)$. Es gilt $(ab)^{rs}=(a^r)^s(b^s)^r=e$, daher
  ist die Ordnung von $ab$ h\"ochstens $rs$. Sei nun $m\in\NNN$ mit
  $(ab)^m=e$. Dann gilt $a^m=(b^{-1})^m$. Nach dem Satz von Lagrange haben die
  beiden Gruppen $\gen{a}$ und $\gen{b}$ trivialen Schnitt. Es folgt
  $a^m=b^m=e$, und somit sind $r$ und $s$ Teiler von $m$. Es folgt $rs=m$.
\end{proof}
\begin{Lemma}
  Sei $G$ eine endliche abelsche Gruppe, und $U$ eine zyklische Untergruppe
  maximaler Ordnung. Dann ist $\ord(g)$ ein Teiler von $\abs{U}$ f\"ur alle
  $g\in G$.
\end{Lemma}
\begin{proof}[Beweis]
  Sei $g\in G$, $p$ eine Primzahl, und $p^r$ ein Teiler von $\ord(g)$. Wir
  werden zeigen, dass $p^r$ ein Teiler von $\abs{U}$ ist. Da das f\"ur alle
  $p$ gilt folgt daraus die Behauptung.
  
  Da $p^r$ ein Teiler von $\ord(g)$ ist gibt es ein Element $a\in\gen{g}$ der
  Ordnung $p^r$. Sei $\abs{U}=p^em$ mit $\ggT(p,m)=1$, und $b\in U$ mit
  Ordnung $m$. Nach obigem Lemma hat $ab$ die Ordnung $p^rm$, wegen der
  Maximalit\"at von $\abs{U}$ folgt $p^rm\le p^em$, also $r\le e$ und die
  Behauptung folgt.
\end{proof}
Als Vorbereitung des Hauptsatzes ben\"otigen wir
\begin{Lemma}
  Sei $G$ eine endliche abelsche Gruppe, und $U$ eine zyklische Untergruppe
  maximaler Ordnung. Dann hat $U$ ein Komplement in $G$.
\end{Lemma}
\begin{proof}[Beweis]
  Wir beweisen die Aussage durch vollst\"andige Induktion \"uber $\abs{G}$.
  
  Ist $G=U$, dann gibt es nichts zu beweisen. Sei also $\abs{U}<\abs{G}$. Wir
  w\"ahlen $g\in G\setminus U$ von minimaler Ordnung. Sei $u$ ein Erzeuger von
  $U$. Nach obigem Lemma besitzt $U$ eine Untergruppe $W$ der Ordnung
  $\ord(g)$. Sei $w$ ein Erzeuger von $W$. Sei $p$ ein Primteiler von
  $\ord(g)$. Dann gilt $g^p\in U$ (wegen der Minimalit\"at von $\ord(g)$ mit
  $g\notin U$), aber $\gen{g^p}=\gen{w^p}$, da die Untergruppen der zyklischen
  Gruppe $U$ schon durch ihre Ordnungen eindeutig bestimmt sind. Es gibt also
  $i\in\NNN$ mit $g^p=w^{pi}$, und daher $(gw^{-i})^p=e$. Aber $gw^{-i}\notin
  U$ ist somit ein Element der Ordnung $p$.
  
  Somit ist $N:=\gen{g}w^{-i}$ eine Untergruppe mit $\abs{N}>1$ und $U\cap
  N=\{e\}$.  Sei $\phi:G\to G/N$, $g\mapsto Ng$ der nat\"urliche
  Epimorphismus. F\"ur $x\in G$ folgt aus $x^n=e$ nat\"urlich
  $(x^\phi)^n=e^\phi$, also $\ord(x^\phi)\le\ord(x)$. Wegen $UN/N\cong
  U/(U\cap N)\cong U$ gilt $\abs{U^\phi}=\abs{U}$. Somit ist $U^\phi$ die
  gr\"o\ss te zyklische Untergruppe von $G^\phi$. Nach Induktionsannahme hat
  $U^\phi$ in $G^\phi$ ein Komplement. Sei $V\ge N$ das Urbild in $G$ dieses
  Komplements in $G$.  Dann gilt $UN\cap V\le N$, also $U\cap V\le U\cap
  N=\{e\}$. Ferner gilt $G^\phi=U^\phi V^\phi$, also $G=(UN)V=U(NV)=UV$, und
  die Behauptung folgt.
\end{proof}
Aus diesem Lemma folgt nun, wiederum durch Induktion \"uber $\abs{G}$, das
Hauptergebnis dieses Abschnitts.
\begin{Satz}
  Jede endliche abelsche Gruppe ist eine direkte Summe zyklischer
  Untergruppen.
\end{Satz}
Im folgenden wollen wir noch die Eindeutigkeit solcher Zerlegungen in
zyklische Gruppen kl\"aren. In einem fr\"uheren Beispiel sahen wir schon, dass
jede endliche zyklische Gruppe eine direkte Summe zyklischer Gruppen von
Primpotenzordnung ist. An obigem Satz sehen wir schon, dass es zu jedem Teiler
$m$ der Ordnung einer endlichen abelschen Gruppe auch eine Untergruppe der
Ordnung $m$ gibt. Diese Untergruppe muss nicht eindeutig sein, aber es gilt
\begin{Lemma}
  Sei $p$ eine Primzahl, und $p^m$ die h\"ochste Potenz von $p$, die die
  Ordnung einer endlichen abelschen Gruppe teilt. Dann gibt es genau eine
  Untergruppe der Ordnung $p^m$.
\end{Lemma}
\begin{proof}[Beweis]
  Seien $A$ und $B$ Untergruppen der Ordnung $p^m$. Wegen
  $\abs{AB}=\abs{A}[B:A\cap B]$ ist auch die Ordnung der Untergruppe $AB$ eine
  Potenz von $p$, die aber $p^m$ teilen muss. Daher gilt $B=A\cap B$, also
  $B\le A$ und aus Symmetriegr\"unden $A=B$.
\end{proof}
Hieraus folgt
\begin{Lemma}
  Sei $n=\prod p_i^{e_i}$ die Primfaktorzerlegung der Ordnung einer endlichen
  abelschen Gruppe $G$, und $G_{p_i}$ die Untergruppe der Ordnung
  $p_i^{e_i}$. Dann ist $G\cong\bigoplus G_{p_i}$.
\end{Lemma}
Um eine bis auf Isomorphie eindeutige Darstellung einer abelschen Gruppe zu
bekommen m\"ussen wir nur noch zyklische Gruppen von Primpotenzordnung
betrachten. Wir wissen schon, dass solche Gruppen direkte Summen zyklischer
Gruppen sind. Im folgenden erhalten wir auch eine Eindeutigkeitsaussage.
\begin{Lemma}
  Sei $G$ eine abelsche Gruppe von Primpotenzordnung, und $G=G_1G_2\dots
  G_n\cong G_1\times G_2\times\dots\times G_n$ ein direktes Produkt zyklischer
  Untergruppen $G_i$ der Ordnung $>1$. Dann sind bis auf Reihenfolge die
  Ordnungen $\abs{G_1}$, $\abs{G_2}$, \dots, $\abs{G_n}$ eindeutig gegeben.
\end{Lemma}
\begin{proof}[Beweis]
  Sei $\abs{G}$ eine Potenz von $p$. Wir zeigen die Aussage durch
  vollst\"andige Induktion \"uber die Gruppenordnung. Sei $N=\{g\in
  G|\,g^p=1\}$. Man sieht sofort, dass $N$ eine Untergruppe ist. Ferner
  besteht $N$ aus genau den Elementen $g_1g_2\dots g_n$ mit $g_i\in G_i$,
  f\"ur die $g_i^p=1$ gilt. Insbesondere ist $N$ das direkte Produkt der
  Gruppen $N_i:=G_i\cap N$, und da die Gruppen $G_i$ zyklisch sind, gilt
  $\abs{N_i}=p$ f\"ur alle $i$.
  
  Der nat\"urliche Epimorphismus $G\to\bigoplus G_i/N_i$, der
  $(g_1,\dots,g_n)$ auf $(N_1g_1,\dots,N_ng_n)$ abbildet, hat Kern $N$. Daher
  gilt $G/N\cong\bigoplus G_i/N_i$. Nach Induktionsannahme sind die Ordnungen
  $\abs{G_i/N_i}$ eindeutig gegeben, und wegen $\abs{G_i}=p\abs{G_i/N_i}$ gilt
  das dann auch f\"ur die Ordnungen $\abs{G_i}$.
\end{proof}
Wir fassen die Ergebnisse zusammen
\begin{Satz}
  Eine endliche abelsche Gruppe ist isomorph zu einem direkten Produkt
  zyklischer Gruppen von Primpotenzordnung, und die Ordnungen dieser
  zyklischer Gruppen sind bis auf Reihenfolge eindeutig gegeben.
\end{Satz}
\paragraph{Beispiele}
\begin{itemize}
\item Um (bis auf Isomorphie) die abelschen Gruppen der Ordnung $8$ zu
  bestimmen m\"ussen wir alle Darstellungen $\bigoplus_{i=1}^rC_{m_i}$
  betrachten, wo $C_{m_i}$ eine zyklische Gruppe der Ordnung $m_i$ ist und
  $m_1m_2\dots m_r=8$ gilt. Da es auf die Reihenfolge nicht ankommt, k\"onnen
  wir ferner $m_1\le m_2\le m_r$ annehmen. Wir sehen, dass es f\"ur diese
  Tupel $(m_1,m_2,\dots,m_r)$ genau die M\"oglichkeiten $(2,2,2)$, $(2,4)$,
  und $(8)$ gibt. Bis auf Isomorphie gibt es also genau $3$ abelsche Gruppen
  der Ordnung $8$.
\item Ist $n$ quadratfrei, dann ist die zyklische Gruppe der Ordnung $n$ die
  einzige abelsche Gruppe dieser Ordnung. (Das ist auch schon ohne obigen
  Hauptsatz klar.)
\end{itemize}
\begin{Bemerkung*}
  Der Beweis des Struktursatzes f\"ur abelsche Gruppen (oder auch obiger Satz)
  liefert eine andere Form der Darstellung abelscher Gruppen: Jede endliche
  abelsche Gruppe ist isomorph zu einem direkten Produkt zyklischer Gruppen
  der Ordnungen $n_1$, $n_2$, \dots, $n_k$ mit $n_i>1$, so dass $n_i$ ein
  Teiler von $n_{i+1}$ ist f\"ur alle $n_i$. Diese $n_i$ sind durch die Gruppe
  eindeutig gegeben.
\end{Bemerkung*}
\subsection{Der Satz von Jordan-H\"older}
Ist $N$ ein Normalteiler der Gruppe $G$, so entsprechen die Normalteiler von
$G/N$ bijektiv den Normalteilern $M$ von $G$ mit $M\ge N$. Wir nennen $N$
einen maximalen Normalteiler von $G$, wenn $N<G$ normal in $G$ ist, und es
keinen Normalteiler von $G$ gibt, der echt zwischen $N$ und $G$ liegt. Somit
ist $N$ genau dann ein maximaler Normalteiler von $G$, wenn $G/N$ einfach
ist. Ist $N$ ein maximaler Normalteiler von $G$, dann kann man in $N$ wieder
einen maximalen Normalteiler von $N$ w\"ahlen, und das Verfahren
fortsetzen. Dies f\"uhrt auf den folgenden Begriff.
\begin{Definition}
  Sei $G$ eine Gruppe, und $G=G_0>G_1>G_2>\dots>G_n=\{e\}$ eine Kette von
  Untergruppen, so dass $G_{i+1}$ ein maximaler Normalteiler von $G_i$
  ist. Eine solche Kette hei\ss t \ei{Kompositionsreihe}, und die einfachen
  Gruppen $G_i/G_{i+1}$ hei\ss en \ei{Kompositionsfaktoren}.
\end{Definition}
\begin{Bemerkung*}
  Ist $G$ endlich, dann existiert trivialerweise eine Kompositionsreihe. F\"ur
  unendliche Gruppen muss es das aber nicht geben. Man \"uberlege sich, dass
  z.B.\ $(\ZZZ,+)$ keine Kompositionsreihe besitzt.
\end{Bemerkung*}
Das Beispiel $C_6>C_3>\{e\}$ und $C_6>C_2>\{e\}$ zeigt schon, dass
Kompositionsreihen und die Reihenfolge der Kompositionsfaktoren nicht
eindeutig sein m\"ussen. Die symmetrische Gruppe $S_3$ hat die
Kompositionsreihe $S_3>A_3>\{e\}$, also die gleichen Kompositionsfaktoren wie
$C_6$. Hieran sehen wir, dass die Kompositionsfaktoren den Isomorphietyp einer
Gruppe nicht festlegen. Sie liefern aber dennoch wertvolle Informationen.
Wenigstens h\"angen aber die Kompositionsfaktoren einer Gruppe nicht
wesentlich von der gew\"ahlten Kompositionsreihe ab, wie der folgende Satz
zeigt.
\begin{Satz}[Jordan-H\"older]
  Sei $G$ eine endliche Gruppe. Dann haben alle Kompositionsreihen von $G$ die
  gleiche L\"ange, und die Kompositionsfaktoren stimmen bis auf Reihenfolge
  und Isomorphie \"uberein.
\end{Satz}
\begin{proof}[Beweis]
  Wir zeigen die Behauptung durch vollst\"andige Induktion \"uber die Ordnung
  von $G$. Seien $G=G_0>G_1>G_2>\dots>G_m=\{e\}$ und
  $G=H_0>H_1>H_2>\dots>H_n=\{e\}$ Kompositionsreihen. Ist $G_1=H_1$, so folgt
  die Behauptung durch Induktion.
  
  Sei also von nun an $G_1\ne H_1$. Dann ist $G_1H_1$ ein Normalteiler von
  $G$, der die maximalen und verschiedenen Normalteiler $G_1$ und $H_1$
  enth\"alt. Es folgt $G=G_1H_1$.
  
  Setze $U=G_1\cap H_1$, und $U>U_1>U_2>\dots>U_k=\{e\}$ eine
  Kompositionsreihe von $U$.

Wegen
\[
G_1/U=G_1/G_1\cap H_1\cong G_1H_1/H_1=G/H_1
\]
ist $U$ ein maximaler Normalteiler von $G_1$, insbesondere ist
$G_1>U>U_1>U_2>\dots>U_k=\{e\}$ eine Kompositionsreihe von $G_1$. Aber $G_1$
besitzt auch die Kompositionsreihe $G_1>G_2>\dots>G_m=\{e\}$. Nach
Induktionsannahme stimmen L\"ange und Kompositionsfaktoren dieser beiden
Kompositionsreihen \"uberein.

Analog gilt die Behauptung auch f\"ur die Kompositionsreihen
$H_1>U>U_1>U_2>\dots>U_k=\{e\}$ und $H_1>H_2>\dots>H_n=\{e\}$. Es folgt
zun\"achst $m=n$. Oben sahen wir bereits $G/H_1\cong G_1/U$, und analog
$G/G_1\cong H_1/U$, und daraus folgt schlie\ss lich die Behauptung, wobei
zur besseren Orientierung folgendes Diagramm dient:
\dgARROWLENGTH=0.7em
\[
\begin{diagram}
  \node[3]{G_2} \arrow{e,-} \node{G_3} \arrow{e,-} \node{\dots} \arrow{e,-}
  \node{G_m=\{e\}}\\
\node[2]{G_1} \arrow{ne,-} \arrow{se,-}\\
\node{G} \arrow{ne,-} \arrow{se,-} \node[2]{U} \arrow{e,-} \node{U_1} \arrow{e,-} \node{\dots} \arrow{e,-} \node{U_k=\{e\}}\\
\node[2]{H_1} \arrow{ne,-} \arrow{se,-}\\
\node[3]{H_2} \arrow{e,-} \node{H_3} \arrow{e,-} \node{\dots} \arrow{e,-}
  \node{H_m=\{e\}}
\end{diagram}
\]
\end{proof}
\paragraph{Beispiele von Kompositionsfaktoren}
\begin{itemize}
\item Ist $G$ eine abelsche Gruppe der Ordnung $\prod p_i^{e_i}$, dann
  bestehen die Kompositionsfaktoren aus den zyklischen Gruppen der Ordnung
  $p_i$, wobei die Gruppe der Ordnung $p_i$ genau $e_i$ Male auftritt.
\item F\"ur eine Primzahl $p$ sei $G$ eine Gruppe der Ordnung $p^m>1$. Wir
  behaupten, dass alle Kompositionsfaktoren Ordnung $p$ haben (und damit
  zyklisch sind). Dazu muss man lediglich sehen, dass eine einfache Gruppe mit
  $p$-Potenzordnung $>1$ zyklisch der Ordnung $p$ ist. Wegen Korollar
  \ref{K:Zentrum_p-Gruppe} hat eine solche Gruppe ein nichtriviales
  Zentrum. Eine Untergruppe der Ordnung $p$ des Zentrums ist aber ein
  Normalteiler der Gruppe, und die Behauptung folgt.
\item Ist die Gruppe $G$ ein direktes Produkt einfacher Gruppen $G_1$, $G_2$,
  \dots, $G_n$, dann sind diese einfachen Gruppen gerade die
  Kompositionsfaktoren von $G$. (Man beweise diese Aussage!)
\end{itemize}
In der Theorie der L\"osbarkeit von Polynomen spielt eine wichtige Eigenschaft
gewisser Gruppen eine gro\ss e Rolle.
\begin{Definition}
  Sei $G$ eine Gruppe. Elemente der Form $a^{-1}b^{-1}ab$ f\"ur $a,b\in G$
  hei\ss en \ei{Kommutatoren}. Die von den Kommutatoren erzeugte Gruppe $G'$
  hei\ss t die \ei{Kommutatorgruppe} von $G$. Die \ei{h\"oheren
    Kommutatorgruppen} $G^{(i)}$ werden rekursiv durch $G^{(0)}=G$,
  $G^{(i+1)}=(G^{(i)})'$ definiert. Die Gruppe $G$ hei\ss t \ei{aufl\"osbar},
  wenn es ein $n\in\NNN$ gibt mit $G^{(n)}=\{e\}$.
\end{Definition}
Sind $A$ und $B$ Normalteiler einer Gruppe $G$, so dass $G/A$ und $G/B$
abelsch sind, dann ist $G/A\cap B$ abelsch, da diese Gruppe isomorph ist zu
einer Untergruppe von $G/A\times G/B$. Daher gibt es genau einen kleinsten
Normalteiler $N$ von $G$, so dass $G/N$ abelsch ist.

Die Kommutatorgruppe $G'$ ist nicht nur normal in $G$, sondern hat eine
st\"arkere Eigenschaft: Es gilt $(G')^\sigma=G'$ f\"ur jeden Automorphismus
von $G$. Das liegt daran, dass ein Automorphismus die Menge der Kommutatoren
permutiert, also das Erzeugnis der Kommutatoren nicht \"andert. Induktiv
folgt, dass auch die h\"oheren Kommutatorgruppen unter Automorphismen $\sigma$
invariant, also insbesondere auch normal in $G$ sind.

Die wichtigste Eigenschaft der Kommutatorgruppe ist
\begin{Satz}
  Sei $G$ eine Gruppe. Dann ist $G'$ der kleinste Normalteiler $N$ von $G$,
  f\"ur den $G/N$ abelsch ist.
\end{Satz}
\begin{proof}[Beweis]
  Sei $N$ der kleinste Normalteiler mit $G/N$ abelsch. Sei $\phi:G\to G/N$ der
  nat\"urliche Epimorphismus. Wegen
  $(a^{-1}b^{-1}ab)^\phi=a^{-\phi}b^{-\phi}a^\phi b^\phi=e$ ist $(G')^\phi$
  die triviale Gruppe, also $G'\le N$. Umgekehrt liegt jedes Element
  $a^{-1}b^{-1}ab$ in $G'$, d.h.\ modulo $G'$ ist $G$ abelsch, und somit $N\le
  G'$. Die Behauptung folgt.
\end{proof}
\begin{Lemma}\label{L:GNaufl}
  Sei $N$ ein Normalteiler von $G$. Dann ist $G$ genau dann aufl\"osbar, wenn
  $N$ und $G/N$ aufl\"osbar sind.
\end{Lemma}
\begin{proof}[Beweis]
  Ist $\phi:G\to H$ ein Epimorphismus, dann bestehen die Kommutatoren von $H$
  gerade aus den Bildern der Kommutatoren von $G$, und somit gilt
  $(G')^\phi=H'$. Per Induktion folgt daraus $(G^{(n)})^\phi=H^{(n)}$ f\"ur
  alle $n\in\NNN$. Speziell f\"ur den Epimorphismus $G\to G/N$ folgt daraus
  $G^{(n)}N/N=(G/N)^{(n)}$.

Ist nun $G$ aufl\"osbar, so sieht man, dass auch $G/N$ aufl\"osbar ist. Wegen
$N^{(n)}\le G^{(n)}$ ist dann auch $N$ aufl\"osbar.

Seien nun $G/N$ und $N$ aufl\"osbar. Wegen der Aufl\"osbarkeit von $G/N$ folgt
zun\"achst die Existenz von $n$ mit $G^{(n)}\le N$. Daraus folgt dann induktiv
$G^{(n+k)}\le N^{(n+k)}$ f\"ur $k\ge0$. Aus der Aufl\"osbarkeit von $N$ ergibt
sich dann die Behauptung.
\end{proof}
Bei endlichen Gruppen l\"asst sich die Aufl\"osbarkeit an den
Kompositionsfaktoren ablesen. 
\begin{Satz}
  Sei $G$ eine endliche Gruppe. Dann ist $G$ genau dann aufl\"osbar, wenn alle
  Kompositionsfaktoren von $G$ zyklisch (von Primzahlordnung) sind.
\end{Satz}
\begin{proof}[Beweis]
  Ist $G$ aufl\"osbar, so sind die Quotienten aufeinanderfolgender Terme in
  der Kommutatorreihe abelsch. Insbesondere l\"asst sich die Kommutatorreihe
  durch Einf\"ugen weiterer Untergruppen zu einer Kompositionsreihe mit
  zyklischen Quotienten verfeinern.
  
  Die Umkehrung folgt unmittelbar per Induktion aus Lemma \ref{L:GNaufl}.
\end{proof}
\paragraph{Beispiele aufl\"osbarer Gruppen}
\begin{itemize}
\item Abelsche Gruppen sind aufl\"osbar.
\item Gruppen von Primpotenzordnung sind aufl\"osbar.
\item Homomorphe Bilder und direkte Produkte aufl\"osbarer Gruppen sind
  aufl\"osbar.
\end{itemize}
\begin{Bemerkung*}
  F\"ur $n\ge5$ sind die Gruppen $A_n$ und $S_n$ wegen der Einfachheit von
  $A_n$ nicht aufl\"osbar. Man \"uberlege sich, dass f\"ur alle $n\in\NNN$
  stets $S_n'=A_n$ gilt.
\end{Bemerkung*}
\subsection{Die S\"atze von Sylow}
Sei $p$ eine Primzahl. Eine endliche Gruppe hei\ss t \ei{$p$-Gruppe}, wenn die
Ordnung der Gruppe eine Potenz von $p$ ist (evtl.\ auch $1$). Wir kommen nun
zu einer wichtigen Klasse von Untergruppen.
\begin{Definition}
  Sei $G$ eine endliche Gruppe, und $p$ eine Primzahl. Eine Untergruppe $H$
  von $G$ hei\ss t \ei{$p$-Sylowgruppe} von $G$, wenn $\abs{U}$ eine
  $p$-Gruppe ist, und $[G:U]$ nicht durch $p$ teilbar ist.
\end{Definition}
Ist also $p^r$ die h\"ochste Potenz von $p$, die $\abs{G}$ teilt, dann sind
die $p$-Sylowgruppen gerade die Untergruppen der Ordnung $p^r$. F\"ur abelsche
Gruppen haben wir gesehen, dass es f\"ur jedes $p$ genau eine $p$-Sylowgruppe
gibt. Der wichtigste Satz in diesem Zusammenhang liefert die Existenz von
Sylowgruppen auch f\"ur nicht abelsche Gruppen.
\begin{Satz}[Sylow]\label{S:Sylow}
  Sei $G$ eine endliche Gruppe, und $p$ eine Primzahl. Dann besitzt $G$ eine
  $p$-Sylowgruppe.
\end{Satz}
F\"ur diesen wichtigen Satz geben wir drei Beweise.
\begin{proof}[1.\ Beweis]
  Sei $Z$ das Zentrum von $G$, und $g_1,g_2,\dots,g_h$ Repr\"asentanten der
  Konjugationsklassen. Die Klassengleichung liefert
\[
\abs{G}=\sum_{i=1}^h[G:C_G(g_i)]=\abs{Z}+\sum_{g_i\notin Z}[G:C_G(g_i)].
\]
Man beachte, dass in der zweiten Summe die Gruppen $C_G(g_i)$ echte
Untergruppen von $G$ sind. Nach Induktion besitzen diese Gruppen
$p$-Sylowgruppen. Ist einer der Indizes $[G:C_G(g_i)]$ teilerfremd zu $p$,
dann ist eine $p$-Sylowgruppe von $C_G(g_i)$ auch eine von $G$. Wir sind also
fertig, au\ss er wenn alle Indizes durch $p$ teilbar sind. In diesem Fall aber
ist $\abs{Z}$ durch $p$ teilbar. $Z$ besitzt also eine Untergruppe $N$ der
Ordnung $p$, diese ist normal in $G$. Nach Induktion besitzt $G/N$ eine
$p$-Sylowgruppe $P/N$, aber dann ist $P$ eine $p$-Sylowgruppe von $G$.
\end{proof}
F\"ur den zweiten Beweis ben\"otigen wir ein kleines
\begin{Lemma}
  Sei $n\in\NNN$, und $p^r$ die h\"ochste Potenz der Primzahl $p$, die $n$
  teilt. Dann ist der Binomialkoeffizient $\binom{n}{p^r}$ nicht durch $p$
  teilbar.
\end{Lemma}
\begin{proof}[Beweis]
  Schreibe $n=p^rm$. Somit ist $m$ nicht durch $p$ teilbar.
  $\binom{n}{p^r}=\binom{p^rm}{p^r}$ ist das Produkt der rationalen Zahlen
  $\frac{p^rm-k}{p^r-k}$ f\"ur $k=0,1,\dots,p^r-1$. Wir zeigen, dass in der
  gek\"urzten Darstellung Z\"ahler und Nenner nicht durch $p$ teilbar sind.
  Das ist klar f\"ur $k=0$. Sei daher $0<k=p^le$ mit $e$ nicht durch $p$
  teilbar. Wegen $k<p^r$ gilt $l<r$, daher ist $p^{r-l}$ durch $p$ teilbar und
  die Behauptung folgt aus
  $\frac{p^rm-k}{p^r-k}=\frac{p^{r-l}m-e}{p^{r-l}-e}$.
\end{proof}
\begin{proof}[2.\ Beweis von Satz \ref{S:Sylow}]
  Sei $p^r$ wieder die h\"ochste $p$-Potenz, die $\abs{G}$ teilt. Sei $M$ die
  Menge der $p^r$-elementigen Teilmengen von $G$. Die M\"achtigkeit
  $\abs{M}=\binom{\abs{G}}{p^r}$ ist nach obigem Lemma nicht durch $p$
  teilbar. Wir lassen nun $G$ von rechts auf $M$ operieren, $g\in G$ schickt
  also die Teilmenge $S\in M$ nach $Sg$. Da $\abs{M}$ nicht durch $p$ teilbar
  ist, muss eine Bahn dabei sein, deren L\"ange nicht durch $p$ teilbar ist.
  Sei $S\in M$ in dieser Bahn, und $P$ der Stabilisator von $S$. Die Bahn
  durch $S$ hat L\"ange $[G:P]$, und da das nicht durch $p$ teilbar ist, ist
  $\abs{P}$ durch $p^r$ teilbar. Sei $s\in S$. Die Abbildung $P\to S$,
  $x\mapsto sx$ ist injektiv, daher gilt $\abs{P}\le\abs{S}=p^r$, und somit
  $\abs{P}=p^r$, und die Behauptung folgt.
\end{proof}
\begin{proof}[3.\ Beweis von Satz \ref{S:Sylow}]
  Diesen Beweis zerlegen wir in eine Serie von \"Ubungsaufgaben.
  \begin{itemize}
  \item[(a)] Sei $S$ eine endliche Gruppe mit $p$-Sylowgruppe $P$, und $G$
    eine Untergruppe von $S$. Dann gibt es ein $s\in S$, so dass $G\cap P^s$
    eine $p$-Sylowgruppe von $G$ ist. Hierzu betrachte man die Operation von
    $G$ auf dem Nebenklassenraum $P\setminus S$.
    
    Um den Satz zu beweisen, gen\"ugt es also, zu unserer gegebeben Gruppe $G$
    eine Gruppe $S$ zu finden, die eine zu $G$ isomorphe Untergruppe
    enth\"alt, und f\"ur die es eine $p$-Sylowgruppe $P$ gibt. Das wird in den
    folgenden Schritten erreicht.
    
  \item[(b)] Die Gruppe $G$ ist isomorph zu einer Untergruppe der
    symmetrischen Gruppe $S_n$. Sei $\FFF_p=\ZZZ/p\ZZZ$ der K\"orper mit $p$
    Elementen. In $\GL_n(\FFF_p)$ betrachte man die Menge $U$ der Matrizen,
    die in jeder Zeile und jeder Spalte genau eine $1$, und sonst nur $0$er
    besitzen. Zeige, dass $U$ eine zu $S_n$ isomorphe Untergruppe von
    $\GL_n(\FFF_p)$ ist. Somit hat $\GL_n(\FFF_p)$ eine zu $G$ isomorphe
    Untergruppe.
  
  \item[(c)] Zeige, dass die Gruppe der echten oberen Dreiecksmatrizen von
    $\GL_n(\FFF_p)$ eine $p$-Sylowgruppe ist durch direkte Berechnung der
    Ordnung dieser Gruppe und der h\"ochsten $p$-Potenz, die die Ordnung von
    $\GL_n(\FFF_p)$ teilt.
  \end{itemize}
\end{proof}
Bevor wir zu weiteren Aussagen von Sylow kommen, wollen wir eine einfache
Eigenschaft von $p$-Gruppen beweisen.
\begin{Lemma}\label{L:p-GruppeIndp}
  Sei $U$ eine echte Untergruppe der $p$-Gruppe $G$. Dann gilt
  $N_G(U)>U$. Ferner hat $G$ f\"ur jeden Teiler $m$ der Gruppenordnung eine
  Untergruppe der Ordnung $m$, und Untergruppen vom Index $p$ sind normal.
\end{Lemma}
\begin{proof}[Beweis]
  $U$ operiere auf dem Nebenklassenraum $U\setminus G$. Da $U$ eine $p$-Gruppe
  ist, hat jede Bahn eine L\"ange, die eine Potenz von $p$ ist. Ferner ist die
  Summe $\abs{U\setminus G}$ durch $p$ teilbar. Da die Bahn durch $U$ L\"ange
  $1$ hat, muss es eine weitere Bahn $Ux\ne U$ der L\"ange $1$ geben. Es gilt
  also $Uxu=Ux$ f\"ur alle $u\in U$, und daraus folgt $U^x=U$, d.h.\ $x$ liegt
  nicht in $U$, normalisiert aber $U$.
  
  Die letzte Aussage ist aus Ordnungsgr\"unden klar, und die Existenz der
  Untergruppen der Ordnung $m$ folgt induktiv aus der ersten Aussage, oder
  alternativ aus der bereits gesehenen Tatsache, dass die Kompositionsfaktoren
  von $G$ alle Ordnung $p$ haben.
\end{proof}
Das Hauptergebnis von Sylow, das wir im folgenden Abschnitt st\"andig
verwenden werden, ist
\begin{Satz}[Sylow]\label{S:Sylow2}
  Sei $G$ eine endliche Gruppe, und $p$ eine Primzahl. Dann gilt:
  \begin{itemize}
  \item[(a)] Jede $p$-Untergruppe von $G$ liegt in einer $p$-Sylowgruppe von
  $G$.
\item[(b)] Die $p$-Sylowgruppen von $G$ sind konjugiert.
\item[(c)] Die Anzahl der $p$-Sylowgruppen ist von der Form $1+kp$ mit
  $k\in\NNN_0$, und gleich $[G:N_G(P)]$, wobei $P$ eine $p$-Sylowgruppe von
  $G$ ist.
  \end{itemize}
\end{Satz}
\begin{proof}[Beweis]
  Sei $P$ eine $p$-Sylowgruppe von $G$, und $S$ die Menge der zu $P$
  konjugierten Untergruppen von $P$. Nat\"urlich ist jede Gruppe aus $S$
  wieder eine $p$-Sylowgruppe. $G$ operiert auf $S$ durch Konjugation. Wegen
  $\abs{S}=[G:N_G(P)]$ ist die $\abs{S}$ nicht durch $p$ teilbar. Sei $U\le G$
  eine $p$-Gruppe. Auch $U$ operiert durch Konjugation auf $S$. Da jede
  Bahnl\"ange von $U$ eine $p$-Potenz ist, und $\abs{S}$ nicht durch $p$
  teilbar ist, muss $U$ eine Bahn $\{Q\}$ der L\"ange $1$ haben. Dann gilt
  $uQ=Qu$ f\"ur alle $u\in U$, insbesondere ist $QU$ eine Gruppe mit
  Normalteiler $Q$. Wegen $QU/Q\cong U/Q\cap U$ ist $QU$ eine $p$-Gruppe der
  Ordnung $\abs{Q}[U:Q\cap U]$. Aber $Q$ ist schon eine $p$-Untergruppe von
  $G$ gr\"o\ss tm\"oglicher Ordnung, daher gilt $U=Q\cap U$, also $U\le Q$.
  Hieraus folgt (a), und damit auch (b).

Um (c) zu zeigen, verwenden wir obiges Argument mit $U=P$. Ein Fixpunkt von
$P$ ist $P$ selber. Weitere Fixpunkte $Q$ kann es nicht geben, denn wie oben
folgt $P\le Q$, also $P=Q$ aus Ordnungsgr\"unden. Alle anderen Bahnen von $p$
haben also eine L\"ange der Form $p^r$ mit $r\ge1$, insbesondere sind diese
L\"angen durch $p$ teilbar. Daher hat $\abs{S}$ die Form $1+kp$.

Der Stabilisator von $P$ in $G$ ist $N_G(P)$, daher gilt $\abs{S}=[G:N_G(P)]$,
und alles ist gezeigt.
\end{proof}
\subsection{Gruppen kleiner Ordnung}
In diesem Abschnitt sollen examplarisch einige Anwendungen der Sylow-S\"atze
und der Techniken der Gruppenoperationen vorgestellt werden, die dem Studium
von Gruppen dienen, die klein sind oder deren Ordnungen nur wenige Teiler
besitzen.

Ist $p$ eine Primzahl, dann sind die Gruppen der Ordnung $p$ nat\"urlich
zyklisch. Eine Gruppe der Ordnung der Ordnung $p^2$ hat ein nicht triviales
Zentrum, wird also vom Zentrum und einem weiteren Element erzeugt, und ist
somit abelsch. Daher ist jede Gruppe der Ordnung $p^2$ isomorph zu $C_{p^2}$
oder $C_p\times C_p$.

Wenn man in einer Gruppe $G$ mit wenigen Teilern von $\abs{G}$ nach
Normalteilern sucht, dann ist h\"aufig folgendes Vorgehen erfolgreich:

Man nimmt den gr\"o\ss ten Primteiler $p$ von $\abs{G}$, oder die Primzahl $p$
mit der gr\"o\ss ten $p$-Sylowgruppe $P$. Nun versucht man mittels der
Teilbarkeitseigenschaften aus den S\"atzen von Sylow zu zeigen, dass $P$
normal in $G$ ist. Gelingt das nicht, sollte man andere Primzahlen
probieren. Funktioniert auch das nicht, so kann man versuchen, die
Nichtnormalit\"at einer $p$-Sylowgruppe $P$ dadurch zu einem Widerspruch zu
f\"uhren, dass $P$ zu viele Konjuguierte hat. Ich n\"amlich $P$ nicht normal,
dann hat $P$ mindestens $1+p$ Konjugierte. Hat beispielsweise $P$ die Ordnung
$p$, dann schneiden sich die Konjugierten trivial, man erh\"alt also
mindestens $(p-1)(p+1)$ Elemente der Ordnung $p$. Wenn man das auch noch f\"ur
andere Sylowgruppen macht, kann es beispielsweise passieren, dass man durch
solche Z\"ahlungen mehr Elemente bekommt als die Gruppenordnung ist.

Hilft auch das nicht weiter, kann man probieren, kleine
Permutationsoperationen zu finden. Ist z.B.\ $P$ nicht normal, dann erh\"alt
man einen nicht trivialen Homomorphismus $G\to\text{Sym}(P\setminus G)$. Oft
muss aus Anzahl- oder Teilbarkeitsgr\"unden ein nicht-trivialer Kern vorhanden
sein. Manchmal wei\ss\ man schon, dass $N_G(P)>P$ gilt, dann bekommt man meist
bessere Aussagen, wenn man $G\to\text{Sym}(N_G(P)\setminus G)$ betrachtet.

Nat\"urlich muss man sich oft auch, abh\"angig von der Situation, andere
Tricks einfallen lassen. Beispielsweise kann man aus der Klassengleichung oft
n\"utzliche Informationen gewinnen.

Die folgenden Beispiele sollen einige dieser Prinzipien verdeutlichen.

\subsubsection{$\abs{G}=pq$, $p<q$ Primzahlen.}\label{sss:A}
Sei $P$ und $Q$ eine $p$-Sylow- bzw.\ $q$-Sylowgruppe von $G$, und $n_p$ und
$n_q$ die Anzahl der $p$-Sylow- bzw.\ $q$-Sylowgruppen. Ist $Q$ nicht
normal in $G$, dann gilt $N_G(Q)=Q$, es folgt
$n_q=[G:N_G(Q)]=[G:Q]=\frac{pq}{q}=p$. Anderseits ist $q$ ein Teiler von
$n_q-1=p-1$, also insbesondere $q<p$, im Widerspruch zu $p<q$.

Daher ist $Q$ ein Normalteiler, und $G$ ein semidirektes Produkt von $Q$ mit
$P$.

Ist $P$ auch normal in $G$, dann ist $G$ ein direktes Produkt von $P$ und $Q$,
also $G\cong C_p\times C_q\cong C_{pq}$.

Allerdings, wie das Beispiel $S_3$ schon zeigt, muss $P$ nicht normal sein. In
diesem Fall folgt analog wie oben, dass $p$ ein Teiler von $q-1$ sein
muss. Wenn also $p$ kein Teiler von $q-1$ ist, dann ist $G\cong C_{pq}$.

Als n\"achstes wollen wir den Fall von $3$ verschiedenen Primfaktoren
studieren. Zur Einstimmung erst mal ein einfaches Beispiel.

\subsubsection{$\abs{G}=1001=7\cdot 11\cdot 13$.}
Seien $n_7$, $n_{11}$ und $n_{13}$ die Anzahlen der $7$-, $11$- und
$13$-Sylowgruppen. Nun ist $n_{13}$ ein Teiler von $7\cdot11$, also gleich
$1,7,11$ oder $77$. Aber au\ss er $1$ hat keine dieser Zahlen die Form
$1+13k$, es folgt $n_{13}=1$. Analog ist $n_7$ gleich $1,11,13$ oder $11\cdot
13=143$ und von der Form $1+7k$, also $n_7=1$, und genauso folgt
$n_{11}=1$. Damit sind alle Sylowgruppen normal, und $G$ ist somit ein
direktes Produkt dieser zyklischen Normalteiler mit teilerfremden
Ordnungen. Daher ist $G$ zyklisch.

\subsubsection{$\abs{G}=pqr$, $p<q<r$ Primzahlen.}
Seien nun $P$, $Q$ und $R$ Untergruppen der Ordnungen $p$, $q$, und $r$.
Nat\"urlich sind $P$, $Q$ und $R$ Sylowgruppen. F\"ur eine Primzahl $s$ sei
$n_s$ wieder die Anzahl der $s$-Sylowgruppen.

\begin{comment}
Zun\"achst wollen wir sehen, dass $R$ normal ist. Das sei also nicht der Fall.
Wegen $R\le N_G(R)<G$ gilt $n_r=p,q$ oder $pq$. Die ersten beiden F\"alle
treten aber nicht auf, da $r$ ein Teiler von $n_r-1$ ist, aber $r>p-1$ und
$r>q-1$ nach Voraussetzung. Daher gilt $n_r=pq$. Somit gibt es
$pq(r-1)=pqr-pq$ Elemente der Ordnung $r$. Die Anzahl der Elemente von einer
Ordnung $\ne r$ ist daher $pq$. Die Konjugierten von $P$ und $Q$ steuern
$1+n_p(p-1)+n_q(q-1)$ Elemente der Ordnung $\ne r$ bei, daher gilt
$1+n_p(p-1)+n_q(q-1)\le pq$. Wir wollen sehen, dass $n_p=1$ oder $n_q=1$ gilt.
Ist das nicht der Fall, dann ist $n_p>1$ ein Teiler von $qr$, also $n_p\ge q$,
und analog $n_q=\ge p$. Es folgt $1+q(p-1)+p(q-1)\le pq$, also $(p-1)(q-1)\le
0$, was Unsinn ist. Damit ist $P$ oder $Q$ normal in $G$, insbesondere ist
$U=PQ$ ein semidirektes Produkt der Ordnung $pq$, und $U$ enth\"alt keine
Elemente der Ordnung $r$. Damit besteht $U$ genau aus den Elementen von $G$
der Ordnung $\ne r$, insbesondere ist $U$ normal in $G$. Nach Beispiel
\ref{sss:A} ist $Q$ normal in $U$. Insbesondere ist $Q$ die einzige
$q$-Sylowgruppe von $U$, und da $U$ normal in $G$ ist, ist $Q$ sogar normal in
$G$. Insbesondere ist $QR$ ein semidirektes Produkt, aber da $r$ kein Teiler
von $q-1$ ist, ist das Produkt direkt. Wegen $G=QRP$, sind die Konjugierten
von $R$ von der Form $R^x$ mit $x\in P$, insbesondere gilt $n_r\le p$, im
Widerspruch zu $n_r=pq$.

Damit ist $R$ als Normalteiler erkannt. Nach Beispiel \ref{sss:A} ist $QR/R$
normal in $G/R$, und somit ist $QR$ normal in $G$. Es folgt
\[
G\cong P\ltimes (Q\ltimes R).
\]
Als \"Ubung \"uberlege man sich, dass auch
\[
G\cong (P\ltimes Q)\ltimes R
\]
gilt. Dazu muss man etwas aufpassen, denn $PQ$ muss keine Gruppe
sein. Stattdessen muss man $Q$ durch eine Konjugierte ersetzen, die von $P$
normalisiert werden.

F\"uhrt man diese etwas technischen Betrachtungen noch etwas weiter, so sieht
man, dass stets $G\cong C_p\ltimes C_{qr}$ oder $G\cong C_{pq}\ltimes C_r$
gilt.
\end{comment}
Zun\"achst wollen wir sehen, dass eine der Gruppen $P$, $Q$ oder $R$ normal in
$G$ ist. Dazu nehmen wir an, das sei nicht der Fall. Wegen $R\le N_G(R)<G$
gilt $n_r=p,q$ oder $pq$. Die ersten beiden F\"alle treten aber nicht auf, da
$r$ ein Teiler von $n_r-1$ ist, aber $r>p-1$ und $r>q-1$ nach Voraussetzung.
Daher gilt $n_r=pq$. Ferner ist $1<n_p$ ein Teiler von $qr$, also $n_p\ge q$,
und analog $n_q\ge p$. Die Konjugierten der Gruppen $P$, $Q$ und $R$ schneiden
sich paarweise trivial. Daher gibt es $n_r(r-1)$ Elemente der Ordnung $r$,
$n_p(p-1)$ Elemente der Ordnung $p$ und $n_q(q-1)$ Elemente der Ordnung $q$.
Es folgt
\[
pqr-1\ge n_r(r-1)+n_p(p-1)+n_q(q-1)\ge pq(r-1)+q(p-1)+p(q-1),
\]
also $(p-1)(q-1)\le 0$, was Unsinn ist. Damit ist eine der Gruppen $P$, $Q$
oder $R$ normal in $G$.

Wir wollen sehen, dass $R$ normal in $G$ ist. Nach obigem ist $P$ oder $Q$
normal in $G$, sei z.B.\ $P$ normal in $G$. Nach Beispiel \ref{sss:A} ist dann
$PR/P$ normal in $G/P$, also $PR$ normal in $G$. Aber wiederum nach Beispiel
\ref{sss:A} ist $R$ normal in $PR$. Somit ist $R$ die einzige $r$-Sylowgruppe
in $PR$, und daher muss $R$ in $G$ normal sein.

Damit ist $R$ als Normalteiler erkannt. Nach Beispiel \ref{sss:A} ist $QR/R$
normal in $G/R$, und somit ist $QR$ normal in $G$. Es folgt
\[
G\cong P\ltimes (Q\ltimes R).
\]

F\"uhrt man diese etwas technischen Betrachtungen noch weiter, so sieht man,
dass stets $G\cong C_p\ltimes C_{qr}$ oder $G\cong C_{pq}\ltimes C_r$ gilt.
\subsubsection{$\abs{G}=p^aq^b$, $p<q$ Primzahlen, $0\le a,b\le 2$.}
Sei wieder $P$ und $Q$ eine $p$- und $q$-Sylowgruppe. Ist $Q$ normal in $G$,
dann ist $G$ ein semidirektes Produkt von $Q$ mit $P$. Wir wollen den Fall
analysieren, dass $Q$ nicht normal ist. Mit den gewohnten Bezeichnungen folgt,
dass $1<n_q$ ein Teiler von $p^2$ ist, und $q$ ein Teiler von $n_q-1$
ist. Wegen $q>p$ kann nicht $n_q=p$ gelten, daher gilt $n_q=p^2$, und $q$ ist
ein Teiler von $p^2-1=(p-1)(p+1)$. Daher ist $p-1$ oder $p+1$ durch $q$
teilbar. Wegen $q>p$ folgt $q=p+1$, also $p=2$, $q=3$.

Somit gilt $n_3=4$ und $\abs{G}=12$ oder $\abs{G}=36$. Wegen $n_3=2^2$ gilt
$N_G(Q)=Q$. Betrachte die Operation durch Konjugation auf den vier
$3$-Sylowgruppen. Wir erhalten $G\to S_4$, und der Kern ist der Schnitt der
$3$-Sylowgruppen. Im Fall $\abs{G}=12$ ist dieser Schnitt trivial, es folgt
$G\le S_4$. Aber $A_4$ ist die einzige Untergruppe der Ordnung $12$ in
$S_4$, und es folgt $G\cong A_4$. Insbesondere ist $P$ normal in $G$.

Im zweiten Fall ist der Kern von $G\to S_4$ eine Gruppe $Z$ der Ordnung $3$,
wie eben folgt $G/Z\cong A_4$. Somit ist $PZ/Z$ normal in $G/Z$, also $PZ$
normal in $G$ der Ordnung $12$. Es folgt schnell, dass $P$ normal in $G$ ist.
\begin{Bemerkung*}
  F\"ur beliebige nat\"urliche Zahlen $a,b$ gilt der Satz von Burnside, dass
  eine Gruppe der Ordnung $p^aq^b$ aufl\"osbar ist. Auch bis heute ist der
  einfachste Beweis mittels Darstellungstheorie. Hierbei untersucht man die
  m\"oglichen Homomorphismen $G\to\GL_n(\CCC)$.
\end{Bemerkung*}
\subsubsection{Aufl\"osbarkeit von $G$ f\"ur $\abs{G}<60$.}
Wir haben bereits die Gruppe $A_5$ der Ordnung $60$ kennengelernt, die einfach
aber nicht abelsch ist. Hier wollen wir sehen, dass es keine kleinere einfache
und nicht abelsche Gruppen gibt. Insbesondere sind dann alle Gruppen der
Ordnung $<60$ aufl\"osbar.

Sei $G$ einfach und nicht abelsch mit $\abs{G}<60$. Wegen der bisherigen
F\"alle bleiben nur noch die Gruppenordnungen $2\cdot 3^3$, $2^3\cdot 3$,
$2^3\cdot 5$, $2^3\cdot 7$ und $2^4\cdot 3$ zu untersuchen. $G$ kann keine
echte Untergruppe $U$ vom Index $\le4$ besitzen, denn dann bettet $G$ ein in
$\text{Sym}(U\setminus G)\le S_4$, aber $S_4$ ist aufl\"osbar. Damit bleiben
nur noch die Ordnungen $2^3\cdot 5$ und $2^3\cdot 7$. Im ersten Fall folgt aus
den Sylows\"atzen sofort, dass die $5$-Sylowgruppe normal ist, ein
Widerspruch. Bleibt $\abs{G}=8\cdot 7$. Es folgt $n_7=8$, daher hat $G$ genau
$8(7-1)=48$ Elemente der Ordnung $7$. Die Menge der restlichen $8$ Elemente
ist unter Konjugation invariant. Aber auch die $2$-Sylowgruppe von $G$ hat
Ordnung $8$, und stimmt daher mit dieser Menge \"uberein. Sie ist also normal.
\section{Ringe}
\subsection{Definitionen, Beispiele}
Ein \ei{Ring} ist eine Menge $R$ mit zwei zweistelligen Verkn\"upfungen $+$
und $\cdot$ und Elementen $0,1\in R$, so dass die folgenden Ringaxiome gelten:
\begin{itemize}
\item $(R,+)$ ist eine (additiv geschriebene) abelsche Gruppe mit neutralem
  Element $0$.
\item $(R,\cdot)$ ist ein Monoid mit neutralem Element $1$.
\item Es gelten die beiden Distributivgesetze $x\cdot(y+z)=x\cdot y+x\cdot z$
  und $(y+z)\cdot x=y\cdot x+z\cdot x$ f\"ur alle $x,y,z\in R$.
\end{itemize}
Um Klammern zu sparen gelten die \"ublichen Konventionen, dass Potenzbildung
vor Multiplikation, und diese vor Addition geht. Bei Produkten schreibt man
oft $xy$ statt $x\cdot y$.

Zur Gew\"ohnung an die Axiome beginnen wir mit einigen einfachen Aussagen:
\begin{Lemma}
  In einem Ring $R$ gelten die folgenden Aussagen:
  \begin{itemize}
  \item[(a)] $0a=a0=0$.
\item[(b)] $(-a)b=-ab=a(-b)$.
\item[(c)] $a(b-c)=ab-ac$, $(a-b)c=ac-bc$.
\item[(d)] Gilt $0=1$,  so folgt $R=\{0\}$.
\item[(e)] Gilt $ab=ba$, so folgt $(a+b)^n=\sum_{k=0}^n\binom{n}{k}a^kb^{n-k}$
  f\"ur $n\in\NNN$.
  \end{itemize}
\end{Lemma}
\begin{proof}[Beweis]
  Aus $0=0+0$ folgt $a0=a(0+0)=a0+a0$, also $a0=0$, und analog
  $0a=0$. \"Ahnlich einfach folgen die Aussagen (b) und (c). Aussage (d) folgt
  aus (a) und $a=a\cdot1=a\cdot0=0$. Der Beweis von (e) ist wie in der
  Analysis mittels vollst\"andiger Induktion, bei der Rechnung bemerkt man,
  dass man neben den Ringaxiomen noch die multiplikative Vertauschbarkeit von
  $a$ und $b$ ben\"otigt.
\end{proof}
Soweit im folgenden nichts anderes gesagt wird, sei stillschweigend $0\ne1$
vorausgesetzt.
\paragraph{Beispiele von Ringen}
\begin{itemize}
\item Die ganzen Zahlen $\ZZZ$ bilden einen Ring.
\item Ist $n\in\NNN$, so bildet die Menge der Restklassen $\ZZZ/n\ZZZ$ mit den
  Festsetzungen $(a+n\ZZZ)+(b+n\ZZZ)=(a+b)+n\ZZZ$ und
  $(a+n\ZZZ)\cdot(b+n\ZZZ)=(ab)+n\ZZZ$ einen Ring mit $n$ Elementen.
\item Ist $K$ ein K\"orper, so bildet die Menge der Matrizen $M_n(K)$ einen
  Ring.
\item F\"ur eine abelsche Gruppe $A$ ist die Menge der Endomophismen von $A$
  ein Ring, wobei $f+g$ durch $a^{f+g}:=a^f+a^g$ und $a^{f\cdot
  g}:=a^{fg}=(a^f)^g$ erkl\"art werden. (Was passiert, wenn $A$ nicht abelsch
  ist?)
\item Ist $M$ eine Menge, und $R$ ein Ring, so bildet die Menge der
  Abbildungen von $M$ nach $R$ unter elementweiser Addition und Multiplikation
  ein Ring.
\end{itemize}
Nach diesen Beispielen wollen wie einige nachfolgend h\"aufig benutzte
Begriffe definieren.
\begin{Definition}
  Ein Element $a\in R$ hei\ss t \ei{invertierbar} oder \ei{Einheit}, wenn es
  $b,c\in R$ gibt mit $ab=ca=1$. (Dann folgt automatisch $c=c1=cab=1b=b$.)
  Die Menge $R^\times$ der Einheiten bilden eine Gruppe, die
  \ei{Einheitengruppe} von $R$.

Ein Element $a\in R$ hei\ss t \ei{Nullteiler}, wenn es ein $b\ne0$ gibt mit
$ab=0$.

Ein Element $a\in R$ hei\ss t \ei{nilpotent}, wenn es ein $n\in\NNN$ gibt mit
$a^n=0$.

Ein Ring $R$ hei\ss t \ei{kommutativ}, wenn $ab=ba$ gilt f\"ur alle $a,b\in
R$.

Der Ring $R$ ($\ne\{0\}$) hei\ss t \ei{Integrit\"atsbereich} oder
\ei{Integrit\"atsring}, wenn $R$ nullteilerfrei und kommutativ ist.
  
Sind alle Elemente $\ne0$ eines Ringes ($\ne\{0\}$) Einheiten, so hei\ss t der
Ring ein \ei{Schiefk\"orper}.

Kommutative Schiefk\"orper hei\ss en \ei{K\"orper}.
\end{Definition}
Typische Integrit\"atsbereiche sind etwa die ganzen Zahlen $\ZZZ$. Die
Einheitengruppe von $\ZZZ$ besteht allerdings nur aus $\pm1$. Nat\"urlich ist
auch jeder K\"orper ein Integrit\"atsbereich.
Als einfache Anwendung der Begriffe zeigen wir
\begin{Lemma}
  Ein endlicher Integrit\"atsbereich ist ein K\"orper.
\end{Lemma}
\begin{proof}[Beweis]
  Sei $R$ ein endlicher Integrit\"atsbereich, und $a\ne0$. Wir m\"ussen
  zeigen, dass $a$ invertierbar ist. Dazu betrachten wir die Abbildung $R\to
  R$, $x\mapsto ax$. Diese Abbildung ist injektiv, denn aus $ax=ay$ folgt
  $a(x-y)=0$, also $x-y=0$, da $0$ der einzige Nullteiler von $R$ ist.  Da
  aber $R$ endlich ist, folgt aus der Injektivit\"at dieser Abbildung schon
  die Surjektivit\"at. Insbesondere gibt es $x\in R$ mit $ax=1$. Analog folgt
  die Existenz von $y$ mit $ya=1$, die Behauptung folgt.
\end{proof}
Eine wichtige Folgerung ist
\begin{Korollar}\label{K:ZnZ}
  Sei $n\in\NNN$. Dann ist $\ZZZ/n\ZZZ$ genau dann ein K\"orper, wenn $n$ eine
  Primzahl ist.
\end{Korollar}
\begin{proof}[Beweis]
  Sei $R=\ZZZ/n\ZZZ$. Im Falle $n=1$ gilt $R=\{0\}$. Sei nun $n=ab$ mit
  $a,b\in\NNN$ und $a,b>1$. Dann gilt $(a+n\ZZZ)(b+n\ZZZ)=n\ZZZ=0_R$, aber
  $a+n\ZZZ,b+n\ZZZ\ne0_R$. Damit ist $R$ kein K\"orper.
  
  Nun sei $n$ eine Primzahl. Nach obigem Lemma bleibt zu zeigen, dass $R$ ein
  Integrit\"atsbereich ist, d.h, dass $R$ keine Nullteiler au\ss er $0_R$ hat.
  Seien $a+n\ZZZ,b+n\ZZZ\ne0_R$ mit $(a+n\ZZZ)(b+n\ZZZ)=ab+n\ZZZ=0_R$. Damit
  ist $n$ ein Teiler von $ab$, und da $n$ prim ist, gilt z.B., dass $n$ ein
  Teiler von $a$ ist. Aber dann gilt $a+n\ZZZ=0_R$, ein Widerspruch.
\end{proof}
Ein \ei{Teilring} eines Ringes $R$ ist eine Teilmenge $S$ von $R$ mit $0,1\in
S$, so dass $S$ mit den von $R$ auf $S$ eingeschr\"ankten Operationen wieder
einen Ring bildet. Hierzu gen\"ugt es zu zeigen, dass mit $s,t\in S$ auch
$s-t$ und $st$ wieder in $S$ liegen. Ist $M$ eine Teilmenge von $R$, und $R_0$
ein Teilring von $R$, dann bezeichnet $R_0[M]$ den von $R_0$ und $M$ erzeugten
Teilring von $R$. Das ist der kleinste Teilring von $R$, der $R_0$ und $M$
enth\"alt.
\subsection{Homomorphismen und Ideale}
\begin{Definition}
  Eine Abbildung $\phi:R\to S$ zwischen Ringen $R$ und $S$ hei\ss t
  \ei{Ringhomomorphismus}, wenn $1_R^\phi=1_S$ gilt, und
  $(x+y)^\phi=x^\phi+y^\phi$ und $(xy)^\phi=x^\phi y^\phi$ gelten f\"ur alle
  $x,y\in R$.
\end{Definition}
\"Ahnlich einfach wie bei Gruppen rechnet man nach, dass $R^\phi$ ein Teilring
von $S$ ist, und dass Urbilder von Teilringen von $S$ wieder Teilringe von $R$
sind. Die Kerne von Gruppenhomomorphismen waren gerade die Normalteiler von
Gruppen. Der Kern eines Ringhomomorphismus $R\to S$ ist nat\"urlich eine
Untergruppe von $(R,+)$, aber es gilt eine weitere Eigenschaft.
\begin{Definition}
  Eine Untergruppe $I$ von $(R,+)$ hei\ss t \ei{Ideal} von $R$, wenn f\"ur
  alle $r\in R$, $i\in I$ gilt: $ri, ir\in I$. Man schreibt dann auch
  $I\trianglelefteq R$.
\end{Definition}
Sei $I$ der Kern eines Ringhomomorphismus $\phi:R\to S$. Wie schon bemerkt,
ist $I$ eine Untergruppe von $(R,+)$. Sei $r\in R$, $i\in I$. Wegen
$(ri)^\phi=r^\phi i^\phi=r^\phi 0=0$ gilt $ri\in I$, und analog $ir\in
I$. Damit ist $I$ ein Ideal. In K\"urze werden wir sehen, dass umgekehrt jedes
Ideal der Kern eines geeigneten Homomorphismus ist.

Man beachte, dass Ideale im allgemeinen keine Teilringe sind. Gilt n\"amlich
$1\in R$, so folgt bereits $I=R$.

Ist $M$ eine Teilmenge von $R$, so bildet die Menge der endlichen Summen $\sum
r_i m_i s_i$ mit $r_i,s_i\in R$, $m_i\in M$ offenbar das kleinste Ideal von
$R$, welches $M$ enth\"alt. Man nennt das auch das von $M$ erzeugte Ideal. Ein
\ei{Hauptideal} ist ein von einem einzigen Element erzeugtes Ideal. Ist $R$
kommutativ, dann hat das von $m$ erzeugte Hauptideal die Form $mR$. Man
schreibt daf\"ur auch h\"aufig $(m)$. Ist jedes Ideal eines Ringes ein
Haupideal, dann nennt man den Ring \ei{Hauptidealring}.

Wir wissen bereits, dass die Untergruppen von $(\ZZZ,+)$ alle die Form $n\ZZZ$
mit $n\in\NNN_0$ haben. Insbesondere haben alle Ideale von $\ZZZ$ diese Form,
somit ist $\ZZZ$ ein Hauptidealring.

Sei nun $I$ ein Ideal des Rings $R$. Da $(I,+)$ eine normale Untergruppe von
$(R,+)$ ist, erhalten wir eine nat\"urliche Abbildung $\phi:R\to R/I$ via
$x\mapsto x+I$. Wir wollen sehen, dass durch $(x+I)(y+I):=xy+I$ ein
wohldefiniertes Produkt auf $R/I$ gegeben wird. Dazu sei $x=x'+i$, $y=y'+j$
mit $i,j\in I$. Wir m\"ussen sehen, dass $xy+I=x'y'+I$ gilt. Das gilt aber
wegen
\[
xy-x'y'=(x'+i)(y'+j)-x'y'=x'j+iy'+ij\in I.
\]
Damit ist $R/I$ offenbar ein Ring mit Nullelement $I$ und Einselement $1+I$,
da sich die Ringaxiome direkt von $R$ auf $R/I$ \"ubertragen. Wir erhalten
\begin{Satz}
  Sei $I$ ein Ideal des Rings $R$. Dann ist $R\to R/I$, $x\mapsto x+I$ ein
  Epimorphismus von Ringen mit Kern $I$.
\end{Satz}
Das zeigt, dass jedes Ideal der Kern eines Ringhomomorphismus ist. Fast
w\"ortlich genauso wie in der Gruppentheorie (oder linearen Algebra) beweist
man nun
\begin{Satz}[Homomorphiesatz]
  Sei $\phi:R\to S$ ein Homomorphismus von Ringen mit Kern $I$. Dann gilt
  $R/I\cong \text{Bild}(\phi)$.
\end{Satz}
Auch die Isomorphies\"atze der Gruppentheorie haben ein Analogon in der
Ringtheorie.
\begin{Satz}[Isomorphies\"atze]
  \begin{itemize}
  \item[(a)] Sei $R$ ein Ring mit Teilring $S$ und Ideal $I$. Dann gilt
  $S/S\cap I\cong S+I/I$.
\item[(b)] Seien $I$ und $J$ Ideale des Rings $R$ mit $J\subseteq I$. Dann ist
  $I/J$ ein Ideal von $R/J$, und es gilt $R/I\cong (R/J)/(I/J)$.
  \end{itemize}
\end{Satz}
\subsection{Chinesischer Restsatz}
Es seien $I$ und $J$ Ideale eines Rings $R$. Offenbar sind auch $I+J$ und
$I\cap J$ Ideale von $R$. Die multiplikative Struktur von $R$ erlaubt es uns,
aus $I$ und $J$ ein weiteres Ideal zu konstruieren. Mit $IJ$ bezeichnen wir
die Menge der endlichen Summen mit Summanden der Form $ij$ mit $i\in I$, $j\in
J$. Offenbar ist $IJ$ ein Ideal von $R$. Man beachte, dass nicht jedes
Element aus $IJ$ die Form $ij$ haben muss. Es gilt $IJ\subseteq I\cap J$.

Um die n\"achsten Begriffe zu motivieren, beweisen wir zun\"achst
\begin{Lemma}\label{L:mn}
  Es seien $m,n\in\ZZZ$, und nicht beide gleich $0$. Ferner sei $d\in\NNN$ der
  gr\"o\ss te gemeinsame Teiler von $m$ und $n$. Dann gilt $(m)+(n)=(d)$.
\end{Lemma}
\begin{proof}[Beweis]
  Da $\ZZZ$ ein Haupidealring ist, und $(m)+(n)$ nicht das Nullideal ist, gibt
  es $d\in\NNN$ mit $(m)+(n)=(d)$. Daher gibt es $a,b\in\ZZZ$ mit
  $d=am+bn$. Wegen $(m)\subseteq (d)$ ist $d$ ein Teiler von $m$, und analog
  ist $d$ ein Teiler von $n$. Aber die Darstellung $d=am+bn$ zeigt, dass jeder
  gemeinsame Teiler von $m$ und $n$ auch ein Teiler von $d$ ist, und die
  Behauptung folgt.
\end{proof}
Wir sehen also, dass insbesondere $m$ und $n$ genau dann teilerfremd sind,
wenn $(m)+(n)=\ZZZ$ gilt. Das motiviert (hoffentlich) die folgende
\begin{Definition}
  Die Ideale $I$ und $J$ des Rings $R$ hei\ss en \ei{teilerfremd}, wenn
  $R=I+J$ gilt.
\end{Definition}
Sind $R_1,R_2,\dots,R_n$ Ringe, dann ist das direkte Produkt $R_1\times
R_2\times\dots\times R_n$ mit den komponentenweisen Operationen wieder ein
Ring. Im folgenden wollen wir in gewissen Situationen komplizierte Ringe als
direkte Produkte einfacherer Ringe schreiben.

Zuvor ben\"otigen wir noch ein
\begin{Lemma}\label{L:I+JK}
  Es sei $I$ ein Ideal des Rings $R$, das teilerfremd ist zu den Idealen $J$
  und $K$. Dann gilt\begin{itemize}
\item[(a)] $IJ+JI=I\cap J$.
\item[(b)] $I$ ist teilerfremd zu $JK$ und $J\cap K$.
  \end{itemize}
\end{Lemma}
\begin{proof}[Beweis]
  Wegen $IJ,JI\subseteq I\cap J$ folgt $IJ+JI\subseteq I\cap J$. Sei nun
  umgekehrt $x\in I\cap J$. Wegen $I+J=R$ gibt es $i\in I$, $j\in J$ mit
  $1=i+j$. Es folgt
\[
x=(i+j)x=ix+jx\in IJ+JI,
\]
und damit (a). Sei nun weiter $1=i'+k$ mit $i'\in I$, $k\in K$. Wir erhalten
\[
1=(i+j)(i'+k)=ii'+ik+ji'+jk\in I+jk\subseteq I+JK,
\]
und die Behauptung folgt.
\end{proof}
Induktiv folgt aus (b)
\begin{Lemma}\label{L:I_icap}
  Das Ideal $I_1$ des Rings $R$ sei teilerfremd zu den Idealen
  $I_2,\dots,I_n$. Dann ist $I_1$ teilerfremd zu $I_2\cap\dots\cap I_n$.
\end{Lemma}
Das Hauptergebnis dieses Abschnitts ist
\begin{Satz}[Chinesischer Restsatz]
  Seien $I_1,I_2,\dots,I_n$ paarweise teilerfremde Ideale eines Rings
  $R$. Sei $I$ der Schnitt der Ideale $I_i$. Dann ist die Abbildung
\begin{align*}
R/I &\to R/I_1\times R/I_2\times\dots\times R/I_n\\
r+I &\mapsto (r+I_1,r+I_2,\dots,r+I_n)
\end{align*}
ein Isomorphismus von Ringen.
\end{Satz}
\begin{proof}[Beweis]
Offenbar ist $I$ der Kern der Abbildung $R\to R/I_1\times
R/I_2\times\dots\times R/I_n$, $r\mapsto (r+I_1,r+I_2,\dots,r+I_n)$, die
Behauptung folgt also aus dem Homomorphiesatz, sobald wir die Surjektivit\"at
der Abbildung nachgewiesen haben. F\"ur $n=1$ ist nichts zu zeigen.

Sei nun $n=2$. Wegen $R=I_1+I_2$ gibt es $i_1\in I_1$, $i_2\in I_2$ mit
$1=i_1+i_2$. Seien $x,y\in R$ beliebig, und setze $r=xi_2+yi_1$. Es gilt
$r=x(1-i_1)+yi_1=x+(y-x)i_1\in x+I_1$ und $r=xi_2+y(1-i_2)\in y+I_2$. Das
Element $r$ wird also auf $(x+I_1,y+I_2)$ abgebildet. Da $x$ und $y$ beliebig
waren, folgt die Behauptung f\"ur $n=2$.

F\"ur $n\ge3$ benutzen wir vollst\"andige Induktion, wobei wir f\"ur den
Induktionsschritt den Fall $n=2$ benutzen: Nach Lemma \ref{L:I_icap} ist $I_1$
teilerfremd zum Schnitt $I_1'$ der $I_i$ mit $2\le i\le n$. Beachte, dass
$I=I_1\cap I_1'$ gilt. Wegen des Falls $n=2$ ist $R/I_1\to R/I_1\times R/I_1'$
ein Isomorphismus, und $R/I_1'\to R/I_2\times\dots\times R/I_n$ ist ein
Isomorphismus nach der Induktionsannahme f\"ur $n-1$. Die Behauptung folgt.
\end{proof}
Die h\"aufigste Anwendung des chinesischen Restsatzes ist f\"ur kommutative
Ringe $R$. In diesem Fall l\"asst sich der Schnitt paarweiser teilerfremder
Ideale noch anders ausdr\"ucken:
\begin{Lemma}
  Seien $I_1,I_2,\dots,I_n$ paarweise teilerfremde Ideale eines kommutativen
  Rings. Dann gilt $I_1I_2\dots I_n=I_1\cap I_2\cap\dots\cap I_n$.
\end{Lemma}
\begin{proof}[Beweis]
  F\"ur $n=2$ folgt das aus Lemma \ref{L:I+JK}(a), und allgemein per Induktion.
\end{proof}
Eine wichtiger Spezialfall des chinesischen Restsatz ist
\begin{Satz}
  Sei $n=\prod p_i^{e_i}$ die nat\"urliche Primfaktorzerlegung von
  $n\in\NNN$. Dann ist der Ring $\ZZZ/n\ZZZ$ isomorph zum direkten Produkt der
  Ringe $\ZZZ/p_i^{e_i}\ZZZ$.
\end{Satz}
\subsection{Anwendungen der Kongruenzrechnung}
Ein klassisches und wichtiges Gebiet der Mathematik ist die Frage der
ganzzahligen L\"osbarkeit z.B.\ polynomialer Gleichungen oder
Gleichungssysteme. Hier sind Restklassenbetrachtungen h\"aufig von gro\ss em
Nutzen. Wir betrachten exemplarisch einige Beispiele:
\begin{itemize}
\item $4x-1=y^2+z^2$ hat keine ganzzahligen L\"osungen. Dazu sei
  $\ZZZ\to\ZZZ/4\ZZZ$ der nat\"urliche Homomorphismus, und $\bar{a}$ sei das
  Bild von $a\in\ZZZ$. Ist $x,y,z$ eine L\"osung der gegebenen Gleichung, dann
  gilt auch $\bar{3}=-\bar{1}=\overline{4x-1}=\bar{y}^2+\bar{z}^2$. F\"ur
  $a\in\ZZZ$ gilt aber $\bar{a}^2=\bar{0}$ oder $\bar{1}$, also $\bar{3}=u+v$
  mit $u,v\in\{\bar{0},\bar{1}\}$, ein Widerspruch.
\item Die nat\"urliche Zahl $n$ enthalte in der Dezimaldarstellung nur die
  Ziffern $6$ und $0$. Kann $n$ eine Quadratzahl sein? Wir nehmen an, $n$ ist
  eine Quadratzahl. Dann endet $n$ auf eine gerade Anzahl von $0$ern. Streicht
  man diese, dann erh\"alt man eine Quadratzahl, die mit $6$ endet, also
  modulo $4$ den Rest $2$ hat, was aber nicht geht, da die Reste der Quadrate
  modulo $4$ nur $0$ und $1$ sind.
\item $15x^2-7y^2=9$ ist in $\ZZZ$ nicht l\"osbar: Sei $x,y,z$ eine L\"osung.
  Man sieht, dass $7y^2$ durch $3$ teilbar ist. Dann ist $y$ durch $3$
  teilbar, also $7y^2$ durch $9$ teilbar. Somit ist $15x^2$ durch $9$ teilbar,
  also $x$ durch $3$ teilbar. Wir setzen also $x=3a$, $y=3b$ mit $a,b\in\ZZZ$.
  Es folgt $15a^2-7b^2=1$. Eine Betrachtung modulo $4$ bringt hier keinen
  Widerspruch, sie liefert $\bar{b}^2-\bar{a}^2=bar 1$ in $\ZZZ/4\ZZZ$, aber
  das ist nat\"urlich l\"osbar. In $\ZZZ/3\ZZZ$ hingegen erhalten wir
  $-\bar{b}^2=\bar{1}$, aber $\bar{b}^2\in\{\bar{0},\bar{1}\}$, ein
  Widerspruch.
\item F\"ur welche nat\"urlichen Zahlen $n$ ist $a=2^n+65$ eine Quadratzahl?
  Ist $n=2m+1$ ungerade, dann ist $2^n=2\cdot 4^m$ stets kongruent $\pm2$
  modulo $5$, ferner ist $a$ ungerade, also endet $a$ auf $3$ oder $7$. Aber
  ungerade Quadratzahlen enden auf $1$, $5$ oder $9$, ein Widerspruch. Daher
  ist $n=2m$ gerade, und $2^n$ eine Quadratzahl. Wegen $a>2^n=(2^m)^2$ ist $a$
  mindestens so gro\ss\ wie die n\"achste Quadratzahl nach $2^n$, also
  $2^{2m}+65=2^n+65=a\ge (2^m+1)^2$, und somit $65\ge 2^{m+1}+1$, also $2^m\le
  32$ und daher $m\le5$. Durchprobieren liefert die einzigen L\"osungen
  $2^4+65=9^2$ und $2^{10}+65=33^2$.
\item Die Kongruenzmethode f\"uhrt nicht immer zum Ziel: Man kann zeigen, dass
  $3x^3+4y^3+5z^3=0$ modulo jedem $n\in\ZZZ$ nicht trivial l\"osbar ist, und
  dennoch hat diese Gleichung keine ganzzahlige L\"osung au\ss er $x=y=z=0$.
  Beide Aussagen sind aber nicht einfach zu beweisen, und liegen jenseits
  dieser Vorlesung.
\end{itemize}
\subsection{Maximale Ideale}
Sei $R$ ein kommutativer Ring. Ein \ei{maximales Ideal} ist ein Ideal $I$ von
$R$ mit $I\ne R$, so dass kein Ideal $J$ von $R$ existiert mit $I\subsetneq
J\subsetneq R$. Die wichtige Aussage in diesem Zusammenhang ist.
\begin{Satz}
  Sei $I$ ein Ideal eines kommutativen Ringes $R$. Dann ist $I$ genau dann
  maximal, wenn $R/I$ ein K\"orper ist.
\end{Satz}
\begin{proof}[Beweis]
  Sei $I$ maximal, und $a+I\ne 0$ in $R/I$. Damit ist $a\ne I$. Da $Ra+I$ ein
  Ideal ist, welches $I$ echt enth\"alt, muss $Ra+I=R$ gelten. Daher gibt es
  $r\in R$ mit $1\in ra+I$, also $1+I=(a+I)(r+I)$, und $a+I$ ist daher
  multiplikativ invertierbar in $R/I$. Somit ist $R$ ein K\"orper.

Sei nun $I$ nicht maximal, und $J$ ein Ideal echt zwischen $I$ und $R$. F\"ur
$a\in J$ und $r\in R$ gilt dann $(a+I)(r+I)=ar+I\subseteq J$, also $1\ne
ar+I$, und somit ist $a+I$ nicht invertierbar.
\end{proof}
\paragraph{Beispiele maximaler Ideale}
\begin{itemize}
\item Die Ideale $I$ von $\ZZZ$ haben die Form $n\ZZZ$ mit $n=0,1,\dots$. Ein
  solches Ideal ist genau dann maximal, wenn $n$ eine Primzahl ist.
\item Sei $R$ die Menge der reellwertigen Abbildungen $M\to\RRR$ f\"ur eine
  Menge $M$. Offenbar ist $R$ ein Ring. F\"ur $m\in M$ ist die Abbildung
  $R\to\RRR$, $f\mapsto f(m)$ ein Epimorphismus von Ringen. Der Kern $I$ ist
  das Ideal der reellwertigen Abbildungen $f$ mit $f(m)=0$. Wegen
  $R/I\cong\RRR$ ist $R/I$ ein K\"orper, und $I$ damit ein maximales Ideal.
\end{itemize}
Direkt aus dem Zornschen Lemma folgt
\begin{Satz}[Krull]
  Sei $I\subsetneq R$ ein Ideal eines kommutativen Rings $R$. Dann ist $I$ in
  einem maximalen Ideal von $R$ enthalten.
\end{Satz}
\subsection{Primideale}
Sei $R$ wieder ein kommutativer Ring. Ein Ideal $I$ von $R$ hei\ss t
\ei{Primideal}, wenn $R/I$ ein Integrit\"atsring ist. Ist also $ab\in I$, dann
ist $(a+I)(b+I)\subseteq ab+I=0$ in $R/I$, und wegen der Nullteilerfreiheit
gilt $a\in I$ oder $b\in I$. Da K\"orper Integrit\"atsringe sind, sind
maximale Ideale automatisch Primideale. Das einizige nichtmaximale Primideal
von $\ZZZ$ ist $\{0\}$.

Induktiv sieht man sofort, dass nilpotente Elemente in allen Primidealen
liegen. Wir wollen die Umkehrung beweisen.
\begin{Lemma}
  Sei $R$ ein kommutativer Ring (mit $0\ne1$) und $S$ eine multiplikativ
  abgeschlossene Teilmenge von $S$ mit $0\ne S$. Dann hat $R$ ein Primideal
  mit $S\cap I=\emptyset$.
\end{Lemma}
\begin{proof}[Beweis]
  Sei $M$ die Menge der Ideale von $R$, welche disjunkt sind zu $S$. Wegen
  $\{0\}\in M$ ist $M$ nicht leer. Die Vereinigung von Ketten aus $M$ liegt
  wieder in $M$, nach dem Zornschen Lemma hat $M$ daher ein maximales Element
  $I$. Wir wollen sehen, dass $I$ ein Primideal ist. Dazu seien $x,y\in
  R\setminus I$. Wir m\"ussen sehen, dass auch $xy\in R\setminus I$ gilt.
  
  Da die Ideale $Rx+I$ und $Ry+I$ echt gr\"o\ss er als $I$ sind, enthalten sie
  Elemente $s_x$ und $s_y$ aus $S$. Es bestehen also Gleichungen
  $s_x=r_xx+i_x$, $s_y=r_yy+i_y$ mit $r_x,r_y\in R$, $i_x,i_y\in I$. Es gilt
  $s_xs_y+I=(s_x+I)(s_y+I)=(r_xx+I)(r_yy+I)=(r_xr_yxy+I)$. Da $S$
  multiplikativ abgeschlossen ist gibt es $i\in I$ mit $r_xr_yxy+i\in S$.
  W\"are $xy\in I$, dann g\"alte auch $r_xr_yxy+i\in I$, im Widerspruch zu
  $I\cap S=\emptyset$. Die Behauptung folgt.
\end{proof}
\begin{Satz}
  Sei $R$ ein kommutativer Ring (mit $0\ne1$). Die Menge der nilpotenten
  Elemente besteht aus dem Durchschnitt der Primideale.
\end{Satz}
\begin{proof}[Beweis]
  Sei $x\in R$ nicht nilpotent. Wir m\"ussen sehen, dass es ein Primideal $I$
  gibt, das $x$ nicht enth\"alt. Da $x$ nicht nilpotent ist, ist
  $S=\{x^n|n\in\NNN\}$ eine multiplikativ abgeschlossene Menge mit $0\notin
  S$. Die Behauptung folgt aus dem vorigen Lemma.
\end{proof}
\subsection{Polynome}
In diesem Abschnitt sei $R$ stets ein kommutativer Ring. Sei $X$ ein Symbol.
Ein \ei{Polynom} in der \ei{Variablen} $X$ ist eine formale Summe
$r_0+r_1X+r_2X^2+\dots+r_nX^n$ f\"ur ein $n\in\NNN_0$. Hierbei trifft man die
Festsetzung $X^0=1$. Die Menge der Polynome bildet einen Ring unter
koeffizientenweiser Addition, und den Festsetzungen $rX=Xr$ und
$X^mX^n=X^{m+n}$. Man schreibt $R[X]$ f\"ur diesen Ring. Man fasst $R$ als
Teilring von $R[X]$ auf, indem man $r\in R$ mit dem Polynom $r=rX^0$
identifiziert. Die Elemente $r_0,r_1,\dots$ hei\ss en die \ei{Koeffizienten}
des Polynoms $f:=\sum r_iX^i$. Sei $f\ne0$, und $n\in\NNN_0$ maximal mit
$r_n\ne 0$. Dann hei\ss t $n$ der \ei{Grad} von $f$, und $r_n$ der
\ei{Leitkoeffizient}. Man schreibt $n=\grad f$. Man nennt $f$ \ei{normiert},
wenn $r_n=1$ gilt. Der Koeffizient $r_0$ wird \ei{konstanter Term} oder
\ei{Absolutglied} genannt. F\"ur $f=0$ setzt man $\grad f=-\infty$.

Aus den Definitionen folgt unmittelbar

\begin{Lemma}
  Seien $f,g\in R[X]$ Polynome. Dann gilt $\grad(f+g)\le\max(\grad f,\grad g)$
  und $\grad(f\cdot g)\le\grad f+\grad g$. Ist $R$ ein Integrit\"atsbereich,
  dann gilt $\grad(f\cdot g)=\grad f+\grad g$
\end{Lemma}

An der letzten Aussage sieht man, dass wenn $R$ ein Integrit\"atsbereich ist,
dann gilt das auch f\"ur $R[X]$.

Wie bei ganzen Zahlen hat man auch f\"ur Polynome \"uber K\"orpern eine
Division mit Rest. Genauer gilt:

\begin{Satz}
  Sei $g\in R[X]$ ein Polynom mit invertierbarem Leitkoeffizienten, und $f\in
  R[X]$ beliebig. Dann gibt es eindeutig gegebene Polynome $q,r\in R[X]$ mit
  $f=q\cdot g+r$ und $\grad r<\grad g$.
\end{Satz}
\begin{proof}[Beweis]
Wir wollen zun\"chst die Eindeutigkeit beweisen. Dazu sei $f=q'\cdot g+r'$
eine weitere Darstellung der gegebenen Form. Dann folgt $(q-q')g=r'-r$. Da der
Leitkoeffizient von $g$ kein Nullteiler ist, gilt
$\grad(r'-r)=\grad(q-q')+\grad g$. Wegen $\grad(r'-r)<\grad g$ folgt $q=q'$
und $r'=r$.

Es bleibt die Existenz zu zeigen. Sei $a$ die Inverse des Leitkoeffizienten
$g$, und $g'=ag$. Ist $f=q'g'+r$ eine Division durch $g'$ mit Rest der
verlangten Art, dann gilt das auch f\"ur $f=qg+r$ mit $q=aq'$. Wir d\"urfen
also annehmen, dass $g$ normiert ist.

Ist $\grad f<\grad g$, dann gibt es nichts zu zeigen, da wir $q=0$ und $r=f$
setzen k\"onnen. Wir verwenden nun vollst\"andige Induktion \"uber $\grad f$.
Sei $n=\grad f\ge\grad g$, und $a$ der Leitkoeffizient von $f$. Setze
$f'=f-aX^{\grad f-\grad g}g$. Dann gilt $\grad f'<\grad f$. Nach
Induktionsvoraussetzung gibt es also eine Darstellung $f'=q'g+r$ mit $\grad
r<\grad g$. Die Behauptung folgt nun aus $f=f'+aX^{\grad f-\grad
  g}g=q'g+r'+aX^{\grad f-\grad g}g=(q'+aX^{\grad f-\grad g})g+r$.
\end{proof}

Eine wichtige Folgerung ist
\begin{Satz}
  Sei $R$ ein K\"orper, und $I\ne\{0\}$ ein Ideal von $R[X]$. Sei $g\ne0$ ein
  Polynom kleinsten Grades aus $I$. Dann ist $g$ bis auf einen Faktor $0\ne
  r\in R$ eindeutig, und es gilt $I=(g)$. Insbesondere ist $R[X]$ ein
  Hauptidealring.
\end{Satz}
\begin{proof}[Beweis]
  Ist $0\ne r\in R$ und $g\in R[X]$, dann gilt $g\in I$ genau dann wenn $rg\in
  I$. Die Differenz zweier normierter Polynome von gleichem Grad $n$ hat Grad
  $<n$, daraus folgt die Eindeutigkeitsaussage. Sei nun $0\ne g\in I$ von
  kleinstem Grad, und $f\in I$ beliebig. Sei $f=qg+r$ eine Division mit Rest.
  Aus $r\in I$ und $\grad r<\grad g$ folgt $r=0$, also $I\subseteq (g)$. Die
  umgekehrte Inklusion ist sowieso klar.
\end{proof}
\begin{Bemerkung*}
  Obige Aussage wird im allgemeinen falsch, wenn $R$ kein K\"orper mehr
  ist. Sogar ein Integrit\"atsbereich zu sein ist zu schwach, wie das von $2$
  und $X$ in $\ZZZ[X]$ erzeugte Ideal zeigt.
\end{Bemerkung*}
Ist $a\in R$ und $f=\sum r_iX^i\in R[X]$ ein Polynom, dann bedeutet $f(a)$ die
Summe $\sum r_ia^i$, mit der Konvention $a^0=1$ auch dann, wenn $a=0$
gilt. Offenbar ist $R[X]\to R$, $f\mapsto f(a)$ ein Ringhomomorphismus. Man
nennt $a$ eine \ei{Nullstelle} von $f$, wenn $f(a)=0$ gilt.

\begin{Satz}
  Ist $a$ eine Nullstelle des Polynoms $f$, dann gibt es $g\in R[X]$ mit
  $f=(X-a)g$.
\end{Satz}
\begin{proof}[Beweis]
  Schreibe $f=(X-a)g+r$ mit $\deg r<\deg(X-a)=1$. Daher gilt $r\in
  R$. Einsetzen von $a$ f\"ur $X$ liefert $0=f(0)=(a-a)g(a)+r=r$, und die
  Behautptung folgt.
\end{proof}
Eine wichtige Folgerung ist
\begin{Satz}\label{S:Nullstellen}
  Sei $R$ ein Integrit\"atsring, und $0\ne f\in R[X]$. Dann hat $f$
  h\"ochstens $\grad f$ Nullstellen.
\end{Satz}
\begin{proof}[Beweis]
  F\"ur $\grad f\le 1$ ist die Aussage klar. Wir beweisen sie allgemein durch
  vollst\"andige Induktion. Seien $a_1,\dots,a_r$ Nullstellen von
  $f$. Wir schreiben $f=(X-a_1)g$. Wegen $0=f(a_i)=(a_i-a_1)g(a_i)$ und der
  Nullteilerfreiheit von $R$ ist $a_i$ eine Nullstelle von $g$ f\"ur $i\ge
  2$. Daher gilt $r-1\le\grad g=\grad f-1$, und die Behauptung folgt.
\end{proof}
\begin{Bemerkung*}
  Man \"uberlege sich ein Gegenbeispiel zum Satz, wenn $R$ ein geeigneter Ring
  mit Nullteilern ist.
\end{Bemerkung*}
Eine \"uberraschende Folgerung obigen Satzes ist
\begin{Satz}
  Sei $G$ eine endliche Untergruppe der multiplikativen Gruppe eines
  K\"orpers. Dann ist $G$ zyklisch.
\end{Satz}
\begin{proof}[Beweis]
  Sei $n$ die Ordnung von $G$. Ist die abelsche Gruppe $G$ nicht zyklisch,
  dann gibt es nach dem Struktursatz f\"ur abelsche Gruppen eine Primzahl $p$
  und eine zu $C_p\times C_p$ isomorphe Untergruppe $H$ von $G$. Aber f\"ur
  alle Elemente $h\in H$ gilt $h^p=1$. Daher sind die $p^2$ Elemente aus $H$
  Nullstellen des Polynoms $X^p-1$, im Widerspruch zu vorigem Satz.
\end{proof}
\begin{Korollar}\label{K:ZpZzyklisch}
  Sei $p$ eine Primzahl. Dann ist die multiplikative Gruppe des K\"orpers
  $\ZZZ/p\ZZZ$ zyklisch.
\end{Korollar}
\begin{Bemerkung*}
  Das Korollar ist eine reine Existenzaussage in dem Sinne, dass man wei\ss\,
  dass $(\ZZZ/p\ZZZ)^\times$ zyklisch ist, aber keinen Erzeuger explizit
  geliefert bekommt. In der Tat ist keine explizite Formel oder guter
  Algorithmus zum Auffinden solcher Erzeuger bekannt.
\end{Bemerkung*}
Ist $R$ ein kommutativer Ring, und $a$ eine Nullstelle eines Polynoms $0\ne
f\in R[X]$. Sei $e\in\NNN$ mit $f=(X-a)^eg$ f\"ur ein $e\in\NNN$. Wegen
$e+\grad g=\grad f$ gilt $e\le\grad f$. Insbesondere gibt es ein maximales $e$
mit obiger Darstellung. Man nennt $e$ die \ei{Vielfachheit} der Nullstelle
$a$. Man \"uberlegt sich sofort, dass Satz \ref{S:Nullstellen} auch dann noch
gilt, wenn man die Nullstellen mit Vielfachheit z\"ahlt. Ist die Vielfachheit
$1$, dann nennt man die Nullstelle \ei{einfach}. Die Ableitung ist ein
einfaches Hilfsmittel, um Polynome auf vielfache Nullstellen zu testen.
\begin{Definition}
  Auf dem Polynomring $R[X]$ definiert man durch $r':=0$, $(rX^i)':=riX^{i-1}$
  ($r\in R$, $i\ge1$) eine additive Abbildung $R[X]\to R[X]$, $f\mapsto
  f'$. Man nennt diese Abbildung eine \ei{Ableitung} oder
  \ei{Differentiation}.
\end{Definition}
Man rechnet sofort nach, dass die aus der Analysis gewohnte Produktregel
$(fg)'=f'g+fg'$ gilt. Allerdings ist Vorsicht angebracht: Gilt $f'=0$, dann
muss $f$ nicht ein konstantes Polynom sein. Ist beispielsweise $R=\ZZZ/p\ZZZ$,
dann gilt $(X^p)'=pX^{p-1}=0X^{p-1}=0$.
\begin{Satz}
  Sei $R$ ein Integrit\"atsbereich, und $a$ eine Nullstelle des Polynoms $f\in
  R[X]$. Dann ist $a$ eine einfache Nullstelle genau dann, wenn $f'(a)\ne0$
  gilt.
\end{Satz}
\begin{proof}[Beweis]
  Sei $f=(X-a)g$. Ableiten liefert $f'=(X-a)g'+g$, also $f'(a)=g(a)$. Die
  Behauptung folgt.
\end{proof}
\subsection{Einheitengruppe von $\ZZZ/n\ZZZ$}
Als Anwendung des Korollars \ref{K:ZpZzyklisch} wollen wir allgemein die
Einheitengruppe von $\ZZZ/n\ZZZ$ bestimmen. Sei $n=\prod p_e^{e_i}$ die
Primfaktorzerlegung von $n\in\NNN$. Nach dem Chinesischen Restsatz ist
$\ZZZ/n\ZZZ$ isomorph zum direkten Produkt der Ringe $\ZZZ/p_i^{e_i}\ZZZ$, und
damit ist die Einheitengruppe $(\ZZZ/n\ZZZ)^\times$ isomorph zum direkten
Produkt der Einheitengruppen $(\ZZZ/p_i^{e_i}\ZZZ)^\times$. Der folgende Satz
kl\"art die Struktur dieser Gruppen.
\begin{Satz}\label{S:ZnZ}
  \begin{itemize}
  \item[(a)] Sei $p$ eine ungerade Primzahl und $l\in\NNN$. Dann ist
    $(\ZZZ/p^l\ZZZ)^\times$ zyklisch. Sei $w\in\ZZZ$, so dass $w+p\ZZZ$ ein
    Erzeuger von $(\ZZZ/p\ZZZ)^\times$ ist. Dann ist
    $w^{p^{l-1}}(1+p)+p^l\ZZZ$ ein Erzeuger von $(\ZZZ/p^l\ZZZ)^\times$.
  \item[(b)] F\"ur $l\ge2$ ist $(\ZZZ/2^l\ZZZ)^\times$ isomorph zu einem
    direkten Produkt zweier zyklischer Gruppen der Ordnungen $2$ und
    $2^{l-2}$, mit Erzeugern $-1$ und $5$.
  \end{itemize}
\end{Satz}
Dem Beweis schicken wir einige Lemmata voraus.
\begin{Lemma}\label{L:pk}
  Sei $p$ eine Primzahl, und $1\le k\le p-1$. Dann ist $\binom{p}{k}$ durch
  $p$ teilbar.
\end{Lemma}
\begin{proof}[Beweis]
  Wegen $\binom{p}{k}=\frac{p(p-1)\dots(p-k+1)}{1\cdot2\dots k}$ ist der
  Z\"ahler durch $p$ teilbar, aber nicht der Nenner.
\end{proof}
\begin{Lemma}\label{L:ap-bp}
  Sei $p$ eine Primzahl, $l\in\NNN$, und f\"ur $a,b\in\ZZZ$ sei $a-b$ durch
  $p^l$ teilbar. Dann ist $a^p-b^p$ durch $p^{l+1}$ teilbar.
\end{Lemma}
\begin{proof}[Beweis]
  Schreibe $a=b+cp^l$. Dann gilt
  \begin{align*}
    a^p-b^p &= (b+cp^l)^p-b^p\\
&= \sum_{k=1}^p\binom{p}{k}b^{p-k}(cp^l)^k\\
&= b^{p-1}cp^{l+1}+\sum_{k=2}^p\binom{p}{k}b^{p-k}(cp^l)^k.
  \end{align*}
Aber f\"ur $k\ge2$ gilt $kl\ge 2l\ge l+1$, und die Behauptung folgt.
\end{proof}
\begin{Lemma}\label{L:pl}
  Sei $l\ge2$, und $p\ne2$ eine Primzahl. Dann ist
  $(1+p)^{p^{l-2}}-(1+p^{l-1})$ durch $p^l$ teilbar.
\end{Lemma}
\begin{proof}[Beweis]
Wir benutzen vollst\"andige Induktion \"uber $l$. Die Aussage ist klar f\"ur
$l=2$.

Sie gelte nun f\"ur $l$. Nach obigem Lemma ist $(1+p)^{p^{l-1}}-(1+p^{l-1})^p$
durch $p^{l+1}$ teilbar. Wegen
\[
(1+p^{l-1})^p=(1+p^l)+\sum_{k=2}^p\binom{p}{k}(p^{l-1})^k
\]
ist zu zeigen, dass die Summe durch $p^l$ teilbar ist. Wegen Lemma \ref{L:pk}
ist das klar f\"ur jeden Summanden au\ss er dem f\"ur $k=p$. Hierf\"ur
ben\"otigen wir $(l-1)p\ge l+1$. Das gilt aber wegen $p\ge3$ und $l\ge2$.
\end{proof}
\begin{Bemerkung*}
  Man sieht, dass der letzte Schritt im Beweis f\"ur $p=2$ nicht
  funktioniert. In der Tat ist die entsprechende Aussage falsch. Dies f\"uhrt
  dazu, da\ss\ die Primzahl $2$ hier und in der Zahlen- und Gruppentheorie oft
  eine (unangenehme) Sonderrolle spielt.
\end{Bemerkung*}
\begin{Lemma}\label{L:2l}
  F\"ur $l\ge3$ ist $5^{2^{l-3}}-(1+2^{l-1})$ durch $2^l$ teilbar.
\end{Lemma}

\begin{proof}[Beweis]
  Die Behauptung ist klar f\"ur $l=3$. Sie gelte f\"ur $l$. Nach Lemma
  \ref{L:ap-bp} ist dann $5^{2^{l-2}}-(1+2^{l-1})^2$ durch $2^{l+1}$ teilbar.
  Wegen $(1+2^{l-1})^2=1+2^l+2^{2l-2}$ und $2l-2\ge l+1$ folgt die Behauptung.
\end{proof}
\begin{proof}[Beweis von Satz \ref{S:ZnZ}]
  Die Gruppe $(\ZZZ/p^l\ZZZ)^\times$ besteht aus den Restklassen
  $i+p^l\ZZZ$, mit $0\le i\le p^l-1$ und $i$ teilerfremd zu $p$. Daher hat
  $(\ZZZ/p^l\ZZZ)^\times$ die Ordnung $p^{l-1}(p-1)$.
  
  (a) Wir zeigen zun\"achst, dass $1+p$ die Ordnung $p^{l-1}$ hat. W\"are das
  nicht der Fall, dann w\"are die Ordnung von $1+p$ ein echter Teiler von
  $p^{l-1}$. Insbesondere w\"are $(1+p)^{p^{l-2}}-1$ durch $p^l$ teilbar. Nach
  Lemma \ref{L:pl} w\"are dann $p^{l-1}$ durch $p^l$ teilbar, was nat\"urlich
  Unsinn ist.
  
  Da $w$ modulo $p$ die Ordnung $p-1$ hat, ist die Ordnung von $w$ modulo
  $p^l$ ein Vielfaches von $p-1$. Diese Ordnung teilt $p^{l-1}(p-1)$, daher
  hat $w^{p^{l-1}}$ die Ordnung $p-1$, und nach obigem hat $w^{p^{l-1}}(1+p)$
  modulo $p^l$ die Ordnung $(p-1)p^{l-1}$, die Behauptung folgt.
  
  (b) Die Ordnung von $(\ZZZ/2^l\ZZZ)^\times$ ist $2^{l-1}$. Lemma \ref{L:2l}
  zeigt, dass die Ordnung von $5+2^l\ZZZ$ gerade $2^{l-2}$ ist. Zum Beweis der
  Aussage ist also zu zeigen, dass $-1+2^l\ZZZ$ nicht in der von $5+2^l\ZZZ$
  erzeugten zyklischen Gruppe liegt. Wir nehmen das Gegenteil an. Dann gibt es
  $m\in\NNN$, so dass $2^l$ ein Teiler von $5^m-(-1)=5^m+1$ ist. Insbesondere
  ist $4$ ein Teiler von $5^m+1$, ein Widerspruch.
\end{proof}

\subsection{Quotientenk\"orper}
Es sei $R$ ein Integrit\"atsbereich. Analog wie man den K\"orper $\QQQ$ aus
$\ZZZ$ konstruiert, werden wir zu $R$ einen kleinsten K\"orper $K$
konstruieren, in dem $R$ enthalten ist. Auf der Menge der Paare $(r,s)$ mit
$r,s\in R$, $s\ne0$ f\"uhren wir eine \"Aquivalenzrelation ein. Dabei sind
zwei solche Paare $(r,s)$ und $(r',s')$ genau dann \"aquivalent, wenn
$rs'=r's$ gilt. Man rechnet sofort nach, dass man tats\"achlich eine
\"Aquivalenzrelation erh\"alt. Die \"Aquivalenzklasse von $(r,s)$ bezeichnet
man mit $\frac{r}{s}$. Sei $K$ die Menge der \"Aquivalenzklassen. Man rechnet
ebenso schnell nach, dass man durch
\[
\frac{r}{s}+\frac{r'}{s'}=\frac{rs'+r's}{ss'}
\]
und
\[
\frac{r}{s}\cdot\frac{r'}{s'}=\frac{rr'}{ss'}
\]
eine wohldefinierte Addition und Multiplikation auf $K$ bekommt. Ferner
verifiziert man, dass $K$ ein Ring ist, der via $r\mapsto\frac{r}{1}$ den Ring
$R$ als Teilring enth\"alt.

Diw wichtigste Eigenschaft von $K$ ist, dass $K$ ein K\"orper ist: Sei
$\frac{r}{s}\ne0$. Dann gilt $r\ne0$, also $\frac{s}{r}\in K$. Wegen
$\frac{r}{s}\frac{s}{r}=1$ ist $\frac{r}{s}$ multiplikativ invertierbar. Man
nennt daher $K$ den \ei{Quotientenk\"orper} von $R$.
\subsection{Teilbarkeit}
Im folgenden sei $R$ wieder ein Integrit\"atsbereich. Wir wollen einige von
$\ZZZ$ bekannte Eigenschaften \"uber Teilbarkeit auf $R$ \"ubertragen, aber
auch sehen, dass man mit Verallgemeinerungen vorsichtig sein muss.

Sind $r,s\in R$, dann sagt man, $r$ \ei{teilt} $s$ (oder $r$ ist ein
\ei{Teiler} von $s$), wenn es $x\in R$ gibt mit $s=xr$. Insbesondere ist $1$
ein Teiler von allen $r\in R$, und $0$ wird von allen $r\in R$ geteilt. Ist
$K\supseteq R$ der Quotientenk\"orper von $R$, und $s\ne0$, dann ist $s$ ein
Teiler von $r$ genau dann, wenn $\frac{r}{s}\in R$ gilt.

Ist $s$ ein Teiler von $r$, dann schreibt man $r\mid s$, und wenn das nicht
der Fall ist, so schreibt man $r \nmid s$.

Ist $0\ne s\in R$, so besitzt $s$ gewisse triviale Teiler. So ist z.B.\ jede
Einheit $u$ ein Teiler von $s$, denn ist $uu'=1$, so gilt $s=u(u's)$. Ferner
ist wegen $s=(su)u'$ auch $su$ ein Teiler von $s$. Ist $s$ keine Einheit, und
hat $s$ au\ss er diesen trivialen Teilern keine weiteren Teiler, dann hei\ss t
$s$ \ei{unzerlegbar} oder \ei{irreduzibel}. Ist $r$ nicht irreduzibel, dann
ist $r$ \ei{reduzibel}.

Die irreduziblen Elemente aus
$\ZZZ$ sind die Zahlen $\pm p$, wobei $p$ eine Primzahl ist. Man hat sich
daran gew\"ohnt, dass man in $\ZZZ$ eine (bis auf Vorzeichen) eindeutige
Primfaktorzerlegung hat. Erst Gau\ss\ hat erkannt, dass diese Aussage nicht
trivial ist und einen Beweis ben\"otigt, den er auch gegeben hat.
\paragraph{Beispiel f\"ur Versagen der eindeutigen Zerlegung in irreduzible
  Faktoren}
Sei $R=\ZZZ[\sqrt{-5}]=\{a+b\sqrt{-5}|a,b\in\ZZZ\}$. Man rechnet sofort nach,
dass $R$ ein Ring ist. Wir betrachten die Abbildung $N:R\to\NNN_0$,
$a+b\sqrt{-5}\mapsto a^2+5b^2$. Dabei ist also $N(x)$ das Quadrat des Betrags
der komplexen Zahl $x$. Man wei\ss, oder rechnet sofort nach, dass
$N(xy)=N(x)N(y)$ gilt f\"ur alle $x,y\in R$. Ist $r$ eine Einheit, dann gibt
es $r'$ mit $rr'=1$. Hieraus folgt $1=N(1)=N(rr')=N(r)N(r')$, also
$N(r)=1$. Ist umgekehrt $1=N(a+b\sqrt{-5})=a^2+5b^2$, dann gilt $b=0$ und
$a=\pm1$. Somit ist $r$ genau dann eine Einheit, wenn $N(r)=1$ gilt.

Ist $0\ne r$ keine Einheit und reduzibel, dann gibt es $x,y\in R$ mit $r=xy$,
so dass $x$ und $y$ keine Einheiten sind. Es folgt $N(r)=N(x)N(y)$. Da es in
$R$ keine Nichteinheiten $s\ne0$ mit $N(s)=2$ oder $N(s)=3$ gibt, folgt
$N(r)\ge16$. Wegen $N(2)=4$, $N(3)=9$, und $N(1\pm\sqrt{-5})=6$ sind die
Elemente $2$, $3$ und $1\pm\sqrt{-5}$ daher irreduzibel, und $2$ unterscheidet
sich weder von $1+\sqrt{-5}$, noch von $1-\sqrt{-5}$ um eine Einheit. Daher
sind $6=2\cdot3=(1+\sqrt{-5})(1-\sqrt{-5})$ zwei grunds\"atzlich verschiedene
Zerlegungen in irreduzible Faktoren.

Ein mit der Irreduzibilit\"at verwandter Begriff ist der eines primen
Elements. Dabei hei\ss t eine Nichteinheit $0\ne r\in R$ \ei{prim}, wenn aus
$r\mid ab$ stets $r\mid a$ oder $r\mid b$ folgt. Die letzte Bedingung ist
\"aquivalent dazu, dass das Hauptideal $(p)$ ein Primideal von $R$ ist.

Eine erste einfache Beobachtung ist
\begin{Lemma}
  Jedes Primelement ist irreduzibel.
\end{Lemma}
\begin{proof}[Beweis]
  Sei $r$ prim, aber nicht irreduzibel. Dann gibt es Nichteinheiten $x,y\in R$
  mit $r=xy$. Da $r$ prim ist, teilt $r$ einen der Faktoren $x$ und $y$. Wir
  nehmen $r\mid x$ an. Es gibt also $x'\in R$ mit $x=rx'$. Es folgt $r=rx'y$,
  also $r(1-x'y)=0$. Wegen der Nullteilerfreiheit gilt $x'y=1$, daher ist $y$
  eine Einheit, ein Widerspruch.
\end{proof}
  Obiges Beispiel zeigt, dass die Umkehrung im allgemeinen nicht gilt. Die
  Elemente $2$, $3$ und $1\pm\sqrt{-5}$ sind alle irreduzibel in
  $\ZZZ[\sqrt{-5}]$, aber keines davon ist prim. Unter Zusatzvoraussetzungen
  allerdings stimmt die Umkehrung:
\begin{Satz}
  Jedes irreduzible Element eines Hauptidealrings ist ein Primelement.
\end{Satz}
\begin{proof}[Beweis]
  Sei $p$ ein irreduzibles Element, welches das Produkt $ab$ teilt, aber nicht
  $a$ teilt. Das von $a$ und $p$ erzeugte Ideal ist ein Hauptideal $(d)$, es
  gibt also $u,v\in R$ mit $ua+vp=d$. Wegen $(p)\subset (d)$ gilt $d\mid
  p$. Damit ist $d$ eine Einheit, oder von der Form $up$ mit einer Einheit
  $u$. Aber $d\mid a$, und somit $p\mid a$, im Widerspruch zur Voraussetzung
  an $a$. Daher ist $d$ eine Einheit, und ohne Einschr\"ankung
  $d=1$. Multiplikation mit $b$ liefert die Behauptung.
\end{proof}  
Wir kommen nun zum wichtigen Begriff des faktoriellen Rings.
\begin{Definition}
  Ein Integrit\"atsbereich $R$ hei\ss t \ei{faktoriell}, wenn jedes Element
  $\ne0$ entweder eine Einheit ist, oder ein Produkt endlich vieler
  Primelemente ist.
\end{Definition}
Nat\"urlich ist in einem faktoriellen Ring jedes irreduzible Element ein
Primelement. Die wichtigste Eigenschaft eines faktoriellen Rings ist die
eindeutige Primfaktorzerlegung.
\begin{Satz}
  Sei $R$ ein faktorieller Ring, $0\ne a\in R$ keine Einheit, und
  $a=p_1p_2\dots p_r=q_1q_2\dots q_s$ zwei Zerlegungen in Primelemente $p_i,
  q_i$. Dann gilt $r=s$, und die $p_i$ bzw.\ $q_i$ stimmen bis auf Reihenfolge
  und Multiplikation mit Einheiten \"uberein. (D.h., es gibt eine Permutation
  $\sigma$ von $\{1,2,\dots,r\}$ und Einheiten $u_1,\dots,u_r$ mit
  $p_{i}=u_iq_{i^\sigma}$ f\"ur alle $i$.)
\end{Satz}
\begin{proof}[Beweis]
  Ohne Einschr\"ankung gilt $r\ge s$. Wir beweisen die Aussage durch
  vollst\"andige Induktion \"uber $r$. F\"ur $r=1$ gibt es nichts zu
  beweisen. Wir schlie\ss en nun von $r-1$ auf $r$. Da $p_r$ das Produkt
  $q_1q_2\dots q_s$ teilt, folgt durch mehrfache Anwendung der Primeigenschaft
  von $p_r$, dass $p_r$ einen der Faktoren $q_1,q_2,\dots,q_s$ teilt. Durch
  Umbenennung der $q_i$ d\"urfen wir annehmen, dass $p_r\mid q_s$, also
  $q_s=up_r$. Als Primelement ist $q_s$ irreduzibel, daher ist $u$ eine
  Einheit. Die Behauptung folgt nun aus der Richtigkeit f\"ur $r-1$, angewandt
  auf $p_1p_2\dots p_{r-1}=q_1q_2\dots q_{s-2}(uq_{s-1})$. 
\end{proof}
Nat\"urlich hat jede ganze Zahl eine Zerlegung in irreduzible Zahlen. Da
$\ZZZ$ ein Hauptidealring ist, stimmen irreduzible und Primelemente \"uberein,
und obiger Satz liefert die eindeutige Primfaktorzerlegung in $\ZZZ$.

Die Primeigenschaft der Primzahlen hatten wir schon fr\"uher verwendet
(Korollar \ref{K:ZnZ}). Ferner machten wir mehrfach von der eindeutigen
Primfaktorzerlegung in $\ZZZ$ Gebrauch. Man mache sich klar, dass der gerade
nachgetragene Beweis nicht auf einem Zirkelschluss beruht!

Wir wissen auch, dass Polynomringe \"uber einem K\"orper Hauptidealringe
sind. Ferner beweist man mittels der Gradfunktion, dass jedes Polynom ein
Produkt irreduzibler Polynome ist. Das liefert das wichtige
\begin{Korollar}
  Sei $R=\ZZZ$, oder $R$ der Polynomring $K[X]$ f\"ur einen K\"orper $K$. Dann
  stimmen in $R$ irreduzible Elemente und prime Elemente \"uberein.
  Insbesondere ist $R$ faktoriell.
\end{Korollar}
\subsection{Inhalt von Polynomen, Lemma von Gau\ss}
Im folgenden sei $R$ ein faktorieller Integrit\"atsbereich, und $K$ der
Quotientenk\"orper von $R$. Es mag das Verst\"andnis erleichtern, wenn man
stets das Beispiel $R=\ZZZ$, $K=\QQQ$ im Auge beh\"alt.

Sei $0\ne a\in K$, und $p\in R$ ein Primelement von $R$. Aus der eindeutigen
Primfaktorzerlegung folgt schnell, dass man eine Darstellung
$a=p^m\frac{u}{v}$ mit $m\in\ZZZ$, $u\in R$, $0\ne v\in R$ finden kann, so
dass $p$ weder $u$ noch $v$ teilt. Dabei ist $m$ unabh\"angig von der Wahl des
Paares $u,v$. Wir schreiben $v_p(a):=m$, und setzen $v_p(0):=\infty$. Es gilt
$v_p(ab)=v_p(a)v_p(b)$. Ist $0\ne f=a_0+a_1^X+\dots a_nX^n\in K[X]$, dann
setzen wir $v_p(f):=\min_iv_p(a_i)$. Sei $P$ eine Menge von Primelementen von
$R$, so dass jedes Primelement aus $R$ sich von genau einem Element aus $P$
nur um eine Einheit unterscheidet. Ist $0\ne f\in K[X]$, dann nennt man
\[
I(f):=\prod_{p\in P}p^{v_p(f)}
\]
den \ei{Inhalt} von $f$.

Ist beispielsweise $R=\ZZZ$, und $P\subset\NNN$ die Menge der Primzahlen, dann
ist f\"ur $f\in\ZZZ[X]$ der Inhalt $I(f)$ gerade der gr\"o\ss te gemeinsame
Teiler der Koeffizienten von $f$.

Der Inhalt $I(f)$ h\"angt von der Wahl von $P$ ab, und ist nur bis auf
Einheiten eindeutig bestimmt.

Ist $0\ne a\in K$ und $0\ne f\in K[X]$, dann gilt $I(af)=I(a)I(f)$. Hieraus
folgt $I(\frac{1}{I(f)}f)=\frac{1}{I(f)}I(f)=1$. Es gibt also $\gamma\in K$
mit $f=\gamma g$, so dass $I(g)=1$ gilt. Polynome mit Inhalt $1$ nennt man
\ei{primitiv}. Man beachte, dass ein primitives Polynom Koeffizienten aus $R$
hat.

\begin{Satz}[Gau\ss\ Lemma]
  Seien $0\ne f,g\in K[X]$. Dann gilt $I(fg)=I(f)I(g)$.
\end{Satz}
\begin{proof}[Beweis]
  Schreibe $f=\alpha f_1$, $g=\beta g_1$ mit $\alpha,\beta\in K$ und $f_1,g_1$
  primitiv. Wegen $I(\alpha\beta)=I(\alpha)I(\beta)$ bleibt zu zeigen, dass
  auch $f_1g_1$ primitiv ist. Wegen $f_1g_1\in R[X]$ ist das gleichbedeutend
  damit, das $f_1g_1$ f\"ur jedes Primelement $p\in R$ einen nicht durch $p$
  teilbaren Koeffizienten hat. Der Faktorring $\bar{R}=R/(p)$ ist ein
  Integrit\"atsbereich. Der nat\"urliche Homomorphismus $\bar{ }:R\to R/(p)$
  setzt sich fort zu einem Homomorphismus $\bar{ }:R[X]\to\bar{R}[X]$. W\"aren
  alle Koeffizienten von $f_1g_1$ durch $p$ teilbar, dann folgte
  $\overline{f_1g_1}=0$, also $\bar{f_1}\bar{g_1}=0$. Da $\bar{R}$ ein
  Integrit\"atsbereich ist, folgt ohne Einschr\"ankung $\bar f_1=0$. Das
  bedeutet aber, dass alle Koeffizienten von $f_1$ durch $p$ teilbar sind, im
  Widerspruch zu $I(f_1)=1$.
\end{proof}
Das Gau\ss\ Lemma hat eine Reihe von Konsequenzen. Wenn zum Beispiel ein
Polynom $f\in R[X]$ in Faktoren aus $K[X]$ zerf\"allt, dann bekommt man im
wesentlichen eine gleiche Faktorisierung \"uber $R$. Genauer gilt.
\begin{Korollar}
  Sei $f=gh$ mit $f\in R[X]$ und $g,h\in K[X]$. Setze $\gamma_g=I(g)$,
  $\gamma_h=I(h)$, also $g=\gamma_g g_1$, $h=\gamma_h h_1$ mit $g_1,h_1$
  primitiv. Dann gilt $f=(\gamma_g\gamma_h)g_1h_1$ mit $\gamma_g\gamma_h\in
  R$. Insbesondere gilt: Ist ein primitives Polynom irreduzibel in $R[X]$,
  dann ist es auch irreduzibel in $K[X]$.
\end{Korollar}
Eine wichtige Folgerung ist
\begin{Korollar}
  Sei $R$ ein faktorieller Ring. Dann ist $R[X]$ faktoriell. Die Primelemente
  von $R[X]$ bestehen aus den Primelementen von $R$, zusammen mit den
  irreduziblen und primitiven Polynomen aus $R[X]$.
\end{Korollar}
\begin{proof}[Beweis]
  Wir m\"ussen zeigen, dass jedes irreduzible Element $f\in R[X]$ prim ist,
  und die angebene Form besitzt. Wir schreiben $f=\gamma g$ mit $\gamma\in R$
  und $g$ primitiv. Da $f$ irreduzibel ist, ist entweder $g$ irreduzibel und
  $\gamma$ eine Einheit, oder $g$ ein konstantes Polynom, und wegen $I(g)=1$
  eine Einheit, und $\gamma$ ein Primelement.
  
  Es bleibt zu zeigen, dass $f$ prim ist. Sei also $f$ ein irreduzibles
  Element in $R[X]$, und ein Teiler von $uv$ mit $u,v\in R[X]$. Wir
  unterscheiden zwei F\"alle: $f$ ist ein Primelement aus $R$. Dann ist $I(f)$
  auch prim in $R$, und $I(f)$ teilt $I(u)I(v)$. Ohne Einschr\"ankung ist dann
  $I(f)$ ein Teiler von $I(u)$. Dann ist aber $I(\frac{u}{f})\in R$, also
  $\frac{u}{f}\in R[X]$, und somit $f\mid u$.
  
  Im zweiten Fall ist $f$ irreduzibel und primitiv. Wegen obigen Korollars ist
  $f$ auch irreduzibel in $K[X]$, also auch ein Primelement in $K[X]$, da
  $K[X]$ faktoriell ist. In $K[X]$ ist also $f$ ein Teiler von $uv$, also auch
  von $u$ oder $v$. Sei $v=hf$ mit $h\in K[X]$. Wegen $I(v)=I(h)I(f)=I(h)$
  gilt $h\in R[X]$, somit ist $u$ in $R[X]$ durch $f$ teilbar, und die
  Behauptung folgt.
\end{proof}
Durch $K[X_1,X_2,\dots,X_n]:=K[X_1,X_2,\dots,X_{n-1}][X_n]$ definiert man
iterativ einen Polynomring in $n$ Variablen. Ein Folge des obigen Satzes ist
\begin{Korollar}
  Sei $K$ ein K\"orper. Dann ist der Polynomring $K[X_1,X_2,\dots,X_n]$ in $n$
  Variablen faktoriell.
\end{Korollar}
\begin{Bemerkung*}
  Der Polynomring $K[X_1,X_2,\dots,X_n]$ ist f\"ur $n\ge2$ kein
  Hauptidealring, z.B.\ weil das von $X_1,\dots,X_n$ erzeugte Ideal kein
  Hauptideal ist.
\end{Bemerkung*}
\subsection{Das Irreduzibilit\"atskriterium von Eisenstein}
Im allgemeinen ist es von Polynomen hohen Grades oft schwierig, die
Irreduzibilit\"at zu beweisen. Es gibt allerdings ein einfaches Kriterium, das
vor allem in theoretischem Kontext manchmal funktioniert.
\begin{Satz}[Eisenstein]
  Sei $R$ ein Integrit\"atsbereich mit Quotientenk\"orper $K$, $p\in R$ ein
  Primelement und $f(X)=a_0+a_1X+\dots+a_nX^n$ ein Polynom mit folgenden
  Eigenschaften.
  \begin{itemize}
  \item[(i)] F\"ur $0\le i\le n-1$ ist $a_i$ durch $p$ teilbar.
\item[(ii)] $a_n$ ist nicht durch $p$ teilbar.
\item[(iii)] $a_0$ ist nicht durch $p^2$ teilbar.
  \end{itemize}
  Dann ist $f$ irreduzibel in $R[X]$. Ist $R$ faktoriell, dann ist $f$ auch
  irreduzibel in $K[X]$.
\end{Satz}
\begin{proof}[Beweis]
  Sei $f$ reduzibel in $R[X]$, also $f=uv$ mit nicht-konstanten Polynomen
  $u,v\in R[X]$. Setze die Restklassenabbildung $R\to R/(p)=\bar R$ fort zu
  einem Ringhomomorphismus $\bar{}:R[X]\to \bar{R}[X]$ via $X\mapsto X$. Aus
  $f=uv$ folgt $\bar{f}=\bar{u}\bar{v}$. Die Voraussetzungen (i) und (ii)
  liefern $\bar{f}=\overline{a_n}X^n$ mit $\overline{a_n}\ne0$. Wegen der
  eindeutigen Polynomfaktorisierung \"uber dem Integrit\"atsring $\bar{R}$
  haben $\bar u$ und $\bar v$ die Form $\bar{u}=b_rX^r$ und $\bar{v}=c_sX^s$.
  Insbesondere sind die Absolutterme von $u$ und $v$ durch $p$ teilbar, und
  damit ist $a_0$ durch $p^2$ teilbar, im Widerspruch zu (iii).
  
  Der zweite Teil folgt aus dem Gau\ss\ Lemma, denn wenn $R$ faktoriell ist,
  dann sind in $R[X]$ irreduzible Polynome auch in $K[X]$ noch irreduzibel.
\end{proof}
Man nennt ein Polynom mit den Eigenschaften (i)--(iii) auch ein
\ei{Eisenstein-Polynom}. Eine wichige Anwendung ist
\begin{Satz}\label{S:Eisenstein2}
  Sei $p$ eine Primzahl. Dann ist das Polynom $1+X+X^2+\dots+X^{p-1}$ \"uber
  $\QQQ$ irreduzibel.
\end{Satz}
\begin{proof}[Beweis]
  Sei $f(X)=1+X+X^2+\dots+X^{p-1}=\frac{X^p-1}{X-1}$. Nat\"urlich ist $f(X)$
  genau dann irreduzibel, wenn $f(X+1)$ irreduzibel ist. Es gilt
\[
f(X+1)=\frac{(X+1)^p-1}{X}=\sum_{k=1}^p\binom{p}{k}X^{k-1}.
\]
Die Binomialkoeffizienten $\binom{p}{k}$ sind f\"ur $k=1,2,\dots,p-1$ durch
$p$ teilbar, ferner ist der konstante Koeffizient von $f(X+1)$ gleich
$\binom{p}{1}=p$. Alle Voraussetzungen des Eisensteinkriteriums sind also
erf\"ullt.
\end{proof}
\paragraph{Beispiele}
\begin{itemize}
\item Meistens verwendet man das Eisensteinkriterium f\"ur $R=\ZZZ$. So sind
  etwa die Polynome $X^n-7$ ($p=7$), $X^{100}-18X+21$ ($p=3$) und
  $X^4+16X^2+2$ ($p=2$) irreduzibel.
\item Eine andere wichtige Anwendung findet das Kriterium, wenn $R=K[Y]$ ein
  Polynomring ist. So ist zum Beispiel das Polynom $f(X,Y)=X^7+Y^3-1$ \"uber
  $\CCC$ irreduzibel: Sei $R=C[Y]$. Das in $R$ prime Element $Y-1$ teilt
  $Y^3-1$, aber $Y^3-1$ ist nicht durch $(Y-1)^2$ teilbar. Damit ist $f$,
  aufgefasst als Polynom in $X$ \"uber dem Koeffizientenring $R$, ein
  Eisensteinpolynom bez\"uglich $Y-1$. Daher ist $f$ \"uber $R$ irreduzibel.
  Da der Inhalt von $f$ gleich $1$ ist, ist $f$ auch irreduzibel \"uber
  $\CCC$.
\end{itemize}
\subsection{Kryptographie}
Eine wichtige moderne Anwendung der Algebra ist die vertrauliche \"Ubertragung
von Nachrichten. W\"ahrend das fr\"uher fast nur milit\"arisch verwendet
wurde, w\"are unser heutiger Alltag ohne kryptographische Methoden nicht
denkbar. Das Bezahlen mit Geld- oder Kreditkarten, Mobilfunktelefonie, Kauf
und Verkauf \"uber das Internet, elektronische Steuererkl\"arung und vieles
mehr basiert auf kryptographischen Methoden. Dabei geht es stets darum, dass
zwei Partner $S$ und $E$ Nachrichten austauschen, und diese vorher so
verschl\"usseln, dass nur der Partner sie entschl\"usseln kann. Bei den
modernen Anwendungen kommt noch eine weitere wichtige Forderung hinzu: $S$ und
$E$ sollten sich \textbf{\"offentlich} auf eine Methode einigen k\"onnen, also
ohne erst mal \"uber einen sicheren Kanal einen Code festzulegen. Es sieht auf
den ersten Blick sicher so aus, dass so etwas unm\"oglich ist, und bis in die
70er Jahre glaubte man das auch. Dennoch ver\"offentlichten 1976 Diffee und
Hellman eine raffinierte Idee, die diese beiden W\"unsche erf\"ullt. Eine
Variante gaben 1977 \textbf{R}ivest, \textbf{S}hamir und \textbf{A}dleman an,
das heute so genannten \ei{RSA-Verfahren}. Mittlerweile wei\ss\ man, dass
beide diese Durchbr\"uche schon wenige Jahre vorher gefunden wurden, aber ihre
Bedeutung nicht erkannt wurde, und zus\"atzlich noch zur Verschlusssache
erkl\"art wurden. Das RSA-Verfahren z.B.\ entdeckte bereits 1973 der britische
Mathematiker Clifford Cocks. Da beim RSA-Verfahren der Empf\"anger die
Parameter zur Verschl\"usselung \"offentlich macht, nennt man die Methode auch
ein \ei{public key cryptosystem}.

Im folgenden bedeutet $a\equiv b\pmod{m}$ stets, dass $a-b$ durch $m$ teilbar
ist. F\"ur $n\in\NNN$ sei $\varphi(n)$ die Ordnung der Einheitengruppe von
$\ZZZ/n\ZZZ$. Ist $n=\prod p_i^{e_i}$, so wissen wir bereits, dass
$\varphi(n)=\prod (p_i-1)p_i^{e_i-1}$.

Zur Vorbereitung beweisen wir ein kleines
\begin{Lemma}
  Sei $n\in\NNN$ quadratfrei, und $k\in\NNN$. F\"ur alle $a\in\ZZZ/n\ZZZ$ gilt
  dann $a^{k\varphi(n)+1}=a$.
\end{Lemma}
\begin{proof}[Beweis]
  Sei $A\in\ZZZ$ ein Repr\"asentant von $a$. Sei $p\mid n$. Falls $p\mid A$,
  dann gilt nat\"urlich $A^{k\varphi(n)+1}\equiv A\pmod{p}$. Im Fall $p\nmid
  A$ gilt $A^{p-1}\equiv 1\pmod{p}$, also $A^{k\varphi(n)}\equiv 1\pmod{p}$,
  da $p-1$ ein Teiler von $\varphi(n)$ ist. Es folgt $A^{k\varphi(n)+1}\equiv
  A\pmod{p}$. Daher ist $A^{k\varphi(n)+1}-A$ durch alle Primteiler von $n$
  teilbar, und die Behauptung folgt aus der Quadratfreiheit von $n$.
\end{proof}
\paragraph{Das RSA-Verfahren.}
Im folgenden beschreiben wir das RSA-Verfahren. Ein Sender $S$ m\"ochte einem
Empf\"anger $E$ eine Nachricht schicken. Dazu w\"ahlt $E$ zwei gro\ss e
verschiedene Primzahlen $p$ und $q$, und setzt $n=pq$. Ferner w\"ahlt $E$ eine
nat\"urliche Zahl $s\le\varphi(n)$, die teilerfremd zu $\varphi(n)$ ist.  Da
$s$ teilerfremd ist zu $\varphi(n)$, ist $s$ in $\ZZZ/\varphi(n)\ZZZ$
invertierbar, es gibt also $t\in\ZZZ$ mit $st\equiv1\pmod{\varphi(n)}$.  Das
Paar $(n,s)$ macht $E$ \"offentlich (zum Beispiel auf seiner privaten
Homepage), die Primfaktoren $p$ und $q$ von $n$ sowie $t$ h\"alt $E$ geheim.
Aus Sicherheitsgr\"unden kann $E$ sogar $p$ und $q$ nach der Berechnung von
$t$ l\"oschen.

$S$ verwendet zum Versenden seiner Nachricht als ``W\"orter'' die Elemente aus
dem Restklassenring $\ZZZ/n\ZZZ$, sowohl f\"ur den Klartext, als auch f\"ur
die Verschl\"usselung. Das Klartextwort $a\in\ZZZ/n\ZZZ$ verschl\"usselt $S$
hierbei zu $a^s$, der Empf\"anger $E$ (und jeder, der die Nachricht eventuell
mitliest) erh\"alt also $b=a^s$. Der Empf\"anger $E$ berechnet nun $b^t$ in
$\ZZZ/n\ZZZ$.

Wegen des Lemmas gilt in $\ZZZ/n\ZZZ$ also
\[
b^t=a^{st}=a,
\]
$E$ erh\"alt also den Klartext $a$ zur\"uck.

Warum nun h\"alt man das Verfahren bei geeigneter Wahl der Parameter f\"ur
sicher? Ein Zuh\"oher, der das Wort $b=a^s$ abf\"angt, und auch $n$ und $s$
kennt, m\"usste $t$ kennen um $a$ aus $b$ zu rekonstruieren. Wegen
$st\equiv1\pmod{\varphi(n)}$ kann man $t$ bestimmen, wenn man $\varphi(n)$
kennt. Umgekehrt vermutet man, konnte das aber noch nicht beweisen, dass es
keine bessere Methode gibt $t$ zu bestimmen als vorher erst mal $\varphi(n)$
zu ermitteln. Wegen $\varphi(n)=\varphi(pq)=(p-1)(q-1)=n-p-q+1$ ist die
Bestimmung von $\varphi(n)$ \"aquivalent zur Bestimmung von $p$ und $q$, also
der Faktorisierung von $n$.

Man kennt bis heute keine Methode, allgemein Zahlen mit \"uber $150$ Ziffern
(Dezimaldarstellung) zu faktorisieren, auch dann nicht, wenn man schon a
priori wie in unserem Fall wei\ss, dass genau zwei Primfaktoren auftreten.

Andererseits kann man heutzutage von Zahlen mit einigen tausend Stellen
schnell entscheiden, ob sie Primzahlen sind. Diese Methoden beruhen nicht
darauf, dass man versucht, die gegebene Zahl zu faktorisieren, sondern
benutzen z.B.\ die Tatsache, dass $n\in\NNN$ genau dann prim ist, wenn die
Einheitengruppe von $\ZZZ/n\ZZZ$ die Ordnung $n-1$ hat.

Beim RSA-Verfahren sind Potenzen $a^s$ und $b^t$ mit eventuell sehr gro\ss en
Exponenten zu berechnen. Allgemein kann man in einem Monoid eine Potenz $r^n$
mit deutlich weniger als $n-1$ Multiplikationen berechnen: Dazu schreibt man
$n$ in der Bin\"arentwicklung, also $n=\sum 2^{e_i}$ mit verschiedenen
$e_i\in\NNN_0$. Dabei ist $2^{e_i}\le n$, also $e_i\le\log n/\log 2$. Die
Potenz $r^{2^e}$ l\"asst sich durch $e$-faches Quadrieren berechnen, und wegen
$r^n=\prod r^{2^{e_i}}$ l\"asst sich $r^n$ im wesentlichen mit h\"ochstens
$2\log n/\log 2$ Multiplikationen bestimmen.

Diese einfache Beobachtung liefert im \"ubrigen auch einen schnellen Test, mit
dem man in den meisten F\"allen schnell feststellen kann, wenn eine Zahl nicht
prim ist: F\"ur Primzahlen $p$ gilt ja $a^{p-1}\equiv 1\pmod{p}$, wenn $a$
nicht durch $p$ teilbar ist. M\"ochte man eine Zahl $n$ auf Primeigenschaft
testen, dann \"uberpr\"uft man, ob $a^{n-1}\equiv 1\pmod{n}$ gilt z.B.\ f\"ur
$a=2$. Ist das nicht der Fall, dann kann $n$ keine Primzahl sein. Ist es
jedoch der Fall, dann kann man weitere Werte f\"ur $a$ ausprobieren. Leider
gibt es Nicht-Primzahlen $n$, die $a^{n-1}\equiv 1\pmod{n}$ f\"ur alle zu $n$
teilerfremde $a\in\ZZZ$ erf\"ullen, das kleinste Beispiel ist $561=3\cdot
11\cdot 17$.
\section{K\"orper}
Dieser Abschnitt ist den K\"orpern gewidmet. Wichtige Beispiele sind die
K\"orper $\QQQ$, $\RRR$ und $\CCC$ der rationalen, reellen und komplexen
Zahlen. F\"ur jede Primzahl $p$ haben wir schon einen bis auf Isomorphie
eindeutigen endlichen K\"orper $\FFF_p=\ZZZ/p\ZZZ$ mit $p$ Elementen kennen
gelernt. Ist $K$ ein K\"orper, dann hat der Polynomring $K[X]$ keine
Nullteiler, und somit ist der Quotientenk\"orper $K(X)$ selber wieder ein
K\"orper. Auch in der Analysis spielen gewisse Funktionenk\"orper eine
wichtige Rolle. So beweist man etwa in der Funktionentheorie, dass die auf
einem Gebiet meromorphen Funktionen einen K\"orper bilden.
\subsection{K\"orpererweiterungen, K\"orpergrade und Homomorphismen}
Seien $K$ und $L$ K\"orper mit $K\subseteq L$. Man nennt dann $K$ einen
\ei{Teilk\"orper} von $L$, und $L$ einen \ei{Oberk\"orper} von $K$. Das Paar
$L,K$ beschreibt eine \ei{K\"orpererweiterung}, und man schreibt daf\"ur meist
$L|K$ (und spricht von ``$L$ \"uber $K$'').

In dieser Situation ist $L$ ein Vektorraum \"uber dem K\"orper $K$, das
Produkt des ``Skalars'' $\alpha\in K$ mit dem ``Vektor'' $v\in L$ ist einfach
das in $L$ gebildete Produkt $\alpha v$. Die Dimension des Vektorraums $L$
\"uber $K$ hei\ss t der \ei{K\"orpergrad} von $L|K$. F\"ur den K\"orpergrad
$\dim_KL$ schreibt man $[L:K]$. Man nennt die Erweiterung $L|K$ endlich, falls
$[L:K]<\infty$.
\paragraph{Beispiele}
\begin{itemize}
\item Jede komplexe Zahl ist eine eindeutige reelle Linearkombination von $1$
  und $i=\sqrt{-1}$, daher gilt $[\CCC:\RRR]=2$.
\item Die rationalen Zahlen sind abz\"ahlbar, die reellen Zahlen sind es
  nicht. H\"atte $\RRR$ eine endliche $\QQQ$-Basis, dann w\"are mit einem
  einfachen Argument auch $\RRR$ abz\"ahlbar, was nicht der Fall ist. Daher
  gilt $[\RRR:\QQQ]=\infty$.
\item Die Monome $X^0, X^1, X^2,\dots$ des Polynomrings $K[X]$ sind \"uber $K$
  linear unabh\"angig, insbesondere gilt $[K(X):K]=\infty$.
\end{itemize}
Eine fundamentale Eigenschaft des K\"orpergrades ist
\begin{Lemma}
  Seien $M|L$ und $L|K$ K\"orpererweiterungen, und $\{m_i|i\in I\}$ sowie
  $\{l_j|j\in J\}$ Basen von $M$ \"uber $L$ bzw.\ $L$ \"uber $K$. Dann ist
  $\{m_il_j|i\in I,j\in J\}$ eine Basis von $M$ \"uber $K$. Insbesondere gilt
  $[M:K]=[M:L][L:K]$.
\end{Lemma}
\begin{proof}[Beweis]
  Jedes Element aus $M$ ist eine $L$-Linearkombination der $m_i$, und jeder
  der Koeffizienten aus dieser Linearkombination wiederum ist eine
  $K$-Linearkombination der $l_i$. Daher bilden die $m_il_j$ ein
  Erzeugendensystem von $M$ als $K$-Vektorraum.

Es bleibt die lineare Unabh\"angigkeit der $m_il_j$ \"uber $K$ zu zeigen. Dazu
sei
\[
\sum_{i,j}\alpha_{ij}m_il_j=0|\text{ mit }\alpha_{ij}\in K.
\]
Setze $\beta_i:=\sum_{j}\alpha_{ij}l_j$. Dann gilt
\[
\sum_{i}\beta_im_i=0\text{ mit }\beta_i\in L.
\]
Wegen der linearen Unabh\"angigkeit der $m_i$ \"uber $L$ folgt $\beta_i=0$
f\"ur alle $i$. Dies ergibt $\alpha_{ij}=0$ wegen der linearen
Unabh\"angigkeit der $l_j$ \"uber $K$.
\end{proof}
In der Gruppen- und Ringtheorie spielten Kerne von Homomorphismen eine
wichtige Rolle. Bei K\"orpern gibt es so etwas nicht:
\begin{Lemma}
  Jeder Homomorphismus zwischen K\"orpern ist injektiv.
\end{Lemma}
\begin{proof}[Beweis]
  Sei $\phi:K\to L$ ein Homomorphismus zwischen $K$ und $L$. Ist $0\ne a\in
  K$, dann gilt $1=a\frac{1}{a}$, also
  $1=\phi(1)=\phi(a\frac{1}{a})=\phi(a)\phi(\frac{1}{a})$, also
  $\phi(a)\ne0$. Somit besteht der Kern von $\phi$ nur aus der $0$.
\end{proof}
Wir werden h\"aufig die Situation haben, dass $L$ und $M$
K\"orpererweiterungen von $K$ sind. Ein Homomorphismus $\phi:M\to L$ hei\ss t
dann \ei{$K$-Homomorphismus}, wenn die Einschr\"ankung von $\phi$ auf $K$ die
Identit\"at ist. Wegen der Injektivit\"at von $\phi$ spricht man auch oft von
einer \ei{$K$-Einbettung}. Ist $\phi$ sogar bijektiv, dann hei\ss t $\phi$ ein
\ei{$K$-Isomorphismus}.
\subsection{Algebraische und transzendente Elemente}
Sei $L|K$ eine K\"orpererweiterung. F\"ur $a\in L$ bezeichnet man mit $K(a)$
den kleinsten Zwischenk\"orper von $L|K$, der $a$ enth\"alt. Man spricht auch
von dem von $K$ und $a$ erzeugten K\"orper. Allgemeiner sei $a_1,a_2,\dots\in
L$. Dann ist $K(a_1,a_2,\dots)$ der kleinste Zwischenk\"orper von $L|K$, der
alle $a_i$ enth\"alt.

Ein Element $a\in L$ hei\ss t \ei{algebraisch} \"uber $K$, wenn es ein
nicht-triviales Polynom $f(X)\in K[X]$ gibt mit $f(a)=0$. Im folgenden sei das
vorausgesetzt. Die Menge der Polynome $g\in K[X]$ mit $g(a)=0$ bildet ein
Ideal von $K[X]$. Da $K[X]$ ein Hauptidealring ist, gibt es ein normiertes
Polynom $f$ kleinsten Grades mit $f(a)=0$. Man nennt dieses $f$ das
\ei{Minimalpolynom} von $a$ \"uber $K$.  Wegen der Nullteilerfreiheit von
$K[X]$ ist $f$ irreduzibel. Insbesondere ist das Ideal $(f)$ maximal, also
$K[X]/(f)$ ein K\"orper. Sei $K[a]$ der von $K$ und $a$ erzeugte Ring. Die
Abbildung $K[X]\to K[a]$, $g\mapsto g(a)$ ist ein Ringepimomorphismus mit Kern
$(f)$. Da $(f)$ ein maximales Ideal ist, ist $K[a]$ isomorph zum K\"orper
$K[X]/(f)$. Insbesondere ist $K[a]$ ein K\"orper.  Hieraus folgt, dass $K[a]$
der kleinste Zwischenk\"orper von $L|K$ ist, der $a$ enth\"alt.

Da $K[a]$ ein K\"orper ist, liegen auch Elemente wie $\frac{1}{a}$ in $K[a]$.
Wie dr\"uckt man solche Elemente durch $1,a,a^2,\dots$ aus? Dazu verwendet man
das Minimalpolynom $f(X)=\alpha_0+\alpha_1X+\dots+\alpha_nX^n$ von $f$.
Ersetzt man $X$ durch $a$ und dividiert durch $\alpha_0\ne0$, so erh\"alt man
$\frac{1}{a}=-\frac{1}{\alpha_0}(\alpha_1+\alpha_2a+\dots+\alpha_{n}a^{n-1})$.
Sei $h(X)\in K[X]$ mit $h(a)\ne0$. M\"ochte man nun $\frac{1}{h(a)}$
bestimmen, so geht man folgenderma\ss en vor: Wegen $h(a)\ne0$ und der
Irreduzibilit\"at von $f$ sind $h$ und $f$ teilerfremd. Daher findet man
Polynome $u,v\in K[X]$ mit $uh+vf=1$. Es folgt $1=u(a)h(a)+v(a)f(a)=u(a)h(a)$,
also $\frac{1}{h(a)}=u(a)$. Die Polynome $u$ und $v$ findet man z.B.\ mit dem
Euklidischen Algorithmus. (Diese \"Uberlegungen h\"atte man auch alternativ zu
den obigen verwenden k\"onnen, um zu zeigen, dass $K[a]$ ein K\"orper ist.)

Hat das Minimalpolynom $f$ von $a$ den Grad $n$, dann bilden die Elemente
$a^i$, $i=0,1,\dots,n-1$ offenbar eine Basis des $K$-Vektorraums $K[a]$. Wegen
$K[a]=K(a)$ gilt dann also $[K(a):K]=n$. Wir fassen die \"Uberlegungen
zusammen:
\begin{Lemma}
  Sei $L|K$ eine K\"orpererweiterung, und $a\in L$ algebraisch mit
  Minimalpolymom $f\in K[X]$. Dann ist der K\"orper $K(a)=K[a]$ $K$-isomorph
  zu $K[X]/(f)$, und der Grad von $K(a)|K$ ist der Grad von $f$.
\end{Lemma}
Man nennt eine Erweiterung $L|K$ algebraisch, wenn jedes Element aus $L$
algebraisch \"uber $K$ ist.
H\"aufig ben\"otigt man die folgende einfache Beobachtung.
\begin{Lemma}
  Endliche K\"orpererweiterung sind algebraisch.
\end{Lemma}
\begin{proof}[Beweis]
  Sei $L|K$ eine K\"orpererweiterung mit $[L:K]=<\infty$, und $a\in L$. Dann
  sind die $n+1$ Elemente $1,a,a^2,\dots,a^n$ linear abh\"angig \"uber $K$,
  was zu einem Polynom $g\ne0$ f\"uhrt mit $g(a)=0$.
\end{proof}
Sei nun $L|K$ wieder eine K\"orpererweiterung, aber $a\in L$ sei nun nicht
algebraisch \"uber $K$. In diesem Fall nennt man $a$ \ei{transzendent} \"uber
$K$. Der Ringhomorphismus $K[X]\to K[a]$ ist in diesem Fall bijektiv, und
setzt sich daher zu einem $K$-Isomorphismus zwischen den K\"orpern $K(X)$ und
$K(a)$ fort. Eine \ei{transzendente Zahl} ist ein transzendentes Element der
Erweiterung $\CCC|\QQQ$. In der Analysis zeigt man (vielleicht auch sp\"ater
in dieser Vorlesung), dass die Eulersche Konstante $e$ und die Kreiszahl $\pi$
transzendent sind.
\paragraph{Beispiele von K\"orpererweiterungen}
\begin{itemize}
\item Wir sahen schon die Erweiterung $\CCC|\RRR$ vom Grad $2$. Wegen $i^2=-1$
  ist $f(X)=X^2+1$ das Minimalpolynom von $i$ \"uber $\RRR$.
\item Wir wollen die Erweiterung $\QQQ(\sqrt{2},\sqrt{3})|\QQQ$ studieren. Sei
  $L=\QQQ(\sqrt{2})=\QQQ[\sqrt{2}]$. Wegen $\sqrt{2}\notin\QQQ$ gilt
  $[L:\QQQ]\ge2$, und wegen $(\sqrt{2})^2=2=\in \QQQ$ gilt
  $[L:\QQQ]\le2$, also $[L:\QQQ]=2$. Analog folgt
  $[\QQQ(\sqrt{2},\sqrt{3})|L]\le 2$. W\"are $\QQQ(\sqrt{2},\sqrt{3})=L$, dann
  g\"abe es $\alpha,\beta\in\QQQ$ mit $\sqrt{3}=\alpha+\beta\sqrt{2}$. Wegen
  $3+\beta^22-2\beta\sqrt{2}\sqrt{3}=(\sqrt{3}-\beta\sqrt{2})^2=\alpha^2$ und
  der Irrationalit\"at von $\sqrt{2}\sqrt{3}$ folgt $\beta=0$, also
  $\sqrt{3}=\alpha$, im Widerspruch zur Irrationalit\"at von $\sqrt{3}$. Damit
  ist $\{1,\sqrt{3}\}$ eine Basis von $\QQQ(\sqrt{2},\sqrt{3})$ \"uber $L$,
  und $\{1,\sqrt{3},\sqrt{2},\sqrt{2}\sqrt{3}\}$ eine Basis f\"ur
  $\QQQ(\sqrt{2},\sqrt{3})$ \"uber $\QQQ$.
\item Wir behaupten, dass sich die Erweiterung $\QQQ(\sqrt{2},\sqrt{3})|\QQQ$
  schon von einem Element, z.B.\ $a=\sqrt{2}+\sqrt{3}$ erzeugen l\"asst.
  Offenbar gilt $\QQQ(a)\subseteq\QQQ(\sqrt{2},\sqrt{3})$. Die umgekehrte
  Inklusion folgt so: Aus $3=(a-\sqrt{2})^2=a^2-2a\sqrt{2}+2$ folgt
  $\sqrt{2}=\frac{a^2-1}{2a}$, und analog $\sqrt{3}=\frac{a^2+1}{2a}$, also
  $\sqrt{2},\sqrt{3}\in\QQQ(a)$.
\item Die \"Uberlegung aus dem vorigen Beispiel liefert auch ein
  Minimalpolynom von $a$ \"uber $\QQQ$: Quadrieren von $2a\sqrt{2}=a^2-1$
  liefert $8a^2=a^4-2a^2+1$, also $a^4-10a^2+1=0$. Da wir schon wissen, dass
  $\QQQ(a)|\QQQ$ Grad $4$ hat, ist $X^4-10X^2+1$ das Minimalpolynom von $a$
  \"uber $\QQQ$.
\end{itemize}
\subsection{Konstruktion mit Zirkel und Lineal}
In diesem Abschnitt wollen wir pr\"azisieren, was es hei\ss t, mit Zirkel und
Lineal etwas zu konstruieren, und in welchem Sinne die Unl\"osbarkeit einiger
klassischer Konstruktionsaufgaben gezeigt werden kann.

Wir befinden und in der reellen Ebene, in der eine Menge $M$ von Punkten
gegeben ist.

Mit unserem (unendlich langen) Lineal d\"urfen wir eine Gerade durch zwei
verschiedene Punkte ziehen, und mit dem Zirkel d\"urfen wir Kreise zeichnen,
mit Mittelpunkt ein gegebener Punkt, und Radius der Abstand zweier Punkte.

Durch Schnitte der dabei enstehenden Geraden und Kreise erhalten wir weitere
Punkte, die wir zur Menge $M$ hinzunehem, und selbst wieder f\"ur diese
Konstruktionen verwenden.

Sei $c(M)$ die Menge der Punkte, die man aus $M$ in endlich vielen solchen
Schritten gewinnen kann. Da man im ersten Schritt zwei Punkte braucht, setzt
man nat\"urlich $\abs{M}>1$ voraus. Der f\"ur uns wichtige Fall ist
$\abs{M}=2$.

Ist $g$ eine Gerade, und $P$ ein Punkt, so kann man mit Zirkel und Lineal
m\"uhelos ein Lot auf $g$ konstruieren, das durch $P$ geht: Man zeichnet einen
Kreis um $P$ mit einem Radius, der so gro\ss\ ist, dass der Kreis die Gerade
$g$ in $A$ und $P$ schneidet. Die Kreise um $A$ und $B$ durch $P$ schneiden
sich in $Q\ne P$, die Gerade durch $P$ und $Q$ ist das gesuchte Lot. Eine
offensichtliche Modifikation behandelt den Fall, dass $P$ auf $g$ liegt.

\includegraphics[width=80mm]{lot1.eps}

Eine Kleinigkeit m\"ussen wir dabei beachten: Wir d\"urfen nicht mit einem
``beliebigen'' hinreichend gro\ss en Kreis um $P$ starten, da unsere Kreise ja
durch bereits konstruierte Punkte gehen m\"ussen. Den Mangel behebt man
folgenderma\ss en: Die Menge $M$ enth\"alt ja mindestens $2$ Punkte $U$ und
$V$. Durch $U$ und $V$ legen wir eine Gerade, und tragen die Stecke $[UV]$
mehrmals hintereinander auf der Geraden ab. Wir erhalten also Punkte mit
Abst\"anden, die Vielfache der L\"ange $\overline{UV}$ sind. Ein hinreichend
gro\ss es Vielfaches nehmen wir als Radius f\"ur den Kreis um $P$.

Um nun zu einer algebraischen Beschreibung von $M$ zu gelangen, f\"uhren wir
kartesische Koordinaten ein. Dazu nehmen wir zwei verschiedene Punkte $O$ und
$E$ aus $M$. Wir konstruieren $F$, so dass $OF$ senkrecht auf $OE$ steht, und
$\overline{OE}=\overline{OF}$ gilt. Diese drei Punkte bestimmen nun ein
kartesisches Koordinatensystem mit $O=(0,0)$, $E=(1,0)$, $F=(0,1)$.

Nat\"urlich h\"angt das Koordinatensystem von der Wahl der zwei Punkte $O$ und
$E$ (und der zwei M\"oglichkeiten f\"ur $F$) ab, aber wir sehen in K\"urze,
dass das f\"ur unsere Fragen keine Rolle spielt.

Da wir Lote konstruieren k\"onnen, liegt ein Punkt $(a,b)$ genau dann in
$c(M)$, wenn $(a,0)$ und $(b,0)$ in $c(M)$ liegen. Wir setzen
\[
K(M)=\{a|\text{ es gibt ein $b$ mit }(a,b)\in c(M)\}.
\]
Somit besteht $c(M)$ genau aus den Punkten $(a,b)$ mit $a,b\in K(M)$.
\begin{Lemma}
  $K(M)$ ist ein Teilk\"orper von $\RRR$.
\end{Lemma}
\begin{proof}[Beweis]
Es ist klar, dass $K(M)$ eine Untergruppe von $(\RRR,+)$ ist. Ferner ist
$K(M)$ multiplikativ abgeschlossen: Sei $a,b\in K(M)$, und ohne
Einschr\"ankung gelte $a,b>0$. In unserem Koordinatensystem konstruieren wir
die Punkte $A=(a,0)$ und $B=(0,b)$. Die Parallele zu $AF$ durch $B$ schneide
die Gerade $OA$ in $C$. Nach dem Strahlensatz gilt $C=(ab,0)$, d.h.\ wir haben
$ab$ konstruiert.

\includegraphics[width=80mm]{mult.eps}

Um $a^{-1}$ zu konstruieren, konstruiert man wieder $A=(a,0)$. Dann schneidet
man die Parallele zu $AF$, die durch $E$ geht, mit der Geraden $OF$, und
enth\"alt den Schnittpunkt $B=(0,\frac{1}{a})$.
\end{proof}
\begin{Bemerkung*}
  Nat\"urlich liegen auch alle Abst\"ande von Punkten aus $c(M)$ in $K(M)$. Da
  auch die Quotienten dieser Abst\"ande in $K(M)$ liegen, h\"angt $K(M)$
  tats\"achslich nur von $M$ ab, ist also unabh\"angig von der Wahl der Punkte
  aus $M$, die wir zur Festlegung der Koordinaten benutzt hatten.
\end{Bemerkung*}
\begin{Lemma}
  Gilt $0<a\in K(M)$, so ist $\sqrt{a}\in K(M)$.
\end{Lemma}
\begin{proof}[Beweis]
Wir zeichnen auf einer Geraden die Punkte $A$, $B$ und $C$ (mit $B$ zwischen
$A$ und $C$), so dass $\abs{AB}=a$ und $\abs{BC}=1$ gilt. Ein Lot auf $AC$
durch $B$ schneiden wir mit dem Thaleskreis \"uber der Strecke $AC$, und
erhalten $D$. Aus der \"Ahnlichkeit der Dreiecke $ABD$ und $DBC$ folgt
$\abs{BD}=\sqrt{a}$.

\includegraphics[width=80mm]{wurz.eps}
\end{proof}
Wir ben\"otigen ein einfaches Lemma \"uber quadratische K\"orpererweiterungen.
\begin{Lemma}
  Sei $L|K$ eine K\"orpererweiterung mit $[L:K]=2$ und $1+1\ne0$. Dann gibt es
  ein $a\in K$ mit $L=K[\sqrt{a}]$.
\end{Lemma}
\begin{proof}[Beweis]
  Sei $\alpha\in L\setminus K$. Wegen $[K(\alpha):K]>1$ gilt
  $L=K[\alpha]$. Sei $f(X)=X^2+pX+q$ das Minimalpolynom von $\alpha$ \"uber
  $K$. Setze $\beta=\alpha+\frac{p}{2}$. Nat\"urlich gilt
  $L=K[\alpha]=K[\beta]$. Die Behauptung folgt aus
\[
0=\alpha^2+p\alpha+q=(\beta-\frac{p}{2})^2+p(\beta-\frac{p}{2})+q =
\beta^2-(\frac{p^2}{4}-q).
\]
\end{proof}
Das wichtigste vorbereitende Ergebnis ist nun
\begin{Satz}
  Sei $a\in\RRR$. Die folgenden Aussagen sind \"aquivalent.
  \begin{itemize}
  \item[(i)] $a\in K(M)$.
  \item[(ii)] Sei $E$ der von den Koordinaten der Punkte aus $M$ erzeugte
    K\"orper. Dann gibt es eine Kette von K\"orpern $E=E_0\subset E_1\subset
    E_2\subset\dots\subset E_r\subseteq\RRR$ mit $[E_{i+1}:E_{i}]=2$ und $a\in
    E_r$.
  \end{itemize}
\end{Satz}
\begin{proof}[Beweis]
  Die Richtung (ii) nach (i) folgt unmittelbar aus den obigen Betrachtungen.
  Um zu sehen, dass (ii) aus (i) folgt, m\"ussen wir sehen, dass wir bei
  unserer Konstruktion von Schnittpunkten schlimmstenfalls quadratische
  Gleichungen l\"osen m\"ussen. Das ist offensichtlich au\ss er f\"ur den
  Fall, wo wir zwei Kreise schneiden. Sei $(x,y)$ ein Schnittpunkt des Kreises
  mit Radius $r$ um $(a,b)$ mit dem Kreis um $(c,d)$ mit Radius $s$. Dabei
  m\"ogen $a,b,c,d,r$ und $s$ im K\"orper $E_r$ liegen. Wir m\"ussen also
  zeigen, dass es einen K\"orper $E_{r+1}\supseteq E^{r}$ gibt mit
  $[E_{r+1}:E_r]=2$ und $x,y\in E_{r+1}$.

Es gilt also das Gleichungssystem
\begin{align*}
  (x-a)^2+(y-b)^2 &= r^2\\
  (x-c)^2+(y-d)^2 &= s^2
\end{align*}
zu l\"osen. Auf den ersten Blick sieht es so aus, dass die Elimination z.B.\ 
von $y$ auf eine Gleichung vom Grad $4$ f\"ur $x$ f\"uhrt. Das ist jedoch
nicht der Fall: Subtrahieren wir die beiden Gleichungen, dann erhalten wir
eine lineare Gleichung in $x$ und $y$. Diese, zusammen mit einer der
Ausgangsgleichungen, zeigt aber, dass wir zur Berechnung von $x$ und $y$
schlimmstenfalls eine quadratische K\"orpererweiterung ben\"otigen.
\end{proof}
Als Anwendung des Satzes notieren wir zwei Folgerungen. Die erste ist ein
Spezialfall des Satzes.
\begin{Korollar}\label{K:konstr}
  In der reellen Ebene seien zwei Punkte mit Abstand $1$ gegeben. Ferner sei
  $0<a\in\RRR$. Die folgenden Aussagen sind \"aquivalent.
  \begin{itemize}
  \item[(i)] Ausgehend von diesen zwei Punkten k\"onnen wir mit Zirkel und
  Lineal zwei Punkte mit Abstand $a$ konstruieren.
\item[(ii)] Es gibt eine Kette von K\"orpern $\QQQ=E_0\subset E_1\subset
  E_2\subset\dots\subset E_r$ mit $[E_{i+1}:E_{i}]=2$ und $a\in E_r$.
  \end{itemize}
\end{Korollar}
\begin{Korollar}
  In der reellen Ebene seien zwei Punkte mit Abstand $1$ gegeben. Ausgehend
  von diesen zwei Punkten sei $a$ der Abstand zweier Punkte, die wir mit
  Zirkel und Lineal konstruieren k\"onnen. Dann ist $[\QQQ(a):\QQQ]$ eine
  $2$-Potenz.
\end{Korollar}
\begin{proof}[Beweis]
  Mit den Bezeichnungen des vorigen Korollars gilt $[E_r:\QQQ]=2^r$. Da
  $\QQQ(a)$ ein Zwischenk\"orper von $E_r|\QQQ$ ist, folgt die Behauptung aus
  $2^r=[E_r:\QQQ]=[E_r:\QQQ(a)][\QQQ(a):\QQQ]$.
\end{proof}
\begin{Bemerkung*}
  Im letzten Korollar gilt die Umkehrung nicht. Wir werden sp\"ater sehen,
  dass es algebraische $a\in\RRR$ gibt mit $[\QQQ(a):\QQQ]=4$, so dass
  $\QQQ(a)|\QQQ$ keinen echten Zwischenk\"orper hat.
\end{Bemerkung*}
\subsubsection{W\"urfelverdoppelung}
Verdoppelt man das Volumen eines W\"urfels mit Kantenl\"ange $1$, so hat der
neue W\"urfel die Kantenl\"ange $x=\sqrt[3]{2}$. Da das Polynom $X^3-2$ keine
ganze und damit (Gau\ss\ Lemma) keine rationale Nullstelle hat, ist es
irreduzibel. Daher gilt $[\QQQ(x):\QQQ]=3$, und $x$ ist nach Korollar
\ref{K:konstr} nicht mit Zirkel und Lineal konstruierbar.
\subsubsection{Winkeldreiteilung}
Um zu zeigen, dass man nicht alle Winkel mit Zirkel und Lineal dritteln kann,
zeigen wir das speziell f\"ur den $60^\circ$-Winkel.  Wir nehmen an, wir
k\"onnten einen $20^\circ$-Winkel konstruieren. Das ist \"aquivalent zur
Konstruktion von $x=\cos 20^\circ$. Das Additionstheorem $\cos
3\phi=4\cos^3\phi-3\cos\phi$ (man gewinnt das z.B.\ als Realteil von
$\cos3\phi+i\sin3\phi=e^{3i\phi}=(e^{i\phi})^3=(\cos\phi+i\sin\phi)^3$)
liefert, zusammen mit $\cos60^\circ=\frac{1}{2}$, die Beziehung
\[
8x^3-6x-1=0.
\]
Ist $f(X)=8X^3-6X-1$ \"uber $\QQQ$ irreduzibel ist, dann gilt
$[\QQQ(x):\QQQ]=3$, und $x$ ist nach Korollar \ref{K:konstr} nicht
konstruierbar. Bleibt zu zeigen, dass $f$ irreduzibel ist. Dazu zeigen wir,
dass $g(X)=f(X/2)=X^3-3X-1$ irreduzibel ist. W\"are $g$ reduzibel, dann
h\"atte $g$ eine ganzzahlige Nullstelle, die zudem ein Teiler von $1$ ist.
Aber weder $-1$ noch $1$ ist eine Nullstelle von $g(X)$.
\subsubsection{Regul\"are $n$-Ecke}
Da wir nach der vorigen Betrachtung keine $20^\circ$-Winkel konstruieren
k\"onnen, ist es auch nicht m\"oglich, einen $40^\circ$-Winkel zu
konstruieren. Daher kann man auch kein regul\"ares $9$-Eck konstruieren. Man
\"uberlegt sich schnell, dass man regul\"are $n$-Ecke f\"ur $n=2,3,4,5,6$
konstruieren kann. Das regul\"are $7$-Eck kann man hingegen nicht
konstruieren. Dazu zeigen wir zun\"achst
\begin{Lemma}\label{L:nEck}
  Sei $1<n\in\NNN$. Die komplexe Zahl $\zeta=e^{\frac{2\pi i}{n}}$ ist wegen
  $\zeta^n=1$ algebraisch \"uber $\QQQ$. Falls man ein regul\"ares $n$-Eck mit
  Zirkel und Lineal kosntruieren kann, dann ist der Grad von $\zeta$ \"uber
  $\QQQ$ eine Potenz von $2$.
\end{Lemma}
\begin{proof}[Beweis]
  Das $n$-Eck sei konstruierbar. Sei $\zeta=\alpha+i\beta$ mit
  $\alpha,\beta\in\RRR$. Dann ist $[\QQQ(\alpha):\QQQ]$ eine Potenz von $2$.
  Wegen $1=\abs{\zeta}=\alpha^2+\beta^2=\alpha^2-(i\beta)^2=
  \alpha^2-(\zeta-\alpha)^2$ liegt $\zeta$ entweder in $\QQQ(\alpha)$, oder in
  einer quadratischen Erweiterung von $\QQQ(\alpha)$. Die Behauptung folgt.
\end{proof}
Ist zum Beispiel $n=p$ eine Primzahl, dann ist nach Satz \ref{S:Eisenstein2}
das Polynom $f(X)=\frac{X^p-1}{X-1}$ irreduzibel \"uber $\QQQ$. Wegen
$f(\zeta)=0$ ist $f$ das Minimalpolynom von $\zeta$. Daher hat $\zeta$ den
Grad $p-1$. Ist also das regul\"are $p$-Eck konstruierbar, dann ist $p-1$ eine
Potenz von $2$. Daher ist z.B.\ das regul\"are $7$-Eck nicht konstruierbar.

Gau\ss\ produzierte als Jugendlicher eine kleine Sensation, als er zeigte,
dass man ein regul\"ares $17$-Eck mit Zirkel und Lineal konstruieren kann. Auf
diese Konstruktionsaufgaben kommen wir zur\"uck, wenn wir die K\"orpertheorie
schon weiter entwickelt haben.
\subsection{Algebraische K\"orpererweiterungen}
In diesem Abschnitt wollen wir algebraische Erweiterungen genauer untersuchen.
Wir haben gesehen, dass endliche Erweiterungen algebraisch sind. Die Umkehrung
muss nicht stimmen (man nehme z.B.\ die Vereinigung der K\"orper
$\QQQ(2^{1/2^n})$, $n=0,1,2,\dots$, \"uber $\QQQ$.) In einem wichtigen Fall
gilt allerdings die Umkehrung.
\begin{Lemma}
  Sei $L|K$ eine Erweiterung, und $\alpha_1,\alpha_2,\dots,\alpha_n\in L$ alle
  algebraisch \"uber $K$. Dann ist
  $K(\alpha_1,\alpha_2,\dots,\alpha_n)=K[\alpha_1,\alpha_2,\dots,\alpha_n]$
  eine endliche Erweiterung von $K$.
\end{Lemma}
\begin{proof}[Beweis]
  Ist $\alpha\in L$ algebraisch \"uber $K$, dann ist $\alpha$ erst recht
  algebraisch \"uber jedem Zwischenk\"orper von $L|K$. Die Behauptung folgt
  dann per Induktion und der bereits bewiesenen Version f\"ur $n=1$.
\end{proof}
Sind $\alpha,\beta\in L$ algebraisch \"uber $K$, dann liegen die Elemente
$\alpha+\beta$, $\alpha-\beta$, $\alpha\beta$, $\alpha/\beta$ (falls
$\beta\ne0$) alle in $K(\alpha,\beta)$, sind also algebraisch \"uber $K$.

Eine wichtige Transitivit\"atseigenschaft algebraischer Erweiterungen ist
\begin{Satz}\label{S:algtran}
  Die Erweiterungen $F|E$ und $E|K$ seien algebraisch. Dann ist auch $E|K$
  algebraisch.
\end{Satz}
\begin{proof}[Beweis]
  Sei $\alpha\in F$, und $f(X)=a_nX^n+\dots+a_1X+a_0\in E[X]$ das
  Minimalpolynom von $\alpha$ \"uber $E$. Setze $L=K(a_n,\dots,a_1,a_0)$. Nach
  dem obigen Lemma ist $[L:K]<\infty$. Ferner gilt $[L(\alpha):L]\le n$, also
  auch $[L(\alpha):K]<\infty$, und damit ist $\alpha$ algebraisch \"uber $K$.
\end{proof}
Wie wir oben schon gesehen haben k\"onnen wir die komplexen Zahlen auffassen
als den Restklassenring $\RRR[X]/(X^2+1)$. Bezeichnet $\bar X$ das Bild von
$X$ in diesem Restklassenring, dann gilt $\bar
X^2=\overline{X^2}+\overline{1}-\overline{1}=
\overline{X^2+1}-\overline{1}=-\overline{1}$.  Das rein algebraisch
konstruierte Element $\bar X$ entspricht also der imagin\"aren Einheit
$i\in\CCC$. Kronecker hat diese Idee verallgemeinert.
\begin{Satz}
  Sei $f\in K[X]$ ein irreduzibles Polynom vom Grad $n$, und $L=K[X]/(f)$ der
  Restklassenring. Identifiziert man $a\in K$ mit $a+(f)$ in $L$, so kann man
  $L$ als Oberk\"orper von $K$ auffassen. Setze $\alpha:=X+(f)\in L$. Dann gilt
  $f(\alpha)=0$, $L=K[\alpha]=K(\alpha)$ und $[L:K]=n$.
\end{Satz}
\begin{proof}[Beweis]
  Es gilt $f(\alpha)=f(X)+(f)=(f)$, da $f(X)\in(f)$. Da $(f)$ irreduzibel ist,
  ist $L=K[X]/(f)$ ein K\"orper, der offensichtlich von $\alpha$ erzeugt
  wird. Ferner ist $f$, wieder wegen der Irreduzibilit\"at von $f$, das
  Minimalpolynom von $\alpha$. Die Behauptungen folgen.
\end{proof}
Man nennt $L$ einen \ei{Wurzelk\"orper} von $f$. Ein wichtiger Satz ist
\begin{Satz}
  Seien $E$ und $F$ Erweiterungen eines K\"orpers $K$, und $f(X)\in K[X]$
  irreduzibel. Seien $\alpha\in E$ und $\beta\in F$ mit
  $f(\alpha)=f(\beta)=0$. Dann gibt es einen $K$-Isomorphismus $K(\alpha)\to
  K(\beta)$ mit $\phi(\alpha)=\beta$.
\end{Satz}
\begin{proof}[Beweis]
  Die $K$-Epimorphismen $X[X]\to K(\alpha)$, $X\mapsto\alpha$ und $X[X]\to
  K(\alpha)$, $X\mapsto\beta$ haben beide den Kern $(f)$. Die Behauptung folgt
  aus dem Homomorphiesatz.
\end{proof}
Ein K\"orper $K$ hei\ss t \ei{algebraisch abgeschlossen}, wenn jedes nicht
konstante Polynom $f(X)\in K[X]$ \"uber $K$ in Linearfaktoren zerf\"allt.
Wegen der M\"oglichkeit Linearfaktoren abzudividieren ist das nat\"urlich
\"aquivalent dazu, dass jedes nicht konstante Polynom $f(X)\in K[X]$ eine
Nullstelle in $K$ hat.
\subsubsection{Algebraischer Abschluss}
Ist $K$ ein K\"orper, dann ist ein \ei{algebraischer Abschluss} von $K$ eine
algebraische Erweiterung $\tilde K$, die algebraisch abgeschlossen ist. So ist
z.B.\ $\CCC$ ein algebraischer Abschluss von $\RRR$, aber nicht von $\QQQ$. Im
folgenden wollen wir sehen, dass jeder K\"orper $K$ einen algebraischen
Abschluss $\tilde K$ hat. Wir ben\"otigen zwei vorbereitende Lemmata.
\begin{Lemma}
  Sei $K$ ein K\"orper. Dann gibt es eine Erweiterung $E$ von $K$, so dass
  jedes nicht konstante Polynom $f(X)\in K[X]$ eine Nullstelle in $E$ hat.
\end{Lemma}
\begin{proof}[Beweis]
  F\"ur jedes nicht konstante Polynom $f\in K[X]$ sei $X_f$ eine Variable
  \"uber $K$. Sei $S$ die Menge $X_f$, und $K[S]$ der Polynomring \"uber $K$
  in den Variablen $S$. Sei $I$ das von allen $f(X_f)$ erzeugte Ideal von
  $K[S]$. Wir behaupten $I\ne K[S]$, und nehmen dazu an, dass Gleichheit gilt.
  Dann gibt es $g_1,g_2,\dots,g_n\in K[S]$ mit
\[
g_1f_1(X_{f_1})+g_2f_2(X_{f_2})+\dots+g_nf_n(X_{f_n})=1.
\]
Sei nun $K(\alpha_1)$ ein Wurzelk\"orper von $f_1$, also
$f_1(\alpha_1)=0$. Nun sei $K(\alpha_1,\alpha_2)$ ein Wurzelk\"orper von $f_2$
\"uber $K(\alpha_1)$. Setzt man das induktiv fort, so erh\"alt man eine
Erweiterung $L$ von $K$ mit $\alpha_i\in L$ und $f_i(\alpha_i)=0$. Ersetzt man
nun in obiger Gleichung $X_{f_i}$ durch $\alpha_i$, so folgt der Widerspruch
$0=1$.

Daher ist $I<K[S]$, und es gibt daher ein maximales Ideal $M$ von $K[S]$,
welches $I$ enth\"alt. Betrachte den $K$-Homomorphismus $\phi:K[S]\to K[S]/M$,
und beachte, dass $E:=K[S]/M$ eine K\"orpererweiterung von $K$ ist. Wegen
$f(X_f)\in I\subseteq M$ gilt $0=\phi(f(X_f))=f(\phi(X_f))$, d.h., $f$ hat
eine Nullstelle in $E$.
\end{proof}
\begin{Lemma}
  Sei $K$ ein K\"orper. Dann gibt es einen algebraisch abgeschlossenen
  K\"orper $L\supseteq K$.
\end{Lemma}
\begin{proof}[Beweis]
  Mit Hilfe des obigen Lemmas konstruieren wir eine Kette $K=E_0\subseteq
  E_1\subseteq E_2\subseteq\dots$ von K\"orpern, so dass jedes Polynom aus
  $E_i[X]$ eine Nullstelle in $E_{i+1}$ hat. Sei $L$ die Vereinigung der
  $E_i$. Offenbar ist $L$ ein K\"orper, \"uber dem jedes Polynom eine
  Nullstelle hat. Durch abdividieren von Linearfaktoren folgt daraus, dass
  jedes Polynom aus $L[X]$ sogar in Linearfaktoren zerf\"allt.
\end{proof}
\begin{Satz}
  Jeder K\"orper hat einen algebraischen Abschluss.
\end{Satz}
\begin{proof}[Beweis]
  Sei $K$ ein beliebiger K\"orper, und $L\supseteq K$ algebraisch
  abgeschlossen. Sei $\tilde K$ die Menge der Elemente aus $L$, die
  algebraisch sind \"uber $K$. Wir wissen schon, dass $\tilde K$ ein K\"orper
  ist. Es bleibt zu zeigen, dass auch $\tilde K$ algebraisch abgeschlossen ist.
  Sei $f\in\tilde K[X]$ ein nicht-konstantes Polynom, und $\alpha\in L$ mit
  $f(\alpha)=0$. Damit ist $\alpha$ algebraisch \"uber $\tilde K$, und damit
  auch algebraisch \"uber $K$ wegen Satz \ref{S:algtran}, also
  $\alpha\in\tilde K$. Die Behauptung folgt.
\end{proof}
\begin{Bemerkung*}
Da die irreduziblen Polynome \"uber einem algebraisch abgeschlossenen K\"orper
den Grad $1$ haben, besitzen algebraisch abgeschlossene K\"orper keine echten
algebrische Erweiterungen.
\end{Bemerkung*}
\subsubsection{Fortsetzung von Homomorphismen}
In unserer weiteren Entwicklung spielen Automorphismen von K\"orpern eine
wichtige Rolle. Wir beginnen mit der Frage der Fortsetzbarkeit von
Homomorphismen. Ist $\sigma:K\to L$ ein Homomorphismus und $f\in K[X]$, dann
bezeichnet $f^\sigma$ das Polynom, das durch Anwenden von $\sigma$ auf die
Koeffizienten von $f$ entsteht.
\begin{Lemma}\label{L:injbij}
  Sei $E|K$ eine algebraische Erweiterung, und $\sigma:E\to E$ ein
  $K$-Homomorphismus. Dann ist $\sigma$ ein Automorphismus von $E$.
\end{Lemma}
\begin{proof}[Beweis]
Es ist nur die Surjektivit\"at von $\sigma$ zu zeigen. Sei $\alpha\in E$, und
$f$ das Minimalpolynom von $f$ \"uber $K$. Sei $\beta\in E$ eine Nullstelle
von $f$. Da $\sigma$ die Identit\"at auf $K$
ist, folgt
\[
0=f(\beta)^\sigma=f(\beta^\sigma).
\]
Damit ist $\beta^\sigma\in E$ ebenfalls eine Nullstelle von $f$. Da $f$ nur
endlich viele Nullstellen in $E$ hat, werden diese von $\sigma$ permutiert.
Insbesondere wird eine dieser Nullstellen von $\sigma$ auf $\alpha$
abgebildet.
\end{proof}
\begin{Lemma}\label{L:Fortssep}
  Sei $\sigma:K\to L$ ein K\"orperhomomorphismus. Sei $\alpha$ in einem
  Oberk\"orper von $K$, und algebraisch \"uber $K$. Sei $f(X)\in K[X]$ das
  Minimalpolynom von $\alpha$ \"uber $K$. Dann ist die Anzahl der
  Fortsetzungen von $\sigma$ zu einem Homomorphismus $\tau:K(\alpha)\to L$
  gleich der Anzahl der verschiedenen Nullstellen, die $f^\sigma$ in $L$ hat.
\end{Lemma}
\begin{proof}[Beweis]
  Sei $\tau:K(\alpha)\to L$ eine Fortsetzung von $\sigma$.  Wegen
\[
0=f(\alpha)^\tau=f^\sigma(\alpha^\tau)
\]
ist $\alpha^\tau$ eine Nullstelle von $f^\sigma(X)$. Da $\tau$ durch
$\alpha^\tau$ eindeutig gegeben ist, folgt, dass es f\"ur $\tau$ h\"ochstens
so viele M\"oglichkeiten gibt wie $f^\sigma$ verschiedene Nullstellen in $L$
hat.

Sei nun umgekehrt $\beta\in L$ eine Nullstelle von $f^\sigma$. Jedes Element
aus $K(\alpha)$ hat die Form $g(\alpha)$ f\"ur ein Polynom $g\in K[X]$. Wir
definieren nun $\tau$ durch $g(\alpha)^\tau:=g^\sigma(\beta)$. Zun\"achst
m\"ussen wir sehen, dass diese Definition wohldefiniert ist. Dazu sei $h\in
K[X]$ ein weiteres Polynom mit $g(\alpha)=h(\alpha)$. Aber dann ist $\alpha$
eine Nullstelle von $g-h$, und somit $f$ ein Teiler von $g-h$. Es folgt
$g(X)-h(X)=f(X)q(X)$ f\"ur ein Polynom $q$, also
$g^\sigma(X)-h^\sigma(X)=f^\sigma(X)q^\sigma(X)$. Ersetze $X$ durch $\beta$.
Es folgt $g^\sigma(\beta)=h^\sigma(\beta)$, was die Wohldefiniertheit zeigt.
Offenbar ist $\tau$ additiv und multiplikativ, die Behauptung folgt.
\end{proof}
\begin{Satz}\label{S:Forts}
  Sei $E|K$ eine algebraische K\"orpererweiterung, und $\sigma:K\to L$ ein
  Homomorphismus in einen algebraisch abgeschlossenen K\"orper $L$. Dann
  l\"asst sich $\sigma$ fortsetzen zu einem Homomorphismus $E\to L$.
\end{Satz}
\begin{proof}[Beweis]
  Der Beweis benutzt das Zornsche Lemma. Sei $M$ die Menge der Paare
  $(F,\tau)$, wo $F$ ein Zwischenk\"orper von $E|K$ ist, und $\tau:F\to L$
  eine Fortsetzung von $\sigma$ ist. Die Menge $M$ ist partiell geordnet, mit
  $(F,\tau)\le (F',\tau')$ genau dann, wenn $F\subseteq F'$ und $\tau$ die
  Einschr\"ankung von $\tau'$ auf $F$ ist. Die Ordnung ist induktiv, da die
  Vereinigung einer total geordneten Teilmenge von $M$ wieder zu $M$ geh\"ort.
  Wegen des Lemmas von Zorn besitzt $M$ daher ein maximales Element
  $(F,\tau)$. W\"are $F\subsetneq E$, dann g\"abe es ein $\alpha\in E\setminus
  F$, und nach dem vorigen Lemma eine Fortsetzung des Homomorphismus
  $\tau:F\to L$ zu $F(\tau)\to L$, im Widerspruch zur Maximalit\"at von
  $(F,\tau)$.
\end{proof}
Wegen des folgenden Korollars spricht man oft (etwas unpr\"azise) von ``dem
algebraischen Abschluss'' eines K\"orpers.
\begin{Korollar}
Seien $E$ und $F$ zwei algebraische Abschl\"usse eines K\"orpers $K$. Dann
gibt es eine $K$-Isomorphie zwischen $E$ und $F$.  
\end{Korollar}
\begin{proof}[Beweis]
  Sei $\sigma:K\to F$ die Inklusionsabbildung. Wegen des obigen Satzes besitzt
  $\sigma$ eine Fortsetzung $\tau:E\to F$. Aber $F$ ist algebraisch \"uber
  $K$, und damit erst recht algebraisch \"uber $E^\tau$. Aber $E$ und damit
  $E^\tau$ sind algebraisch abgeschlossen, also $E^\tau=F$. Die Behauptung
  folgt.
\end{proof}
\subsubsection{Zerf\"allungsk\"orper}
Sei $f\in K[X]$ ein nicht konstantes Polynom, und
$f(X)=c(X-\alpha_1)\dots(X-\alpha_n)$ mit $c\in K$ und $\alpha_i$ Elemente
eines Oberk\"orpers von $K$. Der von den Nullstellen von $f$ \"uber $K$
erzeugte K\"orper $K(\alpha_1,\alpha_2,\dots,\alpha_n)$ hei\ss t der
\ei{Zerf\"allungsk\"orper} von $f$ \"uber $K$. Da die Polynomfaktorisierung
\"uber K\"orpern eindeutig ist, gibt es im algebraischen Abschluss $\tilde K$
von $K$ genau einen Zerf\"allungsk\"orper von $f$.

Wir sahen schon, dass Wurzelk\"orper irreduzibler Polynome isomorph sind, und
auch dass algebraische Abschl\"usse eines K\"orpers isomorph sind. Das gleiche
gilt auch f\"ur Zerf\"allungsk\"orper. Genauer gilt:
\begin{Satz}
  Seien $E$ und $F$ Zerf\"allungsk\"orper von $f\in K[X]$. Dann gibt es einen
  $K$-Isomorphismus zwischen $E$ und $F$. Liegt $E$ im algebraischen Abschluss
  $\tilde K$ von $K$, dann ist jeder $K$-Homomorphismus $F\to\tilde K$ eine
  Isomorphie zwischen $E$ und $F$.
\end{Satz}
\begin{proof}[Beweis]
Der erste Teil des Satzes folgt aus dem zweiten Teil. Daher nehmen wir
  $E\subseteq \tilde K$ an.
  
  Nach Satz \ref{S:Forts} gibt es einen $K$-Isomorphismus $\sigma:F\to\tilde
  K$. Sei
\[
f(X)=c(X-\beta_1)(X-\beta_2)\dots(X-\beta_n)\text{ mit $c\in K$, $\beta_i\in
  F$.}
\]
Es folgt
\[
f(X)=f^\sigma(X)=c(X-\beta_1^\sigma)(X-\beta_2^\sigma)\dots(X-\beta_n^\sigma).
\]
Die Elemente $\beta_i^\sigma$ erzeugen den in $\tilde K$ liegenden
Zerf\"allungsk\"orper von $f$, also
$E=K(\beta_1^\sigma,\beta_2^\sigma,\dots\beta_n^\sigma)=
K(\beta_1,\beta_2,\dots\beta_n)^\sigma=K$, und die Behauptung folgt.
\end{proof}
\subsubsection{Separabilit\"at}
In diesem Abschnitt sei $\tilde K$ ein algebraischer Abschluss von $K$. Sei
$\phi$ der eindeutige Ringhomomorphismus $\ZZZ\to K$ mit Kern $p\ZZZ$. Da
$\ZZZ/p\ZZZ$ ein Integrit\"atsbereich ist, ist $p$ eine Primzahl oder $p=0$.
Man nennt $p$ die \ei{Charakteristik} von $K$. Sei $p>0$. Dann gelten in $K$
die Identit\"aten $(a+b)^p=a^p+b^p$ und $a+a+\dots+a=0$ ($p$ Summanden). Diese
Beziehungen haben merkw\"urdige Konsequenzen. Sei z.B.\ $K=\FFF_p(t)$ der
rationale Funktionenk\"orper \"uber $\FFF_p$, und $f(X)=X^p-t$. Nach dem
Eisensteinkriterium ist $f$ irreduzibel \"uber $\FFF_p(t)$. Seien
$\alpha,\beta\in\tilde K$ zwei Nullstellen von $f$. Dann gilt
$0=\alpha^p-\beta^p=(\alpha-\beta)^p$, also $\alpha=\beta$. Damit hat $f$ in
$\tilde K$ genau eine Nullstelle, und diese mit der Vielfachheit $p$.

Man nennt ein nicht konstantes Polynom \ei{separabel}, wenn es keine
mehrfachen Nullstellen hat. Das obige Beispiel $X^p-t$ \"uber $\FFF_p(t)$ ist
also nicht separabel. Wir werden gleich sehen, dass die Existenz inseparabler
Polynome eine sehr spezielle Eigenschaft ist.

\begin{Satz}\label{S:inseppol}
  Sei $f(X)\in K[X]$ ein irreduzibles Polynom, das nicht separabel ist. Dann
  ist die Charakteristik $p$ von $K$ positiv. Ferner gibt es $e\in\NNN$ mit
  $f(X)=g(X^{p^e})$, wobei $g(X)\in K[X]$ separabel und irreduzibel ist.
\end{Satz}
\begin{proof}[Beweis]
  Sei $p$ die Charakteristik von $K$. Falls $p=0$ ist, dann setzen wir $g:=f$,
  und im Fall $p>0$ w\"ahlen wir $e\ge0$ maximal, so dass es ein $g(X)\in
  K[X]$ gibt mit $f(X)=g(X^{p^e})$.  Nat\"urlich ist $g(X)$ irreduzibel. Wir
  nehmen an, dass $g(X)$ nicht separabel ist. Dann haben $g(X)$ und die
  Ableitung $g'(X)$ eine gemeinsame Nullstelle in $\tilde K$. Insbesondere
  haben $g$ und $g'$ einen gemeinsamen Teiler in $K[X]$ von positivem Grad.
  Wegen der Irreduzibilit\"at von $g$ geht das nur, wenn $g$ ein Teiler von
  $g'$ ist. Wegen $\grad g'<\grad g$ folgt $g'=0$. Die Ableitung des Monoms
  $X^m$ ist $mX^{m-1}$. Diese verschwindet nur, wenn $p>0$ ein Teiler von $m$
  ist. Aus $g'=0$ folgt also $p>0$, und dass alle Exponenten von $g'$ durch
  $p$ teilbar sind. Daher hat $g(X)$ die Form $h(X^p)$, also
  $f(X)=h(X^{p^{e+1}})$, im Widerspruch zur Wahl von $e$.
\end{proof}
In der Situation des Satzes sei $g(X)=c(X-\alpha_1)\dots(X-\alpha_n)$ mit
$\alpha_i\in\tilde K$ paarweise verschieden. Die Nullstellen von $f$ sind
daher gerade die Nullstellen der Faktoren $X^{p^e}-\alpha_i$. Wie im obigen
Beispiel sieht man, dass jeder der Faktoren $X^{p^e}-\alpha_i$ genau eine
Nullstelle in $\tilde K$ hat. Daher hat $f$ genau $n=\grad g=\grad f/p^e$
verschiedene Nullstellen. Man nennt $n$ auch den \ei{Separabilit\"atsgrad} von
$f$. Sei $\beta$ eine Nullstelle von $f$.  Wir haben schon gesehen (Lemma
\ref{L:Fortssep}), dass $n$ gleich der Anzahl der $K$-Homomorphismen
$K(\beta)\to\tilde K$ ist. Das motiviert die folgende
\begin{Definition}
  Sei $[E:K]$ eine algebraische Erweiterung. Der \ei{Separabilit\"atsgrad}
  $[E:K]_\text{sep}$ ist die Anzahl der $K$-Homomorphismen $E\to\tilde K$.
\end{Definition}
Das folgende Lemma zeigt unter anderem, dass $[E:K]_\text{sep}$ nicht von der
Wahl von $\tilde K$ abh\"angt.
\begin{Lemma}
  Sei $E|K$ eine algebraische Erweiterung, $L$ ein algebraisch abgeschlossener
  K\"orper, und $\sigma:K\to L$ ein Homomorphismus. Dann hat $\sigma$ genau
  $[E:K]_\text{sep}$ Fortsetzungen $\tau:E\to L$.
\end{Lemma}
\begin{proof}[Beweis]
  Wir wissen schon (Satz \ref{S:Forts}), dass sich $\sigma:K\to L$ fortsetzt
  zu einem Homomorphismus $\rho:\tilde K\to L$. Wir konstruieren eine
  Bijektion zwischen den Fortsetzungen $\tau:E\to L$ von $\sigma$, und den
  $K$-Homomorphismen $\tau':E\to \tilde K$. Sei also $\tau:E\to L$ eine
  Fortsetzung von $\sigma$. Offenbar ist $E^\tau\subseteq L$ eine algebraische
  Erweiterung von $K^\tau=K^\sigma=K^\rho$. Da aber $\tilde K^\rho$
  algebraisch abgeschlossen ist, gilt $E^\tau\subseteq \tilde K^\rho$. Damit
  ist $\tau'=\tau\rho^{-1}$ ein $K$-Homomorphismus $E\to\tilde K$. Umgekehrt
  liefert jeder $K$-Homomorphismus $\tau':E\to\tilde K$ eine Fortsetzung
  $\tau=\tau'\rho:E\to L$ von $\sigma$.
\[
\begin{diagram}
  \node{\tilde K} \arrow[2]{e,t}{\rho} \node[2]{\tilde K^\rho\subseteq L}\\
\node[2]{E} \arrow{nw,t}{\tau'} \arrow{ne,t}{\tau}\\
\node{K} \arrow[2]{n,l}{\subseteq} \arrow{ne,r}{\subseteq}
\arrow[2]{e,b}{\sigma} \node[2]{K^\sigma} \arrow[2]{n,r}{\subseteq} 
\end{diagram}
\]
\end{proof}
Auch f\"ur den Separabilit\"atsgrad gilt eine Multiplikationsformel:
\begin{Satz}
  Seien $E|F$ und $F|K$ algebraische Erweiterungen. Dann gilt
\[
[E:K]_\text{sep}=[E:F]_\text{sep}[F:K]_\text{sep}.
\]
Falls $F|K$ endlich ist, dann ist $[F:K]_\text{sep}$ ein Teiler von $[F:K]$.
Ist der Quotient $>1$, dann hat $K$ eine positive Charakteristik $p$, und der
Quotient ist eine Potenz von $p$.
\end{Satz}
\begin{proof}[Beweis]
  Nach dem vorigen Lemma ist jeder $K$-Homomorphismus $E\to\tilde K$ eine der
  $[E:F]_\text{sep}$ Fortsetzungen der $[F:K]_\text{sep}$ m\"oglichen
  $K$-Homomorphismen $F\to\tilde K$. Daraus folgt der erste Teil der
  Behauptung.

Nun sei $F|K$ endlich. Dann gibt es Elemente $\alpha_1,\dots,\alpha_n\in F$
und eine Kette $K=K_0\subseteq K_1\subseteq\dots\subseteq K_n=F$ mit
$K_i=K_{i-1}(\alpha_i)$. Sei $n_i$ bzw.\ $s_i$ der Grad bzw.\ der
Separabilit\"atsgrad des Minimalpolynoms von $\alpha_i$ \"uber $K_{i-1}$. Dann
gilt $n_i=[K_i:K_{i-1}]$, und es gilt entweder $n_i=s_i$ f\"ur alle $i$, oder
$K$ hat Charakteristik $p>0$ und $n_i/s_i$ ist eine $p$-Potenz nach Satz
\ref{S:inseppol}. Die Behauptung folgt nun aus dem ersten Teil.
\end{proof}
Man nennt eine endliche Erweiterung $F|K$ \ei{separabel}, wenn
$[F:K]=[F:K]_\text{sep}$ gilt. Ein \"uber $K$ algebraisches Element $\alpha$
hei\ss t \ei{separabel}, wenn $K(\alpha)|K$ separabel ist. 

Ist $\alpha$ separabel \"uber $K$, dann ist $\alpha$ auch separabel \"uber
jedem Oberk\"orper $E$ von $K$, da das Minimalpolynom von $\alpha$ \"uber $E$
ein Teiler des Minimalpolynoms \"uber $K$ ist.

Hieraus und aus obigem Satz folgt schnell, dass eine endliche Erweiterung
$F|K$ genau dann separabel ist, wenn alle $\alpha\in F$ separabel \"uber $K$
sind, und das wiederum ist \"aquivalent dazu, dass $F$ von $K$ und endlich
vielen separablen Elementen erzeugt wird (\"Ubung!).

Sind $E$ und $F$ Teilk\"orper von $L$, dann nennt man den von $E$ und $F$
erzeugten Teilk\"orper von $L$ das \ei{Kompositum} von $E$ und $F$. Man
schreibt daf\"ur $EF$. Offenbar ist $EF$ der kleinste Teilk\"orper von $L$,
der $E$ und $F$ enth\"alt.

Die Separabilit\"at hat gute Vererbungseigenschaften:
\begin{Satz}
  Seien $E$ und $F$ Zwischenk\"orper von $L|K$. Dann gilt:
  \begin{itemize}
  \item[(a)] Seien $E|K$ und $L|E$ endlich. Dann ist $E|K$ genau dann
  separabel, wenn $E|K$ und $L|E$ separabel sind.
\item[(b)] $E|K$ sei endlich und separabel. Dann ist auch $EF|F$ separabel.
\item[(c)] $E|K$ und $F|K$ seien endlich und separabel. Dann ist auch $EF|K$
  endlich und separabel.
  \end{itemize}
\end{Satz}
\begin{proof}[Beweis]
  (a) folgt aus der Multiplikativit\"at des Grads bzw.\ des
  Separabilit\"atsgrads.

Sei $E|K$ endlich und separabel. Dann ist $E|K$ von endlich vielen separablen
Elementen erzeugt. Diese Elemente sind separabel \"uber $F$, und erzeugen mit
$F$ gerade $EF$. Hieraus folgt (b).

In (c) sind $E|K$ und $F|K$ jeweils von endlich vielen separablen Elementen
erzeugt. Daher ist $EF$ \"uber $K$ ebenfalls von endlich vielen separablen
Elementen erzeugt, und die Behauptung folgt.
\end{proof}
\subsubsection{Der Satz vom primitiven Element}
Eine f\"ur die weitere Entwicklung unentbehrliche Aussage, die auf Abel
zur\"uckgeht, ist der \ei{Satz vom primitiven Element}.
\begin{Satz}
  Sei $E|K$ eine endliche separable Erweiterung. Dann gibt es $\alpha\in E$
  mit $E=K(\alpha)$. Man nennt $\alpha$ ein \ei{primitives Element} der
  Erweiterung $E|K$.
\end{Satz}
\begin{proof}[Beweis]
Ist $K$ ein endlicher K\"orper, dann ist auch $E$ ein endlicher K\"orper, und
  hat daher eine zyklische multiplikative Gruppe. Ein Erzeuger dieser Gruppe
  ist nat\"urlich auch ein Erzeuger von $E|K$. Daher sei im folgenden
  $\abs{K}=\infty$.

  Per Induktion k\"onnen wir annehmen, dass $E=K(\alpha,\beta)$ gilt. Sei
  $n=[E:K]=[E:K]_\text{sep}$, und sei $S$ die Menge der $n$ verschiedenen
  $K$-Homomorphismen $E\to\tilde K$. Wir betrachten das Polynom
\[
F(X):=\prod_{\sigma\in S}\prod_{\tau\in S\setminus\{\sigma\}}
((\beta^\sigma-\beta^\tau)X+(\alpha^\sigma-\alpha^\tau)).
\]
Man beachte, dass $F(X)$ nicht das $0$-Polynom ist. Denn dann w\"aren $\sigma$
und $\tau$ die gleiche Abbildung auf $K(\alpha,\beta)$.

Da $K$ ein unendlicher K\"orper ist, gibt es ein $c\in K$ mit $F(c)\ne0$, also
\[
(\alpha+c\beta)^\sigma\ne(\alpha+c\beta)^\tau.
\]
Daher hat $K(\alpha+c\beta)$ mindestens $n$ verschiedene $K$-Homomorphismen
nach $\tilde K$, also
\[
[K(\alpha,\beta):K]=n\le
[K(\alpha+c\beta):K]_\text{sep}\le[K(\alpha,\beta):K],
\]
also $K(\alpha+c\beta)=K(\alpha,\beta)$.
\end{proof}
\subsubsection{Normale K\"orpererweiterungen}
Zum Studium der Automorphismen von K\"orpern und Zwischenk\"orpern von
K\"orpererweiterungen ist der Begriff einer normalen Erweiterung
zentral. Hierzu verallgemeinern wir den Begriff eines Zerf\"allungsk\"orpers
etwas: Sei $f_i$, $i\in I$ eine Familie von Polynomen $f_i(X)\in K[X]$. Dann
ist die Erweiterung $E|K$ ein Zerf\"allungsk\"orper der $f_i$, wenn jedes
Polynom $f_i$ \"uber $E$ in Linearfaktoren zerf\"allt, und $E$ von $K$ und den
Nullstellen der $f_i$ erzeugt wird.
\begin{Satz}
  Sei $\tilde K$ ein algebraischer Abschluss von $K$, und $F$ ein
  Zwischenk\"orper von $\tilde K|K$. Dann sind \"aquivalent:
  \begin{itemize}
  \item[(i)] Jeder $K$-Homomorphismus $F\to\tilde K$ ist ein Automorphismus
  von $F$.
\item[(ii)] Jedes irreduzible Polynom aus $K[X]$ mit einer Nullstelle in $F$
  zerf\"allt \"uber $F$ in Linearfaktoren.
\item[(iii)] $F$ ist der Zerf\"allungsk\"orper einer Familie von Polynomen
  \"uber $K$.
  \end{itemize}
\end{Satz}
\begin{proof}[Beweis]
  Es gelte (i). Sei $f(x)\in K[X]$ irreduzibel und $\alpha\in F$ mit
  $f(\alpha)=0$. Sei $\beta\in\tilde K$ eine beliebige Nullstelle von $f$.
  Dann gibt es einen $K$-Homomorphismus $K(\alpha)\to K(\beta)$ mit
  $\alpha\mapsto\beta$. Diesen erweitern wir zu einem Homomorphismus
  $\sigma:F\to\tilde K$. Nach Voraussetzung ist $\sigma$ ein Automorphismus
  von $F$, also $\beta\in F$. Damit zerf\"allt $f$ \"uber $F$ in
  Linearfaktoren. Wir erhalten (ii).
  
  Nun gelte (ii). Dann ist $F$ der Zerf\"allungsk\"orper aller Minimalpolynome
  der Elemente aus $F$, und damit gilt (iii).
  
  Schlie\ss lich gelte (iii). Sei $F$ der Zerf\"allungsk\"orper der $f_i\in
  K[X]$. Sei $\alpha\in F$ die Nullstelle eines der Polynome $f_i$. F\"ur
  jeden $K$-Homomorphismus $\sigma:F\to\tilde K$ ist $\alpha^\sigma$ eine
  Nullstelle von $f_i$. Da aber $f_i$ \"uber $F$ in Linearfaktoren zerf\"allt,
  gilt $\alpha^\sigma\in F$. Daher bildet $\sigma$ alle Nullstellen der $f_i$
  auf Elemente aus $F$ ab, also $F^\sigma\subseteq F$. Die Behauptung folgt
  nun aus Lemma \ref{L:injbij}.
\end{proof}
Eine Erweiterung $F|K$ hei\ss t \ei{normal}, wenn eine der Bedingungen (i),
(ii) oder (iii) erf\"ullt ist. Bei dem Begriff der Normalit\"at ist Vorsicht
angebracht, denn aus $E|F$ und $F|K$ normal folgt nicht, dass $E|K$ normal
ist. So sind z.B.\ $\QQQ(\sqrt[4]{2})|\QQQ(\sqrt{2})$ und
$\QQQ(\sqrt{2})|\QQQ$ normal, aber $\QQQ(\sqrt[4]{2})|\QQQ$ ist nicht normal,
da $\QQQ(\sqrt[4]{2})\subset\RRR$, aber die Nullstellen von $X^4-2$ sind nicht
alle reell.

Allerdings gilt eine andere wichtige Eigenschaft normaler Erweiterungen.
\begin{Satz}
  Seien $F$ und $E$ Erweiterungen von $K$, und enthalten in einem gemeinsamen
  K\"orper. Dann gilt
  \begin{itemize}
  \item[(a)] Ist $F|K$ normal, dann ist auch $FE|E$ normal.
\item[(b)] Sind $E|K$ und $F|K$ normal, dann sind auch $EF|K$ und $E\cap F|K$
  normal.
  \end{itemize}
\end{Satz}
\begin{proof}[Beweis]
  Zu (a): $F$ sei der Zerf\"allungsk\"orper der Polynome $f_i$ \"uber $K$.
  Dann ist $FE$ der Zerf\"allungsk\"orper der selben Polynome \"uber $E$.

Zu (b): Seien $E$ und $F$ die Zerf\"allungsk\"orper der Polynome $f_i$ bzw.\
$g_i$. Dann ist $EF$ der Zerf\"allungsk\"orper von $\{f_i\}\cup\{g_i\}$, und
der erste Teil folgt. Sei $f(X)\in K[X]$ irreduzibel, und mit einer Nullstelle
in $E\cap F$. Dann zerf\"allt $f(X)$ sowohl \"uber $E$, alsauch \"uber $F$ in
Linearfaktoren. Wegen der eindeutigen Faktorisierung liegen die Nullstellen
also in $E\cap F$, und die Behauptung folgt.
\end{proof}
\subsection{Galoistheorie}
\subsubsection{Der Hauptsatz der Galoistheorie}
In diesem Abschnitt erreichen wir das wichtigste Ziel der K\"orpertheorie, die
Galoistheorie. Sie verbindet die Gruppentheorie mit der K\"orpertheorie. Eine
algebraische Erweiterung $E|K$ hei\ss t \ei{galoissch}, wenn sie separabel und
normal ist. In diesem Fall nennt man die $K$-Automorphismen von $E$ die
\ei{Galoisgruppe} von $E|K$, und schreibt daf\"ur $\Gal(E|K)$. Ist $F$ ein
Zwischenk\"orper von $E|K$, dann ist auch $E|F$ normal und separabel, also
auch galoissch. Offenbar ist $\Gal(E|F)$ eine Untergruppe von $\Gal(E|K)$. Ist
$U$ eine Untergruppe von $G$, dann bezeichnet $\Fix(U)$ die Menge der von $U$
fixierten Elemente aus $E$. Offenbar ist $\Fix(U)$ ein Zwischenk\"orper von
$E|K$.

Der Hauptsatz der Galoistheorie besagt nun, dass im Falle einer endlichen
Galoiserweiterung $E|K$ die Abbildung $U\mapsto\Fix(U)$ eine Bijektion der
Menge der Untergruppen mit der Menge der Zwischenk\"orper liefert, und dass
$F\mapsto\Gal(E|F)$ die Umkehrabbildung dazu ist. Diesen Satz, und einige
weitere Kleinigkeiten, werden wir nun in einer Folge von Hilfss\"atzen
beweisen.

In diesem Abschnitt bezeichnet $\tilde K$ stets einen algebraischen Abschluss
von $K$.
\begin{Lemma}
  Sei $E|K$ eine Galoiserweiterung mit $G=\Gal(E|K)$. Dann gilt
  $K=\Fix(G)$. Die Abbildung $F\mapsto\Gal(E|F)$ bildet verschiedene
  Zwischenk\"orper von $E|K$ auf verschiedene Untergruppen von $G$ ab.
\end{Lemma}
\begin{proof}[Beweis]
  Ohne Einschr\"ankung gelte $E\subseteq\tilde K$. Sei $\alpha\in\Fix(G)$. Man
  setze einen $K$-Homomorphismus $K(\alpha)\to\tilde K$ fort zu einem
  Homomorphismus $\sigma:E\to\tilde K$. Wegen der Normalit\"at von $E|K$ ist
  $\sigma$ ein $K$-Automorphismus von $E$, also $\sigma\in G$. Gem\"a\ss\ 
  Voraussetzung wird $\alpha$ von $\sigma$ festgehalten. Daher gilt
  $[K(\alpha):K]_\text{sep}=1$. Aber $\alpha$ ist separabel \"uber $K$, und es
  folgt $[K(\alpha):K]=1$, also $\alpha\in K$. Dies zeigt den ersten Teil der
  Behauptung.
  
  Sei nun $F$ ein Zwischenk\"orper von $E|K$, und $U=\Gal(E|F)$. Da $E|F$
  galoissch ist, folgt aus dem ersten Teil der Behauptung $F=\Fix(U)$, und die
  Behauptung folgt.
\end{proof}
\begin{Lemma}
  Seien $F$ und $F'$ Zwischenk\"orper der Galoiserweiterung $E|K$. Dann gilt
  \begin{itemize}
  \item[(a)] $\Gal(E|F)\cap\Gal(E|F')=\Gal(E|FF')$.
\item[(b)] $\Fix(\gen{\Gal(E|F),\Gal(E|F')})=F\cap F'$.
\item[(c)] Es gilt $F\subseteq F'$ genau dann, wenn $\Gal(E|F')\le\Gal(E|F)$.
  \end{itemize}
\end{Lemma}
\begin{proof}[Beweis]
  Die Aussagen folgen direkt aus den Definitionen und obigem Lemma.
\end{proof}
Ist $E|K$ eine endliche separable Erweiterung, dann gibt es nach dem Satz von
primitiven Element ein $\alpha\in E$ mit $E=K(\alpha)$. Sei $L\supseteq E$ ein
Zerf\"allungsk\"orper des Minimalpolynoms von $\alpha$. Dann ist $L|K$ eine
endliche Galoiserweiterung, und offensichtlich eine kleinste Galoiserweiterung
von $K$ die $E$ enth\"alt. Man nennt $L$ auch eine \ei{Galoish\"ulle} von
$E|K$. Da $\Gal(L|K)$ nur endlich viele Untergruppen hat, hat die Erweiterung
$E|K$ nur endlich viele Zwischenk\"orper.
\begin{Lemma}
  Sei $E|K$ eine algebraische separable Erweiterung, so dass jedes Element aus
  $E$ den Grad $\le n$ \"uber $K$ hat. Dann gilt $[E:K]\le n$.
\end{Lemma}
\begin{proof}
Das Element $\alpha\in E$ habe maximalen Grad $m\le n$. F\"ur $\beta\in E$ ist
$K(\alpha,\beta)|K$ endlich und separabel, also $K(\alpha,\beta)=K(\gamma)$
f\"ur ein geeignetes Element $\gamma\in E$. Es folgt
\[
m=[K(\alpha):K]\le [K(\alpha,\beta):K]=[K(\gamma):K]\le m,
\]
also $\beta\in K(\alpha)$ f\"ur alle $\beta\in E$ und daher $E=K(\alpha)$.
\end{proof}
Wir haben schon gesehen, dass verschiedene Zwischenk\"orper einer
Galoiserweiterung $E|K$ zu verschiedenen Untergruppen von $G=\Gal(E|K)$
f\"uhren. Im allgemeinen ist nicht jede Untergruppe von $G$ von der Form
$\Gal(E|F)$ f\"ur einen Zwischenk\"orper $F$. Der folgende Satz jedoch zeigt,
dass dies wenigstens f\"ur endliche Galoiserweiterungen gilt.
\begin{Satz}[E.\ Artin]
  Sei $E$ ein K\"orper, und $G$ eine endliche Gruppe von Automorphismen von
  $E$. Sei $K$ der Fixk\"orper von $G$. Dann ist $E|K$ eine Galoiserweiterung
  mit Galoisgruppe $G$, und es gilt $[E:K]=\abs{G}$.
\end{Satz}
\begin{proof}[Beweis]
  Sei $\alpha\in E$ und $\{\sigma_1,\sigma_2,\dots,\sigma_m\}\subseteq G$ eine
  maximale Menge von Elementen mit $\sigma_1=1$, so dass die Elemente
  $\alpha^{\sigma_i}$ paarweise verschieden sind. Das hat zur Folge, dass
  f\"ur $\tau\in G$ die Menge der $\alpha^{\sigma_i\tau}$ eine Permutation der
  $\alpha^{\sigma_i}$ ist. Wir bilden das Polynom
\[
f(X)=\prod_{i}(X-\alpha^{\sigma_i}).
\]
Offenbar ist $\alpha$ eine Nullstelle von $f$, und es gilt $f=f^\tau$. Damit
liegen alle Koeefizienten von $f$ im Fixk\"orper von $G$, also in $K$. Ferner
ist $f$ separabel vom Grad $\le\abs{G}$.  Da dies f\"ur alle $\alpha\in E$
gilt, folgt aus dem obigen Lemma $[E:K]\le\abs{G}$. Jedes dieser Polymome $f$
zerf\"allt \"uber $E$ in Linearfaktoren, daher ist $E|K$ nicht nur separabel,
sondern auch normal, also galoissch. Die Galoisgruppe $H$ von $E|K$ hat
Ordnung $\le [E:K]\le\abs{G}$, also $\abs{H}\le\abs{G}$. Aber andererseits
gilt $G\subseteq H$, und somit gilt Gleichheit \"uberall.
\end{proof}
Es sei $F$ ein Zwischenk\"orper der Galoiserweiterung $E|K$, und
$\sigma\in\Gal(E|K)$. Dann besteht $\Gal(E|F^\sigma)$ aus all denen
$\tau\in\Gal(E|K)$, f\"ur die $a^{\sigma\tau}=a^\sigma$ f\"ur alle $a\in F$
gilt. Das ist \"aquivalent zu $a^{\sigma\tau\sigma^{-1}}=a$ f\"ur alle $a$,
also \"aquivalent zu $\sigma\tau\sigma^{-1}\in\Gal(E|F)$. Daher gilt
\[
\Gal(E|F^\sigma)=\sigma^{-1}\Gal(E|F)\sigma=\Gal(E|F)^\sigma.
\]
Diese einfache Beobachtung werden wir im folgenden mehrfach ohne Kommentar
anwenden.

Wenn $F$ ein Zwischenk\"orper der Galoiserweiterung $E|K$ ist, dann muss $F|K$
nicht automatisch galoissch sein. Wir geben nun ein gruppentheoretisches
Kriterium daf\"ur an, wann das der Fall ist.
\begin{Lemma}
  Sei $F$ ein Zwischenk\"orper der Galoiserweiterung $E|K$. Dann ist $F|K$
  genau dann galoissch, wenn $N=\Gal(E|F)$ ein Normalteiler von $G=\Gal(E|K)$
  ist. Ist das der Fall, dann liefert die Einschr\"ankung $G\to\Gal(F|K)$,
  $\sigma\mapsto \sigma|_F)$ einen Epimorphismus mit Kern $N$, also
  $\Gal(F|K)\cong G/N$.
\end{Lemma}

\begin{proof}[Beweis]
  Sei $F|K$ nicht galoissch. Da $F|K$ separabel ist, ist $F|K$ nicht normal.
  Daher gibt es ein $\alpha\in F$, dessen Minimalpolynom $f(X)$ \"uber $F$
  nicht in Linearfaktoren zerf\"allt. Wegen der Normalit\"at von $E|K$
  zerf\"allt aber $f$ \"uber $E$ in Linearfaktoren. Sei $\beta\in E$ eine
  Nullstelle von $f$ mit $\beta\notin F$. Dann gibt es $\sigma\in G$ mit
  $\alpha^\sigma=\beta$, also $F^\sigma\ne F$. Damit gilt
  $N\ne\Gal(E|F^\sigma)=N^\sigma$, d.h., $N$ ist nicht normal in $G$.
  
  Sei nun $F|K$ galoissch, mit $H=\Gal(F|K)$. Jedes Element von $H$ l\"asst
  sich zu einem Element aus $G$ fortsetzen. Daher ist die Einschr\"ankung
  $G\to H$, $\sigma\to\sigma|_F$ surjektiv. Der Kern dieser Abbildung ist
  $N=\Gal(E|F)$, und die Behauptung folgt aus dem Homomorphiesatz.
\end{proof}

Zusammenfassend k\"onnen wir nun unser Hauptergebnis formulieren.
\begin{Satz}[Hauptsatz der Galoistheorie]
  Sei $E|K$ eine endliche Galoiserweiterung. Dann gilt:
  \begin{itemize}
  \item[(a)] Die Abbildung $F\mapsto\Gal(E|F)$ liefert eine Bijektion zwischen
    der Menge der Zwischenk\"orper von $E|K$ und den Untergruppen von
    $\Gal(E|K)$. Die Umkehrabbildung dieser Bijektion ist $U\mapsto \Fix(U)$.
\item[(b)] F\"ur den Zwischenk\"orper $F$ gilt $[E:F]=\abs{\Gal(E|F)}$, und
  daher $[F:K]=[\Gal(E|K):\Gal(E|F)]$.
\item[(c)] Der Zwischenk\"orper $F$ ist genau dann normal \"uber $K$, wenn
  $\Gal(E|F)$ ein Normalteiler von $\Gal(E|K)$ ist. In diesem Fall gilt
  $\Gal(F|K)\cong\Gal(E|K)/\Gal(E|F)$.
  \end{itemize}
\end{Satz}
\begin{Bemerkung*}
  Sei $E|K$ eine galoissche Erweiterung, und $\tilde K\supseteq E$ ein
  algebraischer Abschluss von $K$. Da $E|K$ normal ist, ist jeder
  $K$-Homomorphismus $E\to\tilde K$ ein Automorphismus von $E$. Daher ist die
  von uns schon mehrfach betrachtete Menge der $K$-Homomorphismen $E\to \tilde
  K$ in diesem Fall nichts anderes als die Galoisgruppe $\Gal(E|K)$.
  
  Sei $f(X)\in K[X]$ ein irreduzibles Polynom, das \"uber $E$ in
  Linearfaktoren zerf\"allt. Dann wissen wir bereits, dass es f\"ur je zwei
  Nullstellen $\alpha,\beta\in E$ einen $K$-Homomomorphismus $E\to\tilde K$
  gibt mit $\alpha\mapsto\beta$. Daher operiert $\Gal(E|K)$ transitiv auf der
  Menge der Nullstellen von $f$. Diese einfache aber wichtige Beobachtung
  werden wir im folgenden mehrfach ben\"otigen.
  
  Unter der Galoisgruppe eines separablen (aber nicht notwendig irreduziblen)
  Polynoms $f(X)\in K[X]$ verstehen wir die Galoisgruppe $G=\Gal(L|K)$, wobei
  $L$ ein Zerf\"allungsk\"orper von $f$ \"uber $K$ ist. Da $L$ von den Wurzeln
  von $f$ erzeugt wird, operiert $G$ treu auf den Wurzeln von $f$. H\"aufig
  werden wir daher $G$ ohne Kommentar als Permutationsgruppe auf den Wurzeln
  von $f$ auffassen. Man schreibt dann meist $\Gal(f|K)$ f\"ur diese Gruppe.
\end{Bemerkung*}
\subsubsection{$\CCC$ ist algebraisch abgeschlossen}
Als erste einfache Anwendung der Galoistheorie wollen wir zeigen, dass der
K\"orper $\CCC$ der komplexen Zahlen algebraisch abgeschlossen ist. Der
funktionentheoretische Beweis dieser Aussage (z.B.\ als Korollar des Satzes
von Liouville) ist zwar einfacher und nat\"urlicher und erfordert weniger
Vorbereitungen, aber dennoch ist es im folgenden interessant zu sehen, wie die
Galoistheorie aus zwei einfachen Aussagen \"uber $\RRR$ und $\CCC$ ohne
weitere analytische Hilfsmittel diese nicht triviale Aussage liefert.
\begin{Lemma}\begin{itemize}
  \item[(a)] $\RRR$ besitzt keine echten Erweiterungen ungeraden Grades.
  \item[(b)] $\CCC$ besitzt keine Erweiterungen vom Grad $2$.
\end{itemize}
\end{Lemma}
\begin{proof}[Beweis]
  (a) Sei $E$ eine Erweiterung von $\RRR$ von ungeradem Grad $n$, und
  $\alpha\in E$ ein primitives Element, also $E=\RRR(\alpha)$. Das
  Minimalpolynom $f$ von $\alpha$ hat nach dem Zwischenwertsatz wegen
  $\lim_{x\to\infty}f(x)=\infty$, $\lim_{x\to-\infty}f(x)=-\infty$ eine reelle
  Nullstelle. Da $f$ \"uber $\RRR$ irreduzibel ist, folgt $n=1$.
  
  (b) Hierf\"ur ist zu zeigen, dass jedes Element aus $\CCC$ ein Quadrat in
  $\CCC$ ist. Sei $a+ib\in\CCC$ mit $a,b\in\RRR$, $i^2=-1$. Der Ansatz
  $z^2=(x+iy)^2=a+ib$ mit $x,y\in\RRR$ liefert $x^2-y^2=a$ und $2xy=b$. Wegen
  $(zi)^2=-z^2$ d\"urfen wir $a\ge0$ annehmen. Der Fall $b=0$ ist klar, also
  sei $b\ne0$. Aufl\"osen der zweiten Gleichung nach $y$ und Einsetzen in die
  erste liefert
\[
h(x)=x^2-\frac{b^2}{4x^2}=a.
\]
Wegen $b\ne0$ gilt $\lim_{x\to0}h(x)=-\infty$ und
$\lim_{x\to\infty}h(x)=\infty$. Nach dem Zwischenwertsatz gibt es also $0\ne
x\in\RRR$ mit $h(x)=a$. Setze $y=\frac{b}{2x}$, und es folgt $(x+iy)^2=a+ib$.
\end{proof}
Nach diesen Vorbereitungen folgt nun der
\begin{Satz}[Fundamentalsatz der Algebra]
  Der K\"orper $\CCC$ der komplexen Zahlen ist algebraisch abgeschlossen.
\end{Satz}
\begin{proof}[Beweis]
  Sei $E$ eine endliche Erweiterung von $\CCC$. Dann ist $E$ auch eine
  endliche separable Erweiterung von $\RRR$. Sei $L$ eine Galoish\"ulle von
  $E|\RRR$, und $G=\Gal(L|\RRR)$. Sei $P$ eine $2$-Sylowgruppe von $G$. Dann
  ist $[G:P]=[\Fix(P):\RRR]$ ungerade. Teil (a) des Lemmas liefert $[G:P]=1$,
  d.h., $G$ ist eine $2$-Gruppe. Damit ist auch $H=\Gal(L|\CCC)$ eine
  $2$-Gruppe. Falls $H=1$, dann gilt $L=\CCC$, also $E=\CCC$, und wir sind
  fertig. Sei also $\abs{H}\ge 2$. Nach Lemma \ref{L:p-GruppeIndp} hat $H$
  eine Untergruppe $N$ vom Index $2$. Dann gilt $[\Fix(N):\CCC]=2$, im
  Widerspruch zu (b).
\end{proof}
\subsubsection{Spuren und Normen}
Sei $E|K$ eine endliche separable Erweiterung, und $\tilde K$ ein
algebraischer Abschluss von $K$. Sei $S$ die Menge der $K$-Homomorphismen
$E\to\tilde K$. Wir definieren die \ei{Norm} $N_K^E(\alpha)$ von $\alpha\in E$
durch
\[
N_K^E(\alpha)=\prod_{\sigma\in S}\alpha^\sigma.
\]
Indem wir das Produkt durch eine Summe ersetzen, definieren wir die \ei{Spur}
$T_K^E(\alpha)$ durch
\[
T_K^E(\alpha)=\sum_{\sigma\in S}\alpha^\sigma.
\]
\begin{Satz}
  Sei $E|K$ eine endliche separable Erweiterung. Dann gilt
  \begin{itemize}
  \item[(a)] $N_K^E$ ist ein multiplikativer Homomorphismus $E^\star\to
  K^\star$.
\item[(b)] $T_K^E$ ist ein additiver Homomorphismus $E\to K$.
\item[(c)] Ist $F|E$ eine endliche separable Erweiterung, dann gilt
  $N_K^F=N_K^E\circ N_E^F$ und $T_K^F=T_K^E\circ T_E^F$.
\item[(d)] Ist $f(X)=X^n+a_{n-1}X^{n-1}+\dots+a_1X+a_0$ das Minimalpolynom von
  $\alpha\in E$, dann gilt $N_K^{K(\alpha)}(\alpha)=(-1)^na_0$ und
  $T_K^{K(\alpha)}(\alpha)=-a_{n-1}$.
  \end{itemize}
\end{Satz}
\begin{proof}[Beweis]
  Sei $L$ eine Galoish\"ulle von $E|K$, und $\alpha\in E$. Jedes $\sigma\in S$
  setzt sich fort zu einem Element aus $\Gal(L|K)$. Sei $\tau\in\Gal(L|K)$ und
  $\sigma\in S$. Wegen $E^\sigma\subseteq L$ ist $\sigma\tau$ definiert, und
  ein Element aus $S$. Daher ist $S\to S$, $\sigma\mapsto\sigma\tau$ eine
  Bijektion. Somit permutiert $\tau$ die Elemente $\alpha^\sigma$, $\sigma\in
  S$ untereinander, und es gilt $N_K^E(\alpha)^\tau=N_K^E(\alpha)$ und
  $T_K^E(\alpha)^\tau=T_K^E(\alpha)$ f\"ur alle $\tau\in\Gal(L|K)$. Daher
  liegen $N_K^E(\alpha)$ und $T_K^E(\alpha)$ im Fixk\"orper von $\Gal(L|K)$,
  also in $K$. Hieraus folgen (a) und (b), da die Homomorphieeigenschaften
  sowieso klar sind.
  
  Um (c) zu zeigen, nehmen wir an, dass $F\subseteq\tilde K$ gilt. F\"ur jedes
  $\sigma\in S$ w\"ahlen wir eine Fortsetzung $\hat\sigma:\tilde K\to\tilde
  K$. Sei $T$ die Menge der $E$-Homomorphismen $F\to\tilde K$. Zu jedem
  $K$-Homomorphismus $\rho:F\to\tilde K$ gibt es genau ein $\sigma\in S$, so
  dass $\rho|_E=\sigma$ gilt. Daher ist $\rho\hat\sigma^{-1}$ ein
  $E$-Homomorphismus $F\to\tilde K$, also $\rho\hat\sigma^{-1}=\tau$ f\"ur ein
  $\tau\in T$. Diese \"Uberlegung zeigt, dass die Menge der $\tau\hat\sigma$,
  $\sigma\in S$, $\tau\in T$ genau die Menge der $K$-Homomorphismen
  $F\to\tilde K$ durchl\"auft. Daraus folgt dann die Behauptung aus einer
  einfachen Rechnung und den Teilen (a) und (b).
  
  Sei $f(X)=(X-\alpha_1)\dots(X-\alpha_n)$ mit $\alpha\in\tilde K$. Die
  $\alpha_i$ sind gerade die verschiedenen Bilder von $\alpha$ unter den
  $K$-Homomorphismen $K(\alpha)\to\tilde K$. Die Behauptung folgt durch
  Betrachtung der Koeffizienten von $X^{n-1}$ und $X^0$.
\end{proof}
\subsubsection{Lineare Unabh\"angigkeit von Charakteren}
Sei $(G,\cdot)$ ein Monoid, und $K$ ein K\"orper. Ein \ei{Charakter} ist eine
Abbildung $\phi:G\to K^\star$ mit $\phi(ab)=\phi(a)\phi(b)$ und $\phi(1_G)=1$.
Sei $\phi_i$, $i\in I$ eine Menge von Charakteren. Man sagt, dass die $\phi_i$
\ei{linear unabh\"angig} sind, wenn folgendes gilt: Sei $\sum_i a_i\phi_i=0$
mit $a_i\in K$ und $a_i=0$ f\"ur alle bis auf endlich viele $i$. Dann folgt
schon $a_i=0$ f\"ur alle $i$.
\begin{Satz}[E.\ Artin]\label{S:Charlinunabh}
  Die Charaktere $G\to K^\star$ sind linear unabh\"angig.
\end{Satz}
\begin{proof}[Beweis]
  Sei
\[
a_1\phi_1+\dots+a_n\phi_n=0
\]
eine lineare Abh\"angigkeitsrelation mit $n$ minimal und alle $a_i\ne0$.
Nat\"urlich gilt $n\ge2$. Da $\phi_1$ und $\phi_2$ verschieden sind, gibt es
$g\in G$ mit $\phi_1(g)\ne\phi_2(g)$. F\"ur alle $x\in G$ gilt
\[
0=a_1\phi_1(gx)+\dots+a_n\phi_n(gx)=
a_1\phi_1(g)\phi_1(x)+\dots+a_n\phi_n(g)\phi_n(x)
\]
und
\[
0=\phi_1(g)(a_1\phi_1(x)+\dots+a_n\phi_n(x)).
\]
Subtraktion liefert
\[
b_2\phi_1+\dots+b_n\phi_n=0
\]
mit
\[
b_i=(\phi_i(g)-\phi_1(g))a_i.
\]
Wegen $b_2\ne0$ haben wir eine k\"urzere lineare Abh\"angigkeitsrelation, ein
Widerspruch.
\end{proof}
\subsubsection{Zyklische Erweiterungen}
Sei $n\in\NNN$. Unter einer \ei{primitiven $n$-ten Einheitswurzel} verstehen
wir ein Element $\alpha\in K$, das die multiplikative Ordnung $n$ hat. Die von
$\alpha$ erzeugte Gruppe hat Ordnung $n$, und besteht daher genau aus den
Nullstellen von $X^n-1$. Enth\"alt also $K$ eine primitive $n$-te
Einheitswurzel, dann zerf\"allt $X^n-1$ in Linearfaktoren, und $K$ enth\"alt
alle primitiven Einheitswurzeln aus $\tilde K$.

Eine K\"orpererweiterung $E|K$ hei\ss t zyklisch, abelsch oder aufl\"osbar,
wenn die Erweiterung $E|K$ galoissch ist, und $\Gal(E|K)$ zyklisch, abelsch
oder aufl\"osbar ist.

Der folgende Satz hatte in einer fundamentalen Arbeit von Hilbert die Nummer
90.
\begin{Satz}[Hilberts Satz 90]\label{S:90}
  Sei $E|K$ eine zyklische Erweiterung der Ordnung $n$ mit Galoisgruppe
  $G=\gen{\sigma}$. Sei $\beta\in E$. Dann gilt $N_K^E(\beta)=1$ genau dann,
  wenn es $0\ne\alpha\in E$ gibt mit $\beta=\frac{\alpha}{\alpha^\sigma}$.
\end{Satz}
\begin{proof}[Beweis]
  Setze $N=N_K^E$. Sei $\alpha\ne0$. Es gilt $N(\alpha)=\prod_{\tau\in
    G}\alpha^\tau$, insbesondere folgt
  $N(\alpha^\sigma)=N(\alpha)^\sigma=N(\alpha)\ne0$, und daher
  $N(\frac{\alpha}{\alpha^\sigma})=1$. Damit folgt die eine Richtung.

Sei nun $N(\beta)=1$. F\"ur $0\le i\le n-1$ sind die Abbildungen $x\mapsto
x^{\sigma^i}$ verschiedene Charaktere der additiven Gruppe von $E$. Daher ist
nach Satz \ref{S:Charlinunabh} die Abbildung
\[
x\mapsto \phi(x)=x+\beta x^\sigma+\beta\beta^\sigma x^{\sigma^2}+\dots
+\beta\beta^\sigma\dots\beta^{\sigma^{n-2}}x^{\sigma^{n-1}}
\]
nicht die $0$-Abbildung von $E$ nach $E$. Sei also $\gamma\in E$ mit
$\phi(\gamma)=\alpha\ne0$. Wir rechnen, unter Verwendung von $N(\beta)=1$ und
$\sigma^n=1$,
\[
\alpha-\beta\alpha^\sigma=
\gamma-\beta\beta^\sigma\dots\beta^{\sigma^{n-1}}\gamma^{\sigma^{n}}=
1-N(\beta)\gamma=0,
\]
und die Behauptung folgt.

\end{proof}
Wir hatten schon gesehen, dass eine quadratische Erweiterung $E|K$ f\"ur
K\"orper der Charakteristik $\ne0$ immer von der Form $E=K(\alpha)$ mit
$\alpha^2\in K$ ist. Der folgende Satz verallgemeinert das.
\begin{Satz}\label{S:zykErw}
  Sei $E|K$ eine zyklische Galoiserweiterung der Ordnung $n$. Der K\"orper $n$
  enthalte eine primitive $n$-te Einheitswurzel. Dann gibt es ein $\alpha\in
  E$ mit $E=K(\alpha)$ und $\alpha^n\in K$.
\end{Satz}
\begin{proof}[Beweis]
  Sei $\zeta\in K$ eine primitive $n$-te Einheitswurzel, und $\sigma$ ein
  Erzeuger von $\Gal(E|K)$. Es gilt $N(\zeta)=\zeta^n=1$, also auch
  $N(\zeta^{-1})=1$. Nach Satz \ref{S:90} gibt es $0\ne\alpha\in E$ mit
  $\alpha^\sigma=\alpha\zeta$. Iterierte Anwendung von $\sigma$ liefert
  $\alpha^{\sigma^i}=\zeta^i\alpha$. Da $\zeta$ die multiplikative Ordnung $n$
  hat, hat das Minimalpolynom von $\alpha$ \"uber $K$ mindestens die $n$
  verschiedenen Nullstellen $\zeta^i\alpha$ f\"ur $i=1,2,\dots,n$, also
  $[K(\alpha):K]\ge n$, und daher $E=K(\alpha)$ wegen $[E:K]=n$. Ferner ist
\[
(\alpha^n)^\sigma=(\alpha^\sigma)^n=(\zeta\alpha)^n=\alpha^n,
\]
also $\alpha^n\in\Fix(\Gal(E|K))$ und daher $\alpha^n\in K$.
\end{proof}
\subsubsection{Einheitswurzeln und Kreisteilungspolynome}
Im folgenden sei $n\in\NNN$, und $K$ ein K\"orper entweder der Charakteristik
$0$, oder teilerfremd zu $n$. Dann ist $X^n-1$ separabel. Sei $\tilde K$ ein
algebraischer Abschluss von $K$. Die Nullstellen von $X^n-1$ in $\tilde K$
bilden eine multiplikative Gruppe der Ordnung $n$. Da endliche Untergruppen
multiplikativer Gruppen zyklisch sind, ist die Gruppe zyklisch. Die Erzeuger
dieser Gruppe sind gerade die primitiven $n$-ten Einheitswurzeln.

Das $n$-te \ei{Kreisteilungspolynom} $\Phi_n(X)$ ist das Minimalpolynom einer
primitiven $n$-ten Einheitswurzel \"uber $\QQQ$. Der folgende Satz zeigt unter
anderem, dass $\Phi_n(X)$ unabh\"angig von der gew\"ahlten primitiven
Einheitswurzel ist.
\begin{Satz}
  Die Nullstellen von $\Phi_n(X)$ sind gerade die primitiven $n$-ten
  Einheitswurzeln aus $\CCC$. Ferner gilt $\Phi_n(X)\in\ZZZ[X]$.
\end{Satz}
\begin{proof}[Beweis]
  Offenbar ist $\Phi_n(X)$ ein Teiler von $X^n-1$. Nach dem Lemma von Gau\ss\
  hat jeder normierte Teiler von $X^n-1$ ganzzahlige Koeffizienten, daraus
  folgt der zweite Teil der Behauptung.
  
  Sei nun $X^n-1=\Phi_n(X)g(X)$ mit $g(X)\in\ZZZ[X]$. Da jede primitive $n$-te
  Einheitswurzel die Potenz (mit zu $n$ teilerfremdem Exponenten) einer fest
  gew\"ahlten primitiven $n$-ten Einheitswurzel ist, gen\"ugt es, das folgende
  zu zeigen: Sei $p$ eine Primzahl die $n$ nicht teilt, und $\zeta\in\CCC$ mit
  $\Phi_n(\zeta)=0$. Dann gilt auch $\Phi_n(\zeta^p)=0$.

Wir nehmen also an, dass $\Phi_n(\zeta^p)\ne0$ gilt. Dann folgt
$g(\zeta^p)=0$. Daher ist $\zeta$ eine Nullstelle von $g(X^p)$. Da aber
$\Phi_n(X)$ das Minimalpolynom von $\zeta$ ist, muss $\Phi_n(X)$ ein Teiler
von $g(X^p)$ sein, also
\[
g(X^p)=\Phi_n(X)h(X),
\]
mit $h(X)\in\ZZZ[X]$. F\"ur $u(X)\in\ZZZ[X]$ sei $\bar u(X)$ das
nat\"urliche Bild von $u(X)$ in $\FFF_p[X]$. Wegen der Identit\"at
$(U+V)^p=U^p+V^p$ in $\FFF_p[X]$ gilt $\bar g(X^p)=(\bar g(X))^p$, also
\[
\bar g(X)^p=\bar\Phi_n(X)\bar h(X).
\]
Daher sind die Polynome $\bar\Phi_n(X)$ und $\bar g(X)$ nicht
teilerfremd. Wegen
\[
X^n-1=\bar\Phi_n(X)\bar g(X)
\]
folgt daraus aber, dass $X^n-1$ \"uber $\FFF_p$ nicht separabel ist. Aber das
ist im Widerspruch dazu, dass $p$ kein Teiler von $n$ ist.
\end{proof}
\begin{Korollar}\label{K:zetaphi}
  Sei $\zeta\in\CCC$ eine primitive $n$-te Einheitswurzel. Dann ist
  $\QQQ(\zeta)|\QQQ$ galoissch, mit Gruppe isomorph zu $(\ZZZ/n\ZZZ)^\star$.
\end{Korollar}
\begin{proof}[Beweis]
  Jede primitive $n$-te Einheitswurzel ist von der Form $\zeta^k$ mit
  $k\in\ZZZ$ und $k$ teilerfremd zu $n$. Ferner ist $\zeta^k=\zeta^l$ genau
  dann, wenn $k-l\in n\ZZZ$. Daher ist f\"ur $a\in(\ZZZ/n\ZZZ)^\star$ die
  primitive Einheitswurzel $\zeta^a$ wohldefiniert. Die Galoisbilder von
  $\zeta$ sind wieder primitive $n$-te Einheitswurzeln, daher gibt es f\"ur
  alle $\sigma\in\Gal(\QQQ(\zeta)|\QQQ)$ ein $a_\sigma\in(\ZZZ/n\ZZZ)^\star$
  mit $\zeta^\sigma=\zeta^{a_\sigma}$. Wegen
\[
\zeta^{a_{\sigma\tau}}= \zeta^{\sigma\tau}=(\zeta^\sigma)^\tau=
(\zeta^{a_\sigma})^\tau= (\zeta^\tau)^{a_\sigma}=\zeta^{a_\tau a_\sigma}
\]
gilt $a_{\sigma\tau}=a_\sigma a_\tau$, d.h., $G\to (\ZZZ/n\ZZZ)^\star$,
$\sigma\mapsto a_\sigma$ ist ein Gruppenhomomorphismus. Da $\zeta$ den
K\"orper $\QQQ(\zeta)$ erzeugt, wird $\zeta$ nur von $1\in G$ fixiert. Der
Homomorphismus ist also injektiv. Wegen obigen Satzes kann $\zeta$ durch ein
geeignetes $\sigma\in G$ auf eine beliebige primitive $n$-te Einheitswurzel
abgebildet werden, daher ist der Homomorphismus aus surjektiv.
\end{proof}
\"Ahnlich zeigt man
\begin{Lemma}\label{L:Kz:K}
  Sei $n\in\NNN$, $K$ ein K\"orper $\zeta$ eine primitive $n$-te
  Einheitswurzel in einem Erweiterungsk\"orper von $K$. Dann ist $K(\zeta)|K$
  galoissch, mit einer Gruppe isomorph zu einer Untergruppe von
  $(\ZZZ/n\ZZZ)^\star$. Insbesondere ist $K(\zeta)|K$ eine abelsche
  K\"orpererweiterung.
\end{Lemma}
\subsubsection{Aufl\"osbarkeit durch Radikale}
Eine klassische und wichtige Frage ist die nach der Aufl\"osbarkeit von
Polynomen. Darunter versteht man die Berechnung der Nullstellen durch die vier
Grundrechenarten und das Ziehen von Wurzeln. Dies l\"asst sich
k\"orpertheoretisch pr\"azisieren. Der Einfachheit halber beschr\"anken wir
uns in diesem Kapitel auf K\"orper der Charakteristik $0$.
\begin{Definition*}
  Sei $K$ ein K\"orper der Charakteristik $0$, und $F$ eine Erweiterung von
  $K$. Dann hei\ss t $F$ \ei{aufl\"osbar durch Radikale}, wenn es eine Kette
\[
K=E_0\subseteq E_1\subseteq\dots\subseteq E_r\supseteq F
\]
von K\"orpern gibt, so dass f\"ur jedes $i=0,1,\dots r-1$ die Erweiterung
$E_{i+1}|E_i$ durch Adjunktion einer Wurzel entsteht, also
$E_{i+1}=E_i(\alpha_i)$ mit $\alpha_i^{m_i}\in E_i$ f\"ur ein $m_i\in\NNN$.
\end{Definition*}
Ist also $\alpha$ eine Nullstelle von $f(X)\in K[X]$, dann l\"asst sich
$\alpha$ durch die vier Grundrechenarten und Wurzelziehen genau dann
ausdr\"ucken, wenn $K(\alpha)|K$ aufl\"osbar durch Radikale ist.

Das wichtigste Ergebnis in diesem Zusammenhang ist.
\begin{Satz}
  Sei $F|K$ eine endliche Erweiterung von K\"orpern der Charakteristik $0$,
  und $L$ eine Galoish\"ulle von $F|K$. Dann ist $F|K$ genau dann durch
  Radikale aufl\"osbar, wenn die Galoisgruppe $\Gal(L|K)$ aufl\"osbar ist.
\end{Satz}
\begin{proof}[Beweis]
  Sei $G=\Gal(L|K)$ aufl\"osbar, und $1=N_r\triangleleft
  N_{r-1}\triangleleft\dots\triangleleft N_0=G$ eine Kompositionsreihe. Wegen
  der Aufl\"osbarkeit von $G$ ist $N_{i}/N_{i+1}$ zyklisch von
  Primzahlordnung. Sei $E_i$ der Fixk\"orper von $N_i$. Dann ist
  $E_{i+1}|E_{i}$ galoissch mit Galoisgruppe $N_{i}/N_{i+1}$, also zyklisch.
  Sei $n=[L:K]$ und $\zeta$ eine primitive $n$-te Einheitswurzel in einem
  Erweiterungsk\"orper von $L$. Die Galoisgruppe von
  $E_{i+1}(\zeta)|E_{i}(\zeta)$ ist (via Einschr\"ankung auf $E_{i+1}$
  isomorph zu einer Untergruppe der Galoisgruppe von $E_{i+1}|E_{i}$. Sei
  $m_i=[E_{i+1}(\zeta):E_{i}(\zeta)]$. Da $m_i$ ein Teiler von
  $[E_{i+1}:E_{i}]$ ist, ist $m_i$ ein Teiler von $n$. Daher enth\"alt
  $E_i(\zeta)$ die primitiven $m_i$-ten Einheitswurzeln. Wegen Satz
  \ref{S:zykErw} folgt, dass $F|K$ aufl\"osbar durch Radikale ist.
  
  Sei nun umgekehrt $F|K$ aufl\"osbar durch Radikale, und die K\"orper $E_i$
  wie in der Definition. Sei $n$ das Produkt der $m_i$ in der Bedeutung von
  oben, und $\zeta$ eine primitive $n$-te Einheitswurzel in einer Erweiterung
  von $E_r$. Wir betrachten die K\"orperkette $F_0\subseteq
  F_1\subseteq\dots\subseteq F_{r+1}$ mit $F_0=K$, und $F_i=E_{i-1}(\zeta)$
  f\"ur $i\ge1$. Sei $E$ eine Galoish\"ulle von $F_{r+1}|K$. Setze
  $A=\Gal(E|K)$, und sei $U_i$ die Fixgruppe von $F_i$. Wegen Lemma
  \ref{L:Kz:K} ist $F_1|F_0=K(\zeta)|K$ galoissch mit abelscher Galoisgruppe.
  Nun studieren wir die Erweiterung $F_{i+1}|F_i$. Sei $\alpha\in E_i$ und
  $m\in\NNN$ mit $\alpha^m\in E_{i-1}$ und $E_i=E_{i-1}(\alpha)$. Dann gilt
  $F_{i+1}=F_i(\alpha)$ mit $\alpha^m\in F_i$. Da $m$ ein Teiler von $n$ ist
  enth\"alt $F_i$ die $m$-ten Einheitswurzeln, und daher ist $F_{i+1}$ der
  Zerf\"allungsk\"orper von $X^m-\alpha^m$ \"uber $F_i$. Die Galoisgruppe
  dieser Erweiterung fixiert die $m$-ten Einheitwurzeln, und bildet $\alpha$
  auf $\rho\alpha$ ab, wobei $\rho$ eine $m$-te Einheitswurzel ist. So sieht
  man, dass diese Galoiserweiterung abelsch ist.

Es ist also $U_i/U_{i+1}$ abelsch f\"ur alle $i=0,1,\dots,r$. Da $E$ die
Galoish\"ulle von $F_{r+1}|K$ ist, enth\"alt $U_{r+1}$ keinen Normalteiler
$>1$ von $A$. Der Schnitt der Konjugierten von $U_{r+1}$ in $A$ ist also
trivial. Die Behauptung folgt nun aus dem folgenden rein gruppentheoretischen
Lemma.
\end{proof}
\begin{Lemma}
  Sei $G$ eine endliche Gruppe mit einer Kette von Untergruppen
\[
G=U_0\ge U_1\ge \dots\ge U_n
\]
mit folgender Eigenschaft:
\begin{itemize}
\item[(a)] F\"ur $i=0,1,\dots,n-1$ ist $U_{i+1}$ normal in $U_i$, und
  $U_i/U_{i+1}$ ist aufl\"osbar.
\item[(b)] Der Schnitt der Konugierten $U_n^g$, $g\in G$, ist trivial.
\end{itemize}
Dann ist $G$ aufl\"osbar.
\end{Lemma}
\begin{proof}[Beweis]
  Wir beweisen die Aussage durch vollst\"andige Induktion \"uber $n$. F\"ur
  $n=0$ ist nichts zu zeigen. Sei $N:=\bigcap_{g\in G}U_{n-1}^g$. Dann
  erf\"ullt die um $1$ k\"urzere Kette $G/N>U_1/N>\dots U_{n-1}/N$ die
  Voraussetzungen, und daher ist nach Induktionsannahme $G/N$ aufl\"osbar. Wir
  m\"ussen also noch zeigen, dass $N$ aufl\"osbar ist. Da $U_{n-1}/U_n$
  aufl\"osbar ist, gibt es ein $k$, so dass f\"ur die $k$-te Kommutatorgruppe
  $U_{n-1}^{(k)}\le U_n$ gilt. Wegen $N\le U_{n-1}^g$ f\"ur alle $g\in G$
  folgt $N^{(k)}\le (U_{n}^{(k)})^g\le U_n^g$ f\"ur alle $g$. Da die $U_n^g$
  sich aber trivial schneiden, folgt $N^{(k)}=\{1\}$, und $N$ ist aufl\"osbar.
\end{proof}
L\"osungsformeln f\"ur Polynome vom Grad $\le4$ waren schon im sp\"aten
Mittelalter bekannt. Erst Abel konnte zeigen, dass es f\"ur Polynome vom Grad
$5$ im Allgemeinen keine L\"osungsformel geben kann. Vorher ben\"otigen wir
zwei einfache Aussagen.
\begin{Lemma}
  Sei $p$ eine Primzahl, und $G\le S_p$ eine transitive Permutationsgruppe,
  die eine Transposition enth\"alt. Dann ist $G=S_p$.
\end{Lemma}
\begin{proof}[Beweis]
  Wir nennen zwei Elemente $a,b\in\{1,2,\dots,p\}$ \"aquivalent, wenn entweder
  $a=b$ gilt, oder die Transposition $(a,b)$ in $G$ liegt. Dies definiert eine
  \"Aquivalanzrelation auf $\{1,2,\dots,p\}$: Symmetrie und Reflexivit\"at
  sind klar. Bleibt die Transitivit\"at zu zeigen: Seien $(a,b),(b,c)\in G$
  mit $a\ne c$. Die Behauptung folgt wegen $(a,b)(b,c)(a,b)=(a,c)$. Wegen
  $(a,b)^g=(a^g,b^g)$ sind $a$ und $b$ genau dann \"aquivalent, wenn $a^g$ und
  $b^g$ \"aquivalent sind. Die Gruppe $G$ permutiert also die
  \"Aquivalenzklassen. Da $G$ transitiv operiert, haben alle
  \"Aquivalenzklassen die gleiche Gr\"o\ss e. Da $G$ eine Transposition
  enth\"alt, hat eine \"Aquivalenzklasse mindestens $2$ Elemente. Andereseits
  ist die Gr\"o\ss e einer \"Aquivalenzklasse ein Teiler von $p$, also gleich
  $p$. Das hei\ss t, $G$ enth\"alt alle Transpositionen, und daraus folgt
  $G=S_p$.
\end{proof}
\begin{Lemma}
  Das irreduzible Polynom $f(X)\in\QQQ[X]$ vom Primzahlgrad $p$ habe genau
  $p-2$ reelle Nullstellen. Sei $L$ ein Zerf\"allungsk\"orper von $f(X)$
  \"uber $\QQQ$. Dann gilt $\Gal(L|\QQQ)=S_p$.
\end{Lemma}
\begin{proof}[Beweis]
  Die Galoisgruppe operiert transitiv auf den Nullstellen. Ohne
  Einschr\"ankung gelte $L\subset\CCC$. Dann ist $\Gal(L\RRR|\RRR)$ isomorph
  zu einer Untergruppe von $\Gal(L|\RRR)$. Da $f$ genau $p-2$ reelle
  Nullstellen hat, fixiert $\Gal(L\RRR|\RRR)$ genau $p-2$ Nullstellen, und
  vertauscht die beiden komplex konjugierten Nullstellen. Somit ist
  $\Gal(L\RRR|\RRR)$ eine von einer Transposition erzeugte Gruppe, und die
  Behauptung folgt aus obigem Lemma.
\end{proof}
Sei nun etwa $f(X)=X^5-80X-2$. Dann ist $f$ nach dem Eisenstein-Kriterium
\"uber $\QQQ$ irreduzibel. Eine einfache Kurvendiskussion zeigt, dass $f$
genau $3$ reelle Nullstellen hat. Daher hat der Zerf\"allungsk\"orper $L$ von
$f$ die Galoisgruppe $S_5$. Sei $\alpha\in L$ eine Nullstelle von $f$. Dann
ist $L$ die Galoish\"ulle von $\QQQ(\alpha)|\QQQ$. Da $S_5$ nicht aufl\"osbar
ist, l\"asst sich $\alpha$ nicht durch Radikale ausdr\"ucken.
\subsubsection{Zirkel- und Linealkonstruktion regul\"arer $n$-Ecke}
Wir hatten schon fr\"uher gesehen, dass man z.B.\ das regul\"are $7$-Eck mit
Zirkel und Lineal nicht konstruieren kann. Mit Hilfe der Galoistheorie
k\"onnen wir nun ein Kriterium angeben, f\"ur welche regul\"are $n$-Ecke
konstruierbar sind.
\begin{Lemma}
  Sei $2\le n\in\NNN$, und $\zeta\in\CCC$ eine primitive $n$-te
  Einheitswurzel. Dann ist das regul\"are $n$-Eck genau dann mit Zirkel und
  Lineal konstruierbar, wenn $[\QQQ(\zeta):\QQQ]$ eine $2$-er Potenz ist.
\end{Lemma}
\begin{proof}[Beweis]
  Die eine Beweisrichtung folgt aus Lemma \ref{L:nEck}. Sei nun also
  $[\QQQ(\zeta):\QQQ]$ eine $2$-er Potenz. Setze
  $\alpha=\zeta+\frac{1}{\zeta}=2\Re(\zeta)$. Die Galoisgruppe von
  $\QQQ(\zeta)|\QQQ$ ist abelsch, daher hat wegen $\alpha\in\QQQ(\zeta)$ die
  Erweiterung $\QQQ(\alpha)|\QQQ$ eine abelsche $2$-Gruppe $G$ als
  Galoisgruppe. Daher gibt es eine Kette von Gruppen $1=G_0<G_1<\dots<G_r=G$
  mit $[G_{i+1}:G_i]=2$. Die zugeh\"orige Kette der Fixk\"orper besteht aus
  einer Folge quadratischer K\"orpererweiterungen von $\QQQ$ bis
  $\QQQ(\alpha)$. Daher ist $\alpha$ mit Zirkel und Lineal konstruierbar, und
  damit auch $\zeta$ und schlie\ss lich das regul\"are $n$-Eck.
\end{proof}
Es ist sehr einfach festzustellen, wann $[\QQQ(\zeta):\QQQ]$ eine $2$-er
Potenz ist. Sei $n=\prod p_i^{a_i}$ die Primfaktorzerlegung von $n$. 
Nach Korollar \ref{K:zetaphi} gilt
$[\QQQ(\zeta):\QQQ]=\abs{(\ZZZ/n\ZZZ)^\star}$. Der chinesische Restsatz
liefert
\[
\abs{(\ZZZ/n\ZZZ)^\star}=\prod\abs{(\ZZZ/p_i^{a_i}\ZZZ)^\star}=
  \prod(p_i-1)p_i^{a_i-1}.
\]
Diese Zahl ist offenbar genau dann eine Potenz von $2$, wenn f\"ur jeden
ungeraden Primteiler $p$ von $n$ gilt, dass $p^2$ kein Teiler von $n$ ist, und
$p-1$ eine Potenz von $2$ ist. Es ist nicht schwer zu sehen (\"Ubung!), dass
wenn $p-1$ eine Potenz von $2$ ist, dass dann $p=2^{2^m}+1$ f\"ur ein
$m\in\NNN_0$ gilt. Es ist allerdings nicht bekannt, ob es unendlich viele
Primzahlen dieser Form gibt. Wir erhalten also
\begin{Satz}
  Das regul\"are $n$-Eck ist genau dann mit Zirkel und Lineal konstruierbar,
  wenn das folgende gilt: Sei $n=2^au$ mit $u$ ungerade. Dann ist $u$ das
  Produkt verschiedener Primzahlen der Form $2^{2^m}+1$.
\end{Satz}
\begin{comment}
\subsubsection{Polynome vom Grad $3$}
Sei $K$ ein K\"orper der Charakteristik $\ne2,3$, und $f(X)\in K[X]$ ein
separables Polynom vom Grad $3$. Ohne Einschr\"ankung hat $f$ die Form
$X^3+pX+q$. Die Diskriminante von $f$ ist $D=-4p^3-27q^2$.
Dann gilt
\begin{Satz}
  Es gilt $\Gal(f|K)\le A_3$ genau dann, wenn $D$ ein Quadrat in $K$ ist.
\end{Satz}
\begin{proof}[Beweis]
  Seien $\alpha,\beta,\gamma$ die Nullstellen von $f$ in $\tilde K$, also
  $f(X)=(X-\alpha)(X-\beta)(X-\gamma)$. Das Element
  $(\alpha-\beta)(\beta-\gamma)(\gamma-\alpha)$ ist genau dann fest unter
  $G=\Gal(f|K)$, wenn $G$ nur aus geraden Permutationen besteht. Daher gilt
  $G\le A_3$ genau dann, wenn
  $((\alpha-\beta)(\beta-\gamma)(\gamma-\alpha))^2$ ein Quadrat in $K$ ist.
  Wegen $0=\alpha+\beta+\gamma$, $p=\alpha\beta+\beta\gamma+\gamma\alpha$ und
  $q=-\alpha\beta\gamma$ gilt
  \begin{align*}
    D &= -4p^3-27q^2\\
      &= -4(\alpha\beta+\beta\gamma+\gamma\alpha)^3-27(\alpha\beta\gamma)^2\\
      &= 
  \end{align*}
\end{proof}
\subsubsection{Die Galoisgruppe von $X^n-a$}
In vielen Anwendungen (und so ziemlich jedem Staatsexamen!) kommt die
Berechnung der Galoisgruppe von Polynomen der Form $X^n-a$ vor. Sei $K$ ein
K\"orper dessen Charakteristik $0$ ist oder $n$ nicht teilt. Sei $0\ne a\in
K$. Dann ist $X^n-a$ separabel. Sei $\alpha$ eine Nullstelle von $X^n-a$, und
$\zeta$ eine primitive $n$-te Einheitswurzel. Dann ist $L=K(\zeta,\alpha)$ ein
Zerf\"allungsk\"orper von $X^n-a$, also $\Gal(L|K)$ die Galoisgruppe von
$X^n-a$ \"uber $K$. 
\end{comment}
\subsubsection{Endliche K\"orper}
In diesem Abschnitt wollen wir endliche K\"orper untersuchen. Diese spielen
heutzutage auch in Anwendungen der Algebra, etwa der Codierungstheorie und der
Kryptographie, eine wesentliche Rolle. Die Galoistheorie der endlichen
K\"orper ist sehr \"ubersichtlich.
\begin{Lemma}
  Sei $K$ ein endlicher K\"orper. Dann ist die Charakteristik von $K$ eine
  Primzahl $p$ und die M\"achtigkeit von $K$ ist eine $p$-Potenz.
\end{Lemma}
\begin{proof}[Beweis]
  Der Ringhomomorphismus $\ZZZ\to K$ ist nicht injektiv, daher ist die
  Charakteristik $p>0$. Das Bild von $\ZZZ$ in $K$ ist ein K\"orper $\FFF_p$
  mit $p$ Elementen. Wegen der Endlichkeit von $K$ hat insbesondere $K$
  endliche Dimension $n$ \"uber $\FFF_p$, und besteht daher aus $p^n$
  Elementen.
\end{proof}
\begin{Lemma}
  Sei $E|K$ eine Erweiterung endlicher K\"orper vom Grad $n$ und mit
  $q=\abs{K}$. Dann ist $E$ ein Zerf\"allungsk\"orper von $X^{q^n}-X$ \"uber
  $K$, und $E|K$ ist separabel.
\end{Lemma}
\begin{proof}[Beweis]
  Wegen $\abs{E}=q^n$ hat die multiplikative Gruppe $E^\star$ die Ordnung
  $q^n-1$, daher gilt $x^{q^n-1}=1$ f\"ur alle $0\ne x\in E$, also
  $x^{q^n}-x=0$ f\"ur alle $x\in E$. Daher sind die Elemente von $E$ genau die
  Nullstellen von $X^{q^n}-X$, und der erste Teil der Behauptung folgt.
  
  Die Ableitung von $X^{q^n}-X$ ist $q^nX-1=-1$, hat also keine Nullstellen.
  Insbesondere ist $X^{q^n}-X$ separabel, und das gilt dann erst recht f\"ur
  jedes Minimalpolynom eines Elements aus $E$.
\end{proof}
\begin{Lemma}
  Sei $K$ ein endlicher K\"orper, und $\tilde K$ ein algebraischer Abschluss
  von $K$. F\"ur alle $n\in\NNN$ enth\"alt $\tilde K|K$ genau einen
  Zwischenk\"orper $E$ vom Grad $n$ \"uber $K$.
\end{Lemma}
\begin{proof}[Beweis]
  Nach dem vorigen Lemma gibt es h\"ochstens einen solchen Zwischenk\"orper.
  Es bleibt zu zeigen, dass es auch wenigstens einen solchen K\"orper gibt.
  Dazu sei $q=\abs{K}$ und $F\subseteq\tilde K$ die Menge der $q^n$
  Nullstellen von $X^{q^n}-X$. Wegen $(a+b)^{q^n}-(a+b)=a^{q^n}-a+b^{q^n}-b$
  ist $F$ additiv abgeschlossen, und wegen $(ab)^{q^n}=a^{q^n}b^{q^n}$ ist
  $F^\star$ multiplikativ abgeschlossen. Wegen der Endlichkeit von $F$ sind
  $(F,+)$ und $(F^\star,\cdot)$ Gruppen, und $F$ ist ein Teilk\"orper von
  $\tilde K$. Der Beweis des vorigen Lemmas liefert $K\subseteq F$. Die
  Behauptung folgt.
\end{proof}
\begin{Korollar}
  F\"ur jede Primpotenz $q$ gibt es bis auf Isomorphie genau einen K\"orper
  $\FFF_q$ mit $q$ Elementen.
\end{Korollar}
\begin{proof}[Beweis]
  Sei $q=p^n$ f\"ur eine Primzahl $p$. Dann ist nach obigem Lemma und wegen
  der Eindeutigkeit von Zerf\"allungsk\"orpern bis auf Isomorphie jeder
  K\"orper mit $q$ Elementen isomorph zum Zerf\"allungsk\"orper von $X^q-X$
  \"uber $\FFF_p=\ZZZ/p\ZZZ$.
\end{proof}
Dieser Satz hat eine \"uberraschende Konsequenz, die sich mit gr\"o\ss erem
Aufwand auch mittels der eindeutigen Primfaktorzerlegung im Polynomring $K[X]$
beweisen l\"asst.
\begin{Korollar}
  Sei $K$ ein endlicher K\"orper, und $n\in\NNN$. Dann gibt es ein
  irreduzibles Polynom in $K[X]$ vom Grad $n$.
\end{Korollar}
\begin{proof}[Beweis]
  Sei $E|K$ eine Erweiterung vom Grad $n$. Wegen der Separabilit\"at gilt
  $E=K(\alpha)$ f\"ur ein $\alpha\in E$. Das Minimalpolynom von $\alpha$
  \"uber $K$ hat Grad $n$.
\end{proof}
Die Galoistheorie endlicher K\"orper bereitet keine Schwierigkeiten:
\begin{Satz}
  Sei $E|K$ eine Erweiterung endlicher K\"orper mit $q=\abs{K}$. Dann ist
  $E|K$ galoissch mit zyklischer Galoisgruppe. Ein Erzeuger der Galoisgruppe
  ist $x\mapsto x^q$.
\end{Satz}
\begin{proof}[Beweis]
  Sei $n=[E:K]$. Wir sahen, dass $E|K$ normal und separabel ist. Wie oben
  folgt $x^q=x$ f\"ur alle $x\in K$. Es bleibt zu zeigen, dass der angegebene
  Automorphismus von $E$ die Ordnung $n$ hat. Die Ordnung sei $m$. Dann gilt
  $x^{q^m}=x$ f\"ur alle $x\in E$, also $q^n=\abs{E}\le q^m$, also $m\ge n$.
  Andererseits gilt $m\le n$, also $m=n$.
\end{proof}
In Anwendungen (und Staatsexamensklausuren) st\"o\ss t man oft auf die Frage,
wie sich Kreisteilungspolynome \"uber endlichen K\"orpern zerlegen. Ist
$\Phi_n\in\ZZZ[X]$ das $n$-te Kreisteilungspolynom, dann k\"onnen wir die
Koeffizienten modulo $p$ reduzieren und $\Phi_n(X)$ auch als Polynom \"uber
einem K\"orper der Charakteristik $p$ betrachten.
\begin{Satz}
  Die Charakteristik $p$ des endlichen K\"orpers $K$ teile nicht
  $n\in\NNN$. Dann ist jede Nullstelle von $\Phi_n(X)$ in einem
  Erweiterungsk\"orper von $K$ eine primitive $n$-te Einheitswurzel.
  
  Sei $r\in\NNN$ minimal, so dass $n$ ein Teiler von $\abs{K}^r-1$ ist. Dann
  zerf\"allt $\Phi_n$ \"uber $K$ in irreduzible Faktoren vom Grad $r$.
\end{Satz}
\begin{proof}[Beweis]
  Da jede $n$-te Einheitswurzel aus $\CCC$ eine primitive $d$-te
  Einheitswurzel f\"ur einen Teiler $d$ von $n$ ist, gilt
\[
X^n-1=\prod_{d|n}\Phi_d(X).
\]
Sei nun $\Phi_n(\alpha)=0$ f\"ur $\alpha$ algebraisch \"uber $K$, und $d<n$
die multiplikative Ordnung von $\alpha$. Wegen
\[
X^d-1=\prod_{t|d}\Phi_t(X)
\]
gibt es einen Teiler $t<n$ mit $\Phi_t(\alpha)=0$. Wegen
$\Phi_t(\alpha)=0=\Phi_n(\alpha)$ ist dann $\alpha$ eine mehrfache Nullstelle
von $X^n-1$, also eine Nullstelle der Ableitung $nX^{n-1}$. Da geht aber nur,
wenn $p$ ein Teiler von $n$ ist.

Sei nun $f$ ein irreduzibler Faktor von $\Phi_n(X)$ \"uber $K$, und
$q=\abs{K}$. Sei $r$ der Grad von $f$, und $\alpha$ eine Nullstelle von $f$.
Dann ist $K(\alpha)$ ein K\"orper der Ordnung $q^r$, und somit ist $n$ ein
Teiler von $q^r-1$. Sei nun $n$ ein Teiler von $q^s-1$, und $E|K$ eine
Erweiterung von $K$ vom Grad $s$. Da die multiplikative Gruppe von $E$
zyklisch der Ordnung $q^s-1$ ist, enth\"alt sie eine und daher alle primitiven
$n$-ten Einheitswurzeln. Somit zerf\"allt $f$ \"uber $E$ in Linearfaktoren,
und es gilt $K(\alpha)\subseteq E$, also $r\le s$. Die Behauptung folgt.
\end{proof}
\subsubsection{Primzahlen in $m\NNN+1$ und abelsche Galoisgruppen \"uber
  $\QQQ$} Sind $m$ und $r$ teilerfremde nat\"urliche Zahlen, dann besagt der
Dirichletsche Primzahlsatz, dass die arithmetische Folge $km+r$, $k\in\NNN$,
unendlich viele Primzahlen enth\"alt. Der Beweis dieser Aussage ist sehr
schwierig. Allerdings liefert unser Satz \"uber die Zerlegung der
Kreisteilungsk\"orper modulo Primzahlen einen einfachen Beweis f\"ur den
Spezialfall $r=1$, der f\"ur viele Anwendungen ausreicht.
Wir beginnen mit einem einfachen
\begin{Lemma}
  Sei $f(X)\in\ZZZ[X]$ ein Polynom positiven Grades. Dann ist die Menge der
  Primzahlen, die ein Teiler von $f(a)$ f\"ur ein $a\in\ZZZ$ sind, unendlich.
\end{Lemma}
\begin{proof}[Beweis]
  Der Beweis \"ahnelt Euklids Beweis der Unendlichkeit der Primzahlmenge. Wir
  nehmen, dass es nur endlich viele Primteiler von Zahlen der Form $f(a)$
  gibt. Sei $P$ das Produkt dieser Primzahlen. Wenn wir eventuell $f(X)$ durch
  $f(X-c)$ mit $c\in\ZZZ$ ersetzen, d\"urfen wir $f(0)=a\ne0$ annehmen.  F\"ur
  $m\in\NNN$ ist $\frac{1}{a}f(amP)$ offensichtlich ganzzahlig. W\"ahle
  $m\in\NNN$ so, dass $\frac{1}{a}f(amP)\notin\{-1,0,1\}$ gilt. Dann hat
  $\frac{1}{a}f(amP)$ einen Primteiler $p$, der wegen
  $\frac{1}{a}f(amP)\equiv1\pmod{P}$ kein Teiler von $P$ ist, ein Widerspruch.
\end{proof}
\begin{Satz}
  Sei $m\in\NNN$. Dann gibt es unendlich viele Primzahlen $p$ mit
  $p\equiv1\pmod{m}$.
\end{Satz}
\begin{proof}[Beweis]
  Sei $a\in\ZZZ$ und $p$ ein Primteiler von $\Phi_m(a)$, der $m$ nicht teilt.
  $X-a$ ist ein Teiler von $\Phi_m(X)-\Phi_m(a)$, daher ist $\Phi_m(X)$ modulo
  $p$ durch $X-a$ teilbar. Nach obigem Satz zerf\"allt $\Phi_m(X)$ modulo $p$
  in Linearfaktoren, und $m$ ist ein Teiler von $p-1$. Es gilt also
  $p\equiv1\pmod{m}$, und die Behauptung folgt aus dem vorherigen Lemma.
\end{proof}
Eine \"uberraschende Anwendung des Satzes ist die Tatsache, dass jede endliche
abelsche Gruppe als Galoisgruppe \"uber $\QQQ$ vorkommt.
\begin{Satz}
  Sei $G$ eine endliche abelsche Gruppe. Dann gibt es eine Galoiserweiterung
  $E|\QQQ$ mit $\Gal(E|\QQQ)=G$. Hierbei kann sogar $E\subseteq\QQQ(\zeta)$
  f\"ur eine Einheitswurzel $\zeta$ gew\"ahlt werden.
\end{Satz}
\begin{proof}[Beweis]
Sei $m$ die Ordnung von $G$, und $G=\prod_{i=1}^rC_{m_i}$ f\"ur zyklische
Gruppen $C_{m_i}$ der Ordnung $m_i$. Seien $p_1,p_2,\dots,p_r$ verschiedene
Primzahlen mit $p_i\equiv1\pmod{m}$, und $n=\prod_{i=1}^rp_i$. Sei $\zeta$
eine primitive $n$-te Einheitswurzel. Dann ist
\[
\Gal(\QQQ(\zeta)|\QQQ)\cong(\ZZZ/n\ZZZ)^\star\cong\prod_{i=1}^rC_{p_i-1}.
\]
Da $m_i$ ein Teiler von $p_i-1$ ist, gibt es Gruppenepimorphismen
$C_{p_i-1}\to C_{m_i}$. Daher hat $\Gal(\QQQ(\zeta)|\QQQ)$ einen Normalteiler
$N$ mit
\[
\Gal(\QQQ(\zeta)|\QQQ)/N\cong G.
\]
Der Fixk\"orper von $N$ in $\QQQ(\zeta)$ hat also die Galoisgruppe $G$ \"uber
$\QQQ$.
\end{proof}
\subsection{Quadratische Reste}
Bei theoretischen und praktischen Anwendungen st\"o\ss t man oft auf die
Situation, dass man entscheiden muss, wann eine ganze Zahl $m$ ein Quadrat
modulo einer nat\"urlichen Zahl $n$ ist, das hei\ss t, ob es eine ganze Zahl
$k$ gibt mit $n|m-k^2$. Mit Hilfe des chinesischen Restsatzes kann man das
leicht reduzieren auf den Fall, dass $n$ eine Primpotenz ist. Ist $n=q^r$ eine
ungerade Primpotenz, dann kann man durch Induktion beweisen, dass $m$ genau
dann ein Quadrat modulo $q^r$ ist, wenn $m$ ein Quadrat modulo $q$ ist. Der
Fall $n=2^r$ ist nur geringf\"ugig aufw\"andiger. Man kann die Frage also im
wesentlichen darauf reduzieren, dass $n$ eine Primzahl ist. Das motiviert die
folgende Definition.
\begin{Definition*}
  Sei $p\ne2$ eine Primzahl, und $a\in\ZZZ$ nicht durch $p$ teilbar. Man sagt,
  dass $a$ ein \ei{quadratischer Rest} modulo $p$ ist, wenn die Gleichung
  $X^2=a$ in $\FFF_p$ l\"osbar ist. Im anderen Fall ist $a$ ein
  \ei{quadratischer Nichtrest} modulo $p$.
\end{Definition*}
Eine kompakte Notation f\"ur diesen Begriff bildet das
\ei{Legendre-Symbol}. F\"ur $2\ne p\in\PPP$ und $a\in\ZZZ$ setzt man
\[
\left(\frac{a}{p}\right)=
\begin{cases}
  1, & \text{$a$ ist quadratischer Rest modulo $p$}\\
-1,  & \text{$a$ ist quadratischer Nichtrest modulo $p$}\\
0, & \text{$p$ teilt $a$}
\end{cases}
\]
Nat\"urlich h\"angt $\left(\frac{a}{p}\right)$ nur von der Restklasse von $a$
modulo $p$ ab. F\"ur $a\in\FFF_p$ hat daher $\left(\frac{a}{p}\right)$ die
offensichtliche Bedeutung. Der Kern des Homomorphismus
$\FFF_p^\star\to\FFF_p^\star$, $x\mapsto x^2$ ist $\{-1,1\}$. Daher enth\"alt
$\FFF_p^\star$ genau $\frac{p-1}{2}$ quadratische Reste, und diese Elemente
bilden eine Untergruppe vom Index $2$.
\begin{Satz}[Euler-Kriterium]
  Sei $2\ne p\in\PPP$ und $a\in\ZZZ$. Dann gilt
\[
\left(\frac{a}{p}\right)\equiv a^{\frac{p-1}{2}}\pmod{p}.
\]
\end{Satz}
\begin{proof}[Beweis]
  Die Aussage ist klar f\"ur $p|a$. Sei also $p$ kein Teiler von $a$. Sei
  $\left(\frac{a}{p}\right)=1$, und $\bar a$ das Bild von $a$ in
  $\FFF_p$. Daher ist $\bar a$ ein Quadrat $b^2$ in $\FFF_p$. Es folgt
\[
\bar a^{\frac{p-1}{2}}=b^{p-1}=1.
\]
Daher sind genau die $\frac{p-1}{2}$ Quadrate in $\FFF_p^\star$ die
Nullstellen von $X^{\frac{p-1}{2}}-1$. Die Nichtquadrate sind also keine
Nullstellen dieses Polynoms. Sei $a\in\FFF_p^\star$ ein Nichtquadrat. Wegen
$1=a^{p-1}=(a^{\frac{p-1}{2}})^2$ ist daher $a^{\frac{p-1}{2}}=\pm1$, also
$a^{\frac{p-1}{2}}=-1$, und die Behauptung folgt.
\end{proof}
Ein einfaches Korollar ist die Multiplikativit\"at des Legendre-Symbols.
\begin{Lemma}
  Sei $2\ne p\in\PPP$, $a,b\in\ZZZ$. Dann gilt
\[
\left(\frac{ab}{p}\right)=\left(\frac{a}{p}\right)\left(\frac{b}{p}\right).
\]
\end{Lemma}
Daher sind vor allem die Werte $\left(\frac{p}{q}\right)$ f\"ur Primzahlen $p$
und $q$ interessant. Diese wurden schon zu Eulers Zeiten f\"ur gro\ss e Werte
f\"ur $p$ und $q$ berechnet. Dabei stellte er schon 1740 eine merkw\"urdige
Beziehung zwischen $\left(\frac{p}{q}\right)$ und $\left(\frac{q}{p}\right)$
fest. Sind zum Beispiel $p$ und $q$ beide kongruent $1$ modulo $4$, dann fand
er in allen berechneten Beispielen
$\left(\frac{p}{q}\right)=\left(\frac{q}{p}\right)$. Eigentlich sollte man
denken, dass die Frage, ob $p$ ein Quadrat modulo $q$ ist, nichts damit zu tun
hat, ob $q$ ein Quadrat modulo $p$ ist. Euler konnte seine empirischen
Beobachtungen allerdings nicht beweisen. Auch Legendre bemerkte 1785 diese
merkw\"urdigen Beziehungen, er lieferte einen Beweisansatz, der allerdings auf
dem damals noch unbewiesenen Dirichletschen Primzahlsatz beruhte. Erst Gau\ss\
gab ab 1801 acht verschiedene Beweise. Bis heute sind weit \"uber 100 Beweise
bekannt. Die Vielzahl der Beweise r\"uhrt daher, dass man (im wesentlichen
vergeblich) nach einem nat\"urlichen Beweis suchte. Die von Euler, Legendre
und Gau\ss\ gefundene Beobachtung l\"asst sich so zusammenfassen.
\begin{Satz}[Quadratisches Reziprozit\"atsgesetz]
  Seien $p,q$ ungerade und verschiedene Primzahlen. Dann gilt
  \begin{itemize}
  \item[(a)] $\left(\frac{p}{q}\right)\left(\frac{q}{p}\right)=
    (-1)^{\frac{p-1}{2}\frac{q-1}{2}}$
  \item[(b)] $\left(\frac{-1}{p}\right)=(-1)^{\frac{p-1}{2}}$
  \item[(c)] $\left(\frac{2}{p}\right)=(-1)^{\frac{p^2-1}{8}}$
  \end{itemize}
\end{Satz}
Der Beweis dieses Satzes erfordert einige Vorbereitungen.
\begin{Lemma}
  Sei $\zeta$ eine primitive $q$-te Einheitswurzel in einem K\"orper, und
  $j\in\ZZZ$. Dann gilt
\[
\sum_{a=0}^{q-1}\zeta^{ja}=
\begin{cases}
  q, & \text{falls $q|j$}\\
0, & \text{sonst}
\end{cases}
\]
\end{Lemma}
\begin{proof}[Beweis]
  Falls $q|j$, dann gilt $\zeta^j=1$, und die Behauptung ist klar. Sei also
  $q$ kein Teiler von $j$. Dann gilt $\zeta^j\ne1$. Multiplikation mit
  $\zeta^j$ permutiert die Summanden in $S=\sum_{a=0}^{q-1}\zeta^{ja}$, daher
  gilt $\zeta^jS=S$, und somit $S=0$.
\end{proof}
Sei $q$ eine Primzahl, und $\zeta$ eine primitive $q$-te Einheitswurzel in
einem K\"orper. Dann h\"angt $\zeta^a$ nur von der Restklasse von $a$ modulo
$q$ ab. Daher ist $\zeta^a$ f\"ur $a\in\FFF_q$ wohldefiniert, und es gelten
die \"ublichen Potenzregeln.

Der Beweis des quadratischen Reziprozit\"atsgesetzes beruht auf der Berechnung
des Quadrats von Gau\ss-Summen:
\begin{Satz}
  Sei $q\ne2$ eine Primzahl und $\zeta$ eine primitive $q$-te Einheitswurzel in
einem K\"orper. Setze
\[
\tau=\sum_{a\in\FFF_q^\star}\left(\frac{a}{q}\right)\zeta^a.
\]
Dann gilt $\tau^2=(-1)^{\frac{q-1}{2}}q$.
\end{Satz}
\begin{proof}[Beweis]
  Es gilt
  \begin{align*}
    \tau^2 &=
    \left(\sum_{a\in\FFF_q^\star}\left(\frac{a}{q}\right)\zeta^a\right)
    \left(\sum_{b\in\FFF_q^\star}\left(\frac{b}{q}\right)\zeta^b\right)\\
&=
    \sum_{a,b\in\FFF_q^\star}\left(\frac{a}{q}\right)
    \left(\frac{b}{q}\right)\zeta^{a+b}\\
&= \sum_{a,b\in\FFF_q^\star}\left(\frac{ab}{q}\right)
    \zeta^{a+b}
  \end{align*}
  F\"ur festes $a$ durchl\"auft mit $b$ auch $ab$ alle Elemente aus
  $\FFF_q^\star$. Daher \"andert sich die Summe nicht, wenn wir $b$ durch $ab$
  ersetzen. Wir erhalten also
\[
\tau^2 = \sum_{a,b\in\FFF_q^\star}\left(\frac{a^2b}{q}\right)
    \zeta^{a+ab}.
\]
Wegen
\[
\left(\frac{a^2b}{q}\right)=
\left(\frac{a}{q}\right)^2\left(\frac{b}{q}\right)= \left(\frac{b}{q}\right)
\]
folgt weiter
\begin{align*}
  \tau^2 &=
  \sum_{a,b\in\FFF_q^\star}\left(\frac{b}{q}\right)
    \zeta^{a(1+b)}\\
&= \sum_{b\in\FFF_q^\star}\left(\frac{b}{q}\right) \sum_{a\in\FFF_q^\star}
  \zeta^{a(1+b)}\\
&=  \sum_{b\in\FFF_q^\star}\left(\frac{b}{q}\right) (-1+\sum_{a\in\FFF_q}
  \zeta^{a(1+b)})
\end{align*}
Nach dem obigen Lemma verschwindet die Summe \"uber $a$ f\"ur $b\ne-1$, und
hat den Wert $q$ f\"ur $b=-1$. Wir erhalten also
\[
\tau^2=
-\sum_{b\in\FFF_q^\star}\left(\frac{b}{q}\right)+\left(\frac{-1}{q}\right)q.
\]
Da es genauso viele quadratische Reste wie Nichtreste in $\FFF_q^\star$ gibt,
verschwindet die Summe \"uber die $b$, und die Behauptung folgt, da nach dem
Euler-Kriterium $\left(\frac{-1}{q}\right)=(-1)^{\frac{q-1}{2}}$ gilt.
\end{proof}
Zum Beweis des quadratischen Reziprozit\"atsgesetzes seien nun $p$ und $q$
verschiedene ungerade Primzahlen. In einer Erweiterung des K\"orpers $\FFF_p$
w\"ahlen wir eine primitive $q$-te Einheitswurzel $\zeta$. Es sei wieder
\[
\tau=\sum_{a\in\FFF_q^\star}\left(\frac{a}{q}\right)\zeta^a.
\]
Da $p$ ungerade ist gilt
$\left(\frac{a}{q}\right)^p=\left(\frac{a}{q}\right)$. Wir berechnen nun
$\tau^{p-1}$ auf zwei verschiedene Weisen. Da wir in einem K\"orper der
Charakteristik $p$ rechnen, gilt
\begin{align*}
  \tau^p &= (\sum_{a\in\FFF_q^\star}\left(\frac{a}{q}\right)\zeta^a)^p\\
&= \sum_{a\in\FFF_q^\star}\left(\frac{a}{q}\right)\zeta^{ap}\\
&= \sum_{a\in\FFF_q^\star}\left(\frac{ap^2}{q}\right)\zeta^{ap}\\
&= \left(\frac{p}{q}\right)\sum_{a\in\FFF_q^\star}
\left(\frac{ap}{q}\right)\zeta^{ap}
\end{align*}
Wegen $p\ne q$ ist $p$ eine Einheit in $\FFF_q$, mit $a$ durchl\"auft auch
$ap$ die Elemente aus $\FFF_q$. Es folgt also
\begin{align*}
  \tau^p &= \left(\frac{p}{q}\right)\sum_{a\in\FFF_q^\star}
\left(\frac{a}{q}\right)\zeta^{a}\\
&= \left(\frac{p}{q}\right)\tau,
\end{align*}
also
\[
\tau^{p-1}=\left(\frac{p}{q}\right).
\]
Andererseits gilt, nach dem vorigen Lemma und dem Euler-Kriterium,
\begin{align*}
  \tau^{p-1} &=
(\tau^2)^{\frac{p-1}{2}}\\
&= ((-1)^{\frac{q-1}{2}}q)^{\frac{p-1}{2}}\\
&= (-1)^{\frac{p-1}{2}\frac{q-1}{2}}q^{\frac{p-1}{2}}\\
&= (-1)^{\frac{p-1}{2}\frac{q-1}{2}}\left(\frac{q}{p}\right)
\end{align*}
und Teil (a) des Reziprozit\"atsgesetzes folgt. Teil (b) ist das
Euler-Kriterium f\"ur den Fall $a=-1$. Es bleibt (c) zu zeigen. Dazu sei
$\omega$ eine primitive $8$-te Einheitswurzel in einer Erweiterung von
$\FFF_p$. Wegen $0=\omega^8-1=(\omega^4-1)(\omega^4+1)$ gilt $\omega^4=-1$,
also $(\omega+\frac{1}{\omega})^2=2$. Das Euler-Kriterium liefert
\begin{align*}
  \left(\frac{2}{p}\right) &= 2^{\frac{p-1}{2}}\\
&= ((\omega+\frac{1}{\omega})^2)^{\frac{p-1}{2}}\\
&= (\omega+\frac{1}{\omega})^{p-1}\\
&= \frac{(\omega+\frac{1}{\omega})^p}{\omega+\frac{1}{\omega}}\\
&= \frac{\omega^p+\frac{1}{\omega^p}}{\omega+\frac{1}{\omega}}.
\end{align*}
F\"ur $p\equiv\pm1\pmod{8}$ gilt
$\{\omega,\frac{1}{\omega}\}=\{\omega^p,\frac{1}{\omega^p}\}$, also
$\frac{\omega^p+\frac{1}{\omega^p}}{\omega+\frac{1}{\omega}}=1$, und daher
  $\left(\frac{2}{p}\right)=1$.

Im Fall $p\equiv\pm3\pmod{8}$ gilt
$\{\omega^3,\frac{1}{\omega^3}\}=\{\omega^p,\frac{1}{\omega^p}\}$, also
\[
\frac{\omega^p+\frac{1}{\omega^p}}{\omega+\frac{1}{\omega}}=
\frac{\omega^3+\frac{1}{\omega^3}}{\omega+\frac{1}{\omega}}=
\omega^2-1+\frac{1}{\omega^2}=-1,
\]
also $\left(\frac{2}{p}\right)=-1$.

Daher gilt $\left(\frac{2}{p}\right)=1$ genau dann, wenn $p\equiv\pm1\pmod{8}$
gilt. Schreibe $p=8k\pm r$ mit $r\in\{1,3\}$. Dann gilt $p^2-1\equiv
r^2-1\pmod{16}$, somit ist $\frac{p^2-1}{8}$ genau f\"ur $p\equiv\pm1\pmod{8}$
gerade, und die Behauptung folgt.
\paragraph{Ein Beispiel}
Wenn man wissen will, ob etwa $13$ ein Quadrat modulo $3001$ ist, dann muss
man nicht die $1500$ quadratischen Reste in $\FFF_{3001}$ bestimmen und
\"uberpr\"ufen, ob $13$ dabei ist. Das Reziprozit\"atsgesetz liefert sofort
\[
\left(\frac{13}{3001}\right)=
\left(\frac{3001}{13}\right)=
\left(\frac{11}{13}\right)=
\left(\frac{13}{11}\right)=
\left(\frac{2}{11}\right)=
-1.
\]

\subsection{Symmetrische Polynome}
Sei $R[X_1,\dots,X_n]$ der Polynomring in den Variablen $X_1,X_2,\dots,X_n$
\"uber einem kommutativen Ring $R$. Ist $Z$ eine weitere Variable, dann
definiert man f\"ur $1\le i\le n$ das $i$-te \ei{elementarsymmetrische
  Polynom} $S_i$ als den Koeffizienten von $Z^{n-i}$ von
$(Z+X_1)(Z+X_2)\dots(Z+X_n)$, also z.B.\ $S_1=X_1+X_2+\dots+X_n$ und
$S_n=X_1X_2\dots X_n$. Ein \ei{symmetrisches Polynom} ist ein Element aus
$R[X_1,\dots,X_n]$, das sich bei beliebiger Vertauschung seiner Variablen
nicht \"andert. Die Menge der symmetrischen Polynome bildet einen Teilring von
$R[X_1,\dots,X_n]$. Um die Invarianz dieser Elemente unter Operation der
symmetrischen Gruppe auf den $n$ Variablen anzudeuten schreibt man
$R[X_1,\dots,X_n]^{\text{Sym}_n}$. Offensichtlich gilt
$R[S_1,\dots,S_n]\subseteq R[X_1,\dots,X_n]^{\text{Sym}_n}$. Der folgende Satz
besagt, dass tats\"achlich Gleichheit gilt. Z.B.\ gilt $X_1^2+\dots+X_n^2\in
R[X_1,\dots,X_n]^{\text{Sym}_n}$, und wegen $X_1^2+\dots+X_n^2=S_1^2-2S_2$
gilt $X_1^2+\dots+X_n^2\in R[S_1,\dots,S_n]$.
\begin{Satz}
  Es gilt $R[S_1,\dots,S_n]=R[X_1,\dots,X_n]^{\text{Sym}_n}$.
\end{Satz}
\begin{proof}[Beweis]
  F\"ur $\alpha=(\alpha_1,\dots,\alpha_n)\in\NNN_0^n$ schreibe $X^\alpha$
  f\"ur $X_1^{\alpha_1}\dots X_n^{\alpha_n}$. Wir nennen $\alpha$ den
  Multigrad des Monoms $X^\alpha$. Wir ordnen die Multigrade lexikographisch:
  Seien $\alpha$ und $\beta$ verschiedene Multigrade. Sei $i$ minimal mit
  $\alpha_i\ne\beta_i$. Dann gilt $\alpha<\beta$ (bzw.\ $\alpha>\beta$), falls
  $\alpha_i<\beta_i$ (bzw.\ $\alpha_i>\beta_i$). F\"ur $\alpha\ne\beta$ gilt
  also stets $\alpha>\beta$ oder $\alpha<\beta$, und absteigende Ketten von
  Multigraden brechen ab. F\"ur $c\ne0$ sei $\alpha$ der Multigrad des Monoms
  $cX^\alpha$. Den Multigrad eines Polynoms $f\in R[X_1,\dots,X_n]$ definieren
  wir als das Maximum der Multigrade der Monome von $f$, die einen
  Koeffizienten $\ne0$ haben.
  
  Um den Satz zu beweisen gen\"ugt es zu zeigen, dass jedes symmetrische
  Polynom in $R[S_1,\dots,S_n]$ liegt. Dazu gehen wir indirekt vor, und
  w\"ahlen unter allen symmetrischen Polynomen, welche nicht in
  $R[S_1,\dots,S_n]$ liegen, ein solches aus mit minimalem Multigrad
  $\alpha=(\alpha_1,\dots,\alpha_n)$. Da $f$ symmetrisch ist, kommt f\"ur jede
  Permutation $\alpha'=(\alpha_1',\dots,\alpha_n')$ der $\alpha_i$ auch
  $\alpha'$ als Multigrad eines nicht verschwindenden Monoms in $f$ vor. Da
  aber $\alpha$ der gr\"o\ss te in $f$ vorkommende Multigrad ist, gilt
  $\alpha_1\ge\alpha_2\ge\dots\ge\alpha_n$. Wir setzten nun
\[
g=S_1^{\alpha_1-\alpha_2}S_2^{\alpha_2-\alpha_3}\dots
S_{n-1}^{\alpha_{n-1}-\alpha_n}S_n^{\alpha_n}.
\]
Das Monom von h\"ochstem Multigrad in $g$ ist
\[
X_1^{\alpha_1-\alpha_1}(X_1X_2)^{\alpha_2-\alpha_3}\dots (X_1X_2\dots
X_{n-1})^{\alpha_{n-1}-\alpha_n}(X_1X_2\dots X_n)^{\alpha_n}=
X_1^{\alpha_1}\dots X_n^{\alpha_n}=X^\alpha.
\]
Sei $cX^\alpha$ das Monom in $f$ mit h\"ochstem Multigrad. Dann hat $f-cg$
einen kleineren Multigrad als $f$, und liegt daher in $R[S_1,\dots,S_n]$. Aber
auch $cg$ liegt in diesem Ring, und somit auch $f=(f-cg)+cg$, ein Widerspruch.
\end{proof}
\begin{Bemerkung*}
Der angegebene Widerspruchsbeweis liefert einen Algorithmus, um ein
symmetrisches Polynom als Polynom in den elementarsymmetrischen Polynomen zu
schreiben.
\end{Bemerkung*}
Ist $L|K$ eine K\"orpererweiterung, und $T\subseteq L$ eine Teilmenge von $L$,
dann nennt man $T$ \ei{algebraisch unabh\"angig} \"uber $K$, wenn das folgende
gilt: Ist $0\ne f\in K[X_1,\dots,X_n]$ ein Polynom, dann gibt es keine
Elemente $t_i\in T$, $i=1,\dots,n$, mit $f(t_1,t_2,\dots,t_n)=0$. Das ist
offensichtlich \"aquivalent dazu, dass die Menge der Monome
$t_1^{a_1}t_2^{a_2}\dots t_r^{a_r}$, $r\in\NNN_0$, $a_i\in\NNN$, $t_i\in T$
\"uber $K$ linear unabh\"angig ist.

Eine weitere wichtige Eigenschaft der elementarsymmetrischen Polynome ist
\begin{Satz}
  Sei $K$ ein K\"orper. Dann sind die elementarsymmetrischen Polynome
  $S_1,S_2,\dots,S_n$ in den Variablen $X_1,X_2,\dots,X_n$ algebraisch
  unabh\"angig \"uber $K$.
\end{Satz}
\begin{proof}[Beweis]
  Der Multigrad des Monoms $S_1^{\beta_1}S_2^{\beta_2}\dots S_n^{\beta_n}$ ist
  $(\beta_1+\dots+\beta_n,\beta_1+\dots+\beta_n,\dots,\beta_n)$. Daher haben
  verschiedene Monome in den $S_i$ verschiedene Multigrade, und daher kann nur
  eine triviale $K$-Linearkombination solcher Monome verschwinden.
\end{proof}
\begin{Bemerkung*}
  Wegen des letzten Satzes ist der Ring
  $K[X_1,\dots,X_n]^{\text{Sym}_n}=K[S_1,S_2,\dots,S_n]$ isomorph zu einem
  Polynomring in $n$ Variablen. Ist $G$ eine Untergruppe von $S_n$, dann kann
  man $G$ wieder auffasen als Gruppe von Permutationen der Variablen $X_i$,
  und dann via Fortsetzung wieder als Untergruppe der Automorphismengruppe von
  $K[X_1,\dots,X_n]$. Der Fixring $K[X_1,\dots,X_n]^G$ ist allerdings im
  allgemeinen nicht mehr isomorph zu einem Polynomring.
\end{Bemerkung*}
\subsection{Diskriminanten}
In der Literatur gibt es verschiedene Definitionen der Diskriminante eines
Polynoms, die allerdings meist (aber nicht immer!) f\"ur normierte Polynome
das Gleiche ergeben. Die Diskriminante hat in erster Linie zwei Funktionen:
1.\ Mit ihrer Hilfe kann man erkennen ob ein Polynom separabel ist. 2.\ Hat
der K\"orper eine Charakteristik $\ne2$, dann kann man erkennen, ob die
Galoisgruppe eines separablen Polynoms ungerade Permutationen enth\"alt.

Seien $X,t_1,t_2,\dots,t_n$ Variablen \"uber einem K\"orper $K$. Wir fassen
$F(X)=(X-t_1)(X-t_2)\dots(X-t_n)$ als Polynom in der Variablen $X$ auf. Die
Koeffizienten von $F(X)$ sind, bis auf Vorzeichen, die elementarsymmetrischen
Polynome $s_1=t_1+\dots+t_n$, \dots, $s_n=t_1\dots t_n$ in den $t_i$. Das
Element
\[
D=\prod_{1\le i<j\le n}(t_j-t_i)^2
\]
ist offensichtlich ein symmetrisches Polynom in den $t_i$, und l\"asst sich
daher als Polynom in den elementarsymmetrischen Polynomen $s_i$
schreiben. Ersetzt man die $t_i$ durch die Nullstellen eines Grad $n$ Polynoms
$f$, so erhalten wir
\begin{Satz}
  Sei $n\in\NNN$ und $K$ ein K\"orper. Dann gibt es ein Polynom
  $D(Y_1,Y_2,\dots,Y_n)\in K[Y_1,\dots,Y_n]$ mit der folgenden Eigenschaft:
  Sei $f(X)=X^n+a_1X^{n-1}+\dots+a_{n-1}X+a_n\in K[X]$. Dann gilt
\[
D(f):=D(a_1,a_2,\dots,a_n)=\prod_{1\le i<j\le n}(\alpha_j-\alpha_i)^2,
\]
wobei die $\alpha_i$ die Nullstellen von $f$ in einem Zerf\"allungsk\"orper
von $f$ sind.
\end{Satz}
Das Polynom $D$ h\"angt nicht echt von $K$ ab, sondern es gibt ein
universelles Polynom $\tilde D[Y_1,\dots,Y_n]\in\ZZZ[Y_1,\dots,Y_n]$ das nur
von $n$ abh\"angt, dessen Bild in $K[Y_1,\dots,Y_n]$ gerade $D$ ist.

Man nennt $D$ ein \ei{Diskriminanten-Polynom}, und $D(f)=D(a_1,a_2,\dots,a_n)$
die Diskriminante von $f$. Offensichtlich gilt $D(f)\ne0$ genau dann, wenn $f$
separabel ist.
\paragraph{Beispiele}
(a) Sei $n=2$, also etwa $f(X)=X^2+pX+q$. Sind $\alpha$ und $\beta$ die
Nullstellen von $f$, so gilt $\alpha+\beta=-p$, $\alpha\beta=q$, und man
erh\"alt die aus der Schule vertraute Formel
\[
D(f)=(\alpha-\beta)^2=(\alpha+\beta)^2-4\alpha\beta=p^2-4q.
\]

(b) Sei $n=3$. Um ein Polynom vom Grad $3$ auf Separabilit\"at,
Irreduzibilit\"at, Galoisgruppe etc.\ zu untersuchen, k\"onnen wir durch eine
geeignete Substitution von $X$ durch $X+c$ annehmen, dass es die Form
$f(X)=X^3+pX+q$ hat. Seien $\alpha,\beta,\gamma$ die Nullstellen von
$f$. Wegen $\alpha+\beta+\gamma=0$ berechnet man schnell
\begin{align*}
  p &= -\alpha^2-\beta^2-\alpha\beta\\
  q &= (\alpha+\beta)\alpha\beta\\
  D(f) &= (\alpha-\beta)^2(2\alpha^2+2\beta^2+5\alpha\beta)^2.
\end{align*}
Setze $u:=\alpha+\beta$, $v:=\alpha\beta$. Dann gilt
\begin{align*}
  p &= -u^2+v\\
  q &= uv\\
  D(f) &= (u^2-4v)(2u^2+v)^2.
\end{align*}
Eine einfache Umformung liefert nun
\begin{align*}
  D(f) &= (u^2-4v)(2u^2+v)^2\\
       &= -4(v-u^2)^3-27u^2v^2\\
       &= -4p^3-27q^2.
\end{align*}
Das Polynom $X^3+pX+q$ ist also genau dann separabel, wenn $4p^3+27q^2\ne0$
gilt.

Eine weitere wichtige Anwendung der Diskriminante ist
\begin{Satz}
  Sei $f\in K[X]$ ein separables Polynom \"uber einem K\"orper der
  Charakteristik $\ne2$. Dann enth\"alt $\Gal(f|K)$ genau dann eine ungerade
  Permutation, wenn $D(f)$ kein Quadrat in $K$ ist.
\end{Satz}
\begin{proof}[Beweis]
  Seien $\alpha_1,\dots,\alpha_n$ die Nullstellen von $f$. Wir setzen
\[
\Delta:=\prod_{1\le i<j\le n}(\alpha_j-\alpha_i).
\]
Es gilt $\Delta^2=D(f)$. Somit ist $D(f)$ genau dann ein Quadrat in $K$, wenn
$\Delta\in K$ gilt. Wir werden gleich zeigen, dass eine Transposition der
$\alpha_i$ das Vorzeichen von $\Delta$ wechselt. Da die geraden Permutationen
gerade diejenigen sind, die ein Produkt einer geraden Anzahl von
Transpositionen sind, gilt f\"ur $\sigma\in\Gal(f|K)$ genau dann
$\Delta^\sigma=\Delta$, wenn $\sigma$ gerade ist. Somit liegt $\Delta$ genau
dann in $K$, wenn alle $\sigma$ gerade sind. Bleibt zu zeigen, dass eine
Transposition das Vorzeichen von $\Delta$ wechselt. Die Transposition $\tau$
vertausche $\alpha_k$ mit $\alpha_l$. Es gelte $k<l$. Dann gilt
\[
\Delta^\tau = \prod_{1\le i<j\le n}(\alpha_j^\tau-\alpha_i^\tau).
\]
Die Faktoren in $\Delta^\tau$ stimmen bis auf Vorzeichen mit denen in $\Delta$
\"uberein. Man z\"ahlt schnell ab, dass genau $2(l-k+1)+1$ der Faktoren das
Vorzeichen wechseln. Die Behauptung folgt, da $2(l-k+1)+1$ ungerade ist.
\end{proof}
\subsection{Elementare Integrierbarkeit}
Schon in der Schule lernt man eine beliebige explizit angegebene Funktion zu
differenzieren. Auch f\"ur beliebig kompliziert gebaute und verschachtelte
Funktionen gibt es keine prinzipiellen Schwierigkeiten, mit der Additivit\"at
der Ableitung und der Produkt--, Quotienten-- und Kettenregel gelingt das
Differenzieren stets. Bei der umgekehrten Aufgabe, n\"amlich der unbestimmten
Integration, sieht das schon anders aus. Man lernt zwar gewisse Regeln, wie
partielle Integration (Produktregel) und Substitutionsregel (Kettenregel),
aber dennoch gelingt h\"aufig eine explizite Integration nicht unter
Verwendung der \"ublichen Funktionen. Es sieht so aus, als ob man gezwungen
w\"are, den Funktionenvorrat zu vergr\"o\ss ern. Da man z.B.\ erfolglos
versucht, eine explizite Stammfunktion etwa f\"ur $e^{-x^2}$ anzugeben,
definiert man in der Wahrscheinlichkeitsrechnung einfach die ``neue'' Funktion
$E(x)=\int_{-\infty}^xe^{-t^2}dt$. Dass dies wirklich n\"otig ist konnte
bereits Liouville zeigen, mit Techniken, die ziemlich analog zu denen sind,
die wir bei der Zirkel-- und Linealkonstruktion verwendet haben.

Um die folgenden Betrachtungen auf eine wenigstens halbwegs pr\"azise Form zu
bringen, ben\"otigen wir einige einfache Tatsachen aus der
Funktionentheorie. Aber auch ohne Funktionentheorie sollte das folgende
verst\"andlich sein.
\subsubsection{Holomorphe und meromorphe Funktionen}
Ist $G$ eine offene Teilmenge von $\CCC$, so nennt man eine Funktion
$f:G\to\CCC$ \ei{holomorph}, wenn f\"ur alle $z_0\in G$ der Grenzwert
$\lim_{z\to z_0}\frac{f(z)-f(z_0)}{z-z_0}$ existiert. In diesem Fall nennt man
den Grenzwert die \ei{Ableitung} von $f$ in $z_0$. So sind etwa Polynome oder
die Exponentialfunktion $e^z$ auf ganz $\CCC$ holomorph. Ist $i$ die
imagin\"are Einheit, also $i^2=-1$, so erkennt man anhand der
Potenzreihenentwicklung von $e^{iz}$ die wichtige Beziehung $e^{iz}=\cos
z+i\sin z$. Jede komplexe Zahl $z\ne0$ hat eine eindeutige Darstellung der
Form $r(\cos\phi+i\sin\phi)=re^{i\phi}$ mit $r>0$ und $-\pi\le\phi<\pi$.
Mittels $\log re^{i\phi}=\log r+i\phi$ l\"asst sich der reelle Logarithmus
nach $\CCC\setminus\{0\}$ fortsetzen, und wenn man die negativen reellen
Zahlen herausnimmt, dann ist $\log z$ dort auch holomorph. Wegen Formeln wie
$\frac{e^{iz}+e^{-iz}}{2}=\cos z$ k\"onnen wir die trigonometrischen
Funktionen durch die Exponentialfunktion ausdr\"ucken, und daher die
Umkehrfunktionen der trigonometrischen Funktionen durch die
Logarithmusfunktion.

Eine reell differenzierbare Funktion braucht keine differenzierbare Ableitung
zu haben. Die Holomorphie ist eine wesentlich st\"arkere
Eigenschaft. Insbesondere gilt der wichtige Satz, dass holomorphe Funktionen
holomorphe Ableitungen haben. Die formalen Regeln wie Produktregel und
Kettenregel gelten auch f\"ur holomorphe Funktionen. Insbesondere bildet die
Menge der holomorphen Funktionen auf $G$ einen Ring.

Um richtig arbeiten zu k\"onnen wollen wir nat\"urlich auch dividieren. Dies
f\"uhrt zum Begriff der meromorphen Funktion. Eine Funktion
$f:G\to\CCC\cup\{\infty\}$ auf einer offenen Menge $G$ hei\ss t
\ei{meromorph}, wenn das folgende gilt: Die Menge $B$ der Elemente $a\in G$
mit $f(a)=0$ ist diskret in $G$, und f\"ur alle $b\in B$ gilt $\lim_{z\to
  b}f(z)=\infty$. Ist das der Fall, dann besitzt jeder Punkt aus $G$ eine
Umgebung $U$, so dass die Funktion ein Quotient $f=\frac{u}{v}$ zweier auf $U$
holomorpher Funktionen $u,v$ ist mit $v\ne0$. Die Menge der auf $G$
meromorphen Funktionen bildet einen K\"orper.
\subsubsection{Differentialk\"orper}
Ein Differentialk\"orper ist ein K\"orper $K$ mit einer Abbildung $':K\to K$,
$a\mapsto a'$, die additiv ist und die Produktregel $(ab)'=a'b+ab'$ erf\"ullt.
Die Quotientenregel folgt aus
$a'=(\frac{a}{b}b)'=(\frac{a}{b})'b+(\frac{a}{b})b'$. Per Induktion beweist
man Formeln wie $(a^n)'=na'a^{n-1}$. Ferner gilt
$1'=(1^2)'=2\cdot1'\cdot1=2\cdot1'$, also $1'=0$. Die Elemente $a\in K$ mit
$a'=0$ nennt man \ei{Konstanten}, man rechnet schnell nach, dass sie einen
Teilk\"orper von $K$ bilden.

In positiver Charakteristik erh\"alt man schnell Pathologien, weshalb wir im
folgenden stets voraussetzen, dass die K\"orper die Charakteristik $0$ haben.

In unserem algebraischen Kontext k\"onnen wir nat\"urlich die analytischen
Beschreibungen von Exponential-- und Logarithmusfunktionen nicht
verwenden. Wie wollen wir algebraisch erkennen, dass etwa die Funktion $b(x)$
gleich $e^{a(x)}$ ist? Differenzieren von $b=e^a$ liefert $b'=a'e^a=a'b$. Wir
sagen, dass $b$ ein Exponential von $a$ ist, wenn $a'=\frac{b'}{b}$ gilt, und
sagen in diesem Fall auch, dass $a$ ein Logarithmus von $b$ ist. H\"aufig
ben\"otigen wir die durch vollst\"andige Induktion leicht zu beweisende
Identit\"at der sogenannten logarithmischen Ableitung:
\[
\frac{(a_1a_2\dots a_n)'}{a_1a_2\dots a_n}=
\frac{a_1'}{a_1}+\frac{a_2'}{a_2}+\dots+\frac{a_n'}{a_n}.
\]
Die folgende wichtige Aussage gilt auch f\"ur unendliche algebraische
Erweiterungen. Wir beschr\"anken uns aber auf endliche Erweiterungen.
\begin{Satz}
  Sei $E|K$ eine endliche K\"orpererweiterung und $'$ eine Ableitung auf
  $K$. Dann l\"asst sich $'$ auf genau eine Weise zu einer Ableitung
  auf $E$ fortsetzen.
\end{Satz}
\begin{proof}[Beweis]
  Sei $'$ eine Fortsetzung der Ableitung auf $E$. Sei $\alpha\in E$, und
  $f(X)=a_0+a_1X+\dots+a_nX^n\in K[X]$ das Minimalpolynom von $\alpha$. Wegen
  $f(\alpha)=0$ gilt
\[
0=(f(\alpha))'=\sum_{j=1}^n(a_j'\alpha^j+ja_j\alpha'\alpha^{j-1}).
\]
Diese Gleichung bestimmt $\alpha'$ eindeutig, da
$\sum_{j=1}^nja_j\alpha^{j-1}\ne0$.

Zum Beweis der Existenz einer Fortsetzung der Ableitung auf $E$ definieren wir
zwei Funktionen $D$ und $\delta$ vom Polynomring $K[X]$ in sich. F\"ur
$g(X)=\sum_{i=0}^mb_iX^i$ setze $\delta f:=\sum_{i=0}^mb_i'X^i$ und
$Df:=\sum_{i=0}^mb_iiX^i$. Man rechnet sofort nach, dass $\delta$ und $D$
additiv sind und die Produktregel einer Ableitung erf\"ullen. Sei nun
$E=K(\alpha)$. Obige Ableitung der Gleichung $f(\alpha)=0$ schreibt sich als
$(\delta f)(\alpha)+\alpha'(Df)(\alpha)$. Wir setzen daher
\[
\alpha':=-\frac{(\delta f)(\alpha)}{(Df)(\alpha)}\in E.
\]
Jedes Element $\beta$ ist von der Form $h(\alpha)$ f\"ur ein $h\in K[X]$.
Daher definieren wir
\[
\beta'=h(\alpha)':=(\delta h)(\alpha)+\alpha'(Dh)(\alpha).
\]
Wir m\"ussen zeigen, dass diese Definition wohldefiniert ist, d.h.\ wenn
$h(\alpha)=\bar h(\alpha)$ gilt, dann sollte auch $h(\alpha)'=\bar h(\alpha)'$
gelten. Sei also $h(\alpha)=\bar h(\alpha)$. Dann gibt es ein Polynom $q$ mit
$h-\bar h=qf$. Wegen der Additivit\"at der rechten Seite in obiger Definition
m\"ussen wir also zeigen, dass
\[
(\delta(qf))(\alpha)+\alpha'(D(qf))(\alpha)=0
\]
gilt. Unter Verwendung der Produktregel f\"ur $\delta$, $D$ sowie
$f(\alpha)=0$ und der Definition von $\alpha'$ folgt
\begin{align*}
  (\delta(qf))(\alpha)+\alpha'(D(qf))(\alpha) &=
((\delta q)(\alpha)f(\alpha)+q(\alpha)(\delta f)(\alpha)\\
&\mathrel{\phantom{=}} {}+\alpha'((D q)(\alpha)f(\alpha)+q(\alpha)(D
f)(\alpha))\\
&= q(\alpha)(\delta f)(\alpha)+\alpha'(q(\alpha)(D f)(\alpha))\\
&= 0.
\end{align*}
Damit ist unsere Definition von $'$ auf $E$ wohldefiniert. Mittels der
Additivit\"at und der Produktregel f\"ur $\delta$ und $D$ rechnet man sofort
nach, dass $'$ eine Ableitung auf $E$ ist. Ferner stimmt die Einschr\"ankung
von $'$ auf $K$ mit der urspr\"unglichen Ableitung \"uberein.
\end{proof}
Wir ben\"otigen das folgende
\begin{Korollar}\label{K:diffgal}
  Sei $E|K$ eine endliche Erweiterung von Differentialk\"orpern, und $\sigma$
  ein $K$-Automorphismus von $E$. Dann kommutiert $\sigma$ mit der Ableitung,
  d.h.\ f\"ur alle $a\in E$ gilt ${a'}^\sigma={a^\sigma}'$.
\end{Korollar}
\begin{proof}[Beweis]
  Die Abbildung $a\mapsto ((a^{\sigma^{-1}})')^\sigma$ ist eine Ableitung von
  $E$, die die Ableitung $'$ von $K$ fortsetzt. Wegen des Satzes muss diese
  Ableitung mit der gegebenen Ableitung von $E$ \"ubereinstimmen, und die
  Behauptung folgt.
\end{proof}
Das folgende Lemma ist f\"ur die weitere Entwicklung essentiell.
\begin{Lemma}
  Sei $K(t)|K$ eine transzendente Erweiterung von Differentialk\"orper mit
  Ableitung $'$. Wir nehmen an, dass $K$ und $K(t)$ die gleichen Konstanten
  haben. Dann gilt
  \begin{itemize}
  \item[(a)] Sei $t'\in K$, und $f(t)\in K[t]$ ein Polynom vom Grad $n>0$.
    Dann ist $f(t)'$ ein Polynom in $K[t]$. Der Grad von $f(t)'$ ist $n$, wenn
    der h\"ochste Koeffizient von $f$ keine Konstante ist, und andernfalls
    $n-1$.
  \item[(b)] Sei $\frac{t'}{t}\in K$. Dann gilt f\"ur alle $0\ne a\in K$ und
    $0\ne n\in\ZZZ$
\[
\frac{(at^n)'}{t^n}\in K\setminus\{0\}.
\]
Ferner ist f\"ur jedes Polynom $f(t)\in K[t]$ positiven Grades die Ableitung
$f(t)'$ ein Polynom in $K[t]$ vom gleichen Grad. Dabei ist $f(t)'$ ein
Vielfaches von $f(t)$ genau dann, wenn $f(t)$ ein Monom ist.
  \end{itemize}
\end{Lemma}
\begin{proof}[Beweis]
(a) Sei $t'=b\in K$ und $f(t)=a_0+a_1t+\dots+a_nt^n$. Dann gilt
\[
f(t)'=a_n't^n+(a_{n-1}'+na_nb)t^{n-1}+\dots.
\]
Ist $a_n$ keine Konstante, also $a_n'\ne0$, dann ist $f(t)'$ offensichtlich
ein Polynom vom Grad $n$. Sei nun $a_n$ eine Konstante. Dann hat $f(t)'$ den
Grad h\"ochstens $n-1$. Wir nehmen mal an, dass der Grad echt kleiner
ist. Dann gilt $a_{n-1}'+na_nb=0$. Wegen $a_n'=0$ folgt
$(a_{n-1}+na_nt)'=a_{n-1}'+na_nb=0$. Daher ist $a_{n-1}+na_nt$ eine Konstante,
also nach Voraussetzung in $K$ enthalten. Es folgt $t\in K$ im Widerspruch zur
Transzendenz von $t$ \"uber $K$.

(b) Nun sei $b=\frac{t'}{t}\in K$, sowie $a$ und $n$ wie im Lemma. Es gilt
\[
(at^n)'=a't^n+nat't^{n-1}=(a'+nab)t^n.
\]
Daher gilt $\frac{(at^n)'}{t^n}\in K$. Falls dieser Term verschwindet, dann
gilt $(at^n)'=0$, daher ist $at^n$ eine Konstante, also $at^n\in K$. Das
widerspricht wiederum der Transzendenz von $t$. Sei nun $f(t)\in K[t]$ ein
Polynom positiven Grades. Obige Betrachtung zeigt, dass $f(t)'$ ein Polynom
vom gleichen Grad ist. Seien nun $a_rt^r$ und $a_st^s$ zwei verschiedene
Monome von $f$, und $f(t)'=cf(t)$ mit $0\ne c\in K$. Dann folgt
\[
(a_r'+ra_rb)t^r=ca_rt^s, (a_s'+sa_sb)t^s=ca_st^s.
\]
Dies ergibt
\[
\frac{a_r'}{a_r}+rb-\frac{a_s'}{a_s}-sb=0.
\]
Wegen $b=\frac{t'}{t}$ und der Identit\"at der logarithmischen Ableitung folgt
\[
\frac{(a_ra_s^{-1}t^{r-s})'}{a_ra_s^{-1}t^{r-s}}=0.
\]
Daher ist $a_ra_s^{-1}t^{r-s}$ eine Konstante, also in $K$. Wegen $r-s\ne0$
ist das ein Widerspruch zur Transzendenz von $t$.
\end{proof}
Zur Formulierung des Problems der expliziten Integrierbarkeit definieren wir
sogenannte \ei{elementare Funktionen}. Das sind solche Funktionen, die sich
aus rationalen Funktionen, Exponential-- und Logarithmusfunktionen bilden
lassen. Ferner lassen wir endliche algebraische Erweiterungen zu. So ist z.B.\ 
die Funktion $a(x)=\log(17+e^{\sqrt{x}})$ elementar, sowie auch eine Funktion
$b(x)$, die auf einer geeigneten offenen Teilmenge von $\CCC$ eine L\"osung
von $b(x)^{23}-\log(x)b(x)^7+a(x)$ ist, da $b(x)$ algebraisch \"uber
$\CCC(x,\log(x),a(x))$ ist. Solche elementare Funktionen sind stets meromorph
auf geeigneten offenen Teilmengen von $\CCC$. Eine elementare Funktion $f$
nennen wir \ei{elementar integrierbar}, wenn es eine elementare Funktion $F$
gibt mit $f=F'$. Die folgende Definition abstrahiert, pr\"azisiert und
verallgemeinert das Konzept elementarer Funktionen.
\begin{Definition*}
  Eine Erweiterung $E|K$ von Differentialk\"orpern hei\ss t \ei{elementar},
  wenn es eine Kette $K=E_0\subseteq E_1\subseteq\dots\subseteq E_r=E$ von
  Differentialzwischenk\"orpern gibt mit $E_i=E_{i-1}(t_i)$, so dass $t_i$
  algebraisch \"uber $E_{i-1}$ ist, oder ein Exponential oder Logarithmus
  eines Elements aus $E_{i-1}$ ist.
\end{Definition*}
\subsubsection{Die S\"atze von Liouville}
Die Ergebnisse dieses Abschnitts gehen im wesentlichen auf Liouville
zur\"uck. Er verwendete noch weitgehend analytische Methoden, im speziellen
Kontext meromorpher Funktionen. Wir folgen hier einem modernen, auf Rosenlicht
zur\"uckgehenden algebraischen Zugang.

Das Hauptergebnis dieses Abschnitts ist
\begin{Satz}[Liouville]
  Sei $K$ ein Differentialk\"orper und $\alpha\in K$. Falls es in einer
  elementaren Differentialk\"orpererweiterung $E$ ein Element $\beta$ mit
  $\beta'=\alpha$ gibt, und die Konstantenk\"orper von $E$ und $K$
  \"ubereinstimmen, dann gibt es Konstanten $c_i\in K$ und Elemente $u_i,v\in
  K$ mit $\alpha=\sum_{i=1}^nc_i\frac{u_i'}{u_i}+v'$.
\end{Satz}
\begin{Bemerkung*}
  Man kann umgekehrt zeigen, dass ein $\alpha$ der angegebenen Form eine
  Stammfunktion in einer elementaren Erweiterung von $K$ hat. Im Kontext
  meromomorpher Funktionen ist das sowieso klar, da
\[
\sum_{i=1}^nc_i\log(u_i)+v
\]
eine Stammfunktion von $\alpha$ auf einem geeigneten Bereich in $\CCC$ ist.
\end{Bemerkung*}
\begin{proof}[Beweis des Satzes]
  Sei $K=E_0\subseteq E_1\subseteq\dots\subseteq E_r=E$ eine Kette wie in der
  Definition elementarer Erweiterungen, mit $\beta\in E_r$. Wir beweisen die
  Aussage durch vollst\"andige Induktion \"uber $r$. Der Fall $r=0$, d.h.\
  $E=K$ ist trivial. Sei also $r\ge1$, und die Aussage korrekt f\"ur Ketten
  der L\"ange $r-1$. Sei $E_1=K(t)$, wo $t$ algebraisch, ein Exponential oder
  ein Logarithmus eines Elements aus $K$ ist. Die Aussage f\"ur $r-1$,
  angewandt auf die Erweiterung $E(t)|K$, liefert $U_i,V\in K(t)$ und
  Konstanten $c_i\in K$ mit
\[
\alpha=\sum_{i=1}^nc_i\frac{U_i'}{U_i}+V'.
\]
Am einfachsten ist der Fall, dass $t$ algebraisch \"uber $K$ ist. Dann gibt es
Polynome $P_i,Q\in K[t]$ mit $U_i=P_i(t)$ und $V=Q(t)$. Es seien
$t=t_1,t_2,\dots,t_m$ die Konjugierten von $t$ \"uber $K$, d.h.\ die $t_i$
sind die Nullstellen des Minimalpolynoms von $t$ \"uber $K$. Sei $L$ der von
$K$ und den $t_i$ erzeugte K\"orper, und $'$ eine Fortsetzung der Ableitung
auf $L$. Da die Ableitung $'$ mit $\Gal(L|K)$ kommutiert, und $\Gal(L|K)$ die
$t_j$ transitiv permutiert, folgt
\[
\alpha=\sum_{i=1}^nc_i\frac{P_i(t_j)'}{P_i(t_j)}+Q(t_j)'
\]
f\"ur alle $j$. Sei $s$ die Anzahl der $t_j$. Summation \"uber die $j$ und
Verwendung der Identit\"at der logaritmischen Ableitung liefert
\[
\alpha=\sum_{i=1}^n\frac{c_i}{s}\frac{(P_i(t_1)P_i(t_1)\dots P_i(t_s))'}
{P_i(t_1)P_i(t_1)\dots P_i(t_s)}+(\frac{Q(t_1)+Q(t_2)+\dots+Q(t_s)}{s})'.
\]
Jedes der Elemente $P_i(t_1)P_i(t_1)\dots P_i(t_s)$ und
$Q(t_1)+Q(t_2)+\dots+Q(t_s)$ wird von $\Gal(L|K)$ fixiert, und liegt damit in
$K$. Somit haben wir eine Darstellung der verlangten Art gefunden.

Etwas aufw\"andiger ist der Fall, dass $t$ transzendent \"uber $K$ ist. Hier
gibt es nun rationale Funktionen $P_i,Q\in K(t)$ mit $U_i=P_i(t)$ und
$V=Q(t)$, also
\begin{equation}\label{PQ}
\alpha=\sum_{i=1}^nc_i\frac{P_i(t)'}{P_i(t)}+Q(t)'.
\end{equation}
Da $K[t]$ faktoriell ist, gibt es f\"ur $P=P_1$ eine Darstellung der Form
$P(t)=a\prod_{i=1}^kR_i(t)^{e_i}$, mit $0\ne a\in F$, $0\ne e_i\in\ZZZ$, und
$R_i(t)$ ist ein irreduzibles, normiertes Polynom in $K[t]$. Die Identit\"at
der logarithmischen Ableitung liefert
\[
\frac{P(t)'}{P(t)}=\frac{a'}{a}+\sum_{i=1}^k e_i\frac{R_i(t)'}{R_i(t)}.
\]
Das machen wir entsprechend f\"ur alle $P_i$. Daher d\"urfen wir in
der Darstellung \eqref{PQ} von $\alpha$ annehmen, dass alle $P_i$ in
$K$ liegen oder normierte und irreduzible Polynome in $K[t]$ sind. Ferner
seien alle $c_i\ne0$ und die $P_i$ paarweise verschieden. Wir wollen
zun\"achst sehen, dass alle $P_i$ in $K$ liegen. Dazu nehmen wir an, das sei
f\"ur ein Element $P=P_i$ nicht der Fall. Dann ist $P(t)\in K[t]$ ein
irreduzibles, normiertes Polynom. Wir behaupten, dass $P(t)'$ ein zu $P(t)$
teilerfremdes Polynom aus $K[t]$ ist. Dazu sei
$P(t)=t^n+a_{n-1}t^{n-1}+\dots+a_1t+a_0$. Dann ist
\[
P(t)'=nt't^{n-1}+\sum_{k=0}^{n-1}(a_k't^k+a_kkt't^{k-1}).
\]
Ist $t$ ein Logarithmus eines Elements aus $K$, dann gilt $t'\in K$, und
$P(t)'\in K[t]$ hat einen um $1$ kleineren Grad als $P(t)$. Wegen der
Irreduzibilit\"at von $P(t)$ folgt die Teilerfremdheit mit $P(t)'$. Ist
hingegen $t$ ein Exponential, also $\frac{t'}{t}=a'$ f\"ur ein $a\in K$, dann
ist $P(t)$ nach dem Lemma ein Monom. Da aber $P(t)$ normiert und irreduzibel
ist, folgt $P(t)=t$, also $\frac{P(t)'}{P(t)}=\frac{t'}{t}=a'$. Daher kann der
Summand $c_i\frac{P(t)'}{P(t)}$ weggelassen werden, und stattdessen $Q(t)'$
durch $(Q(t)+c_ia)'$ ersetzt werden. Ist also $P(t)=P_i(t)$ ein Polynom, so
ist es teilerfremd zu $P(t)'$. Wegen der Teilerfremdheit der $P_i$ ist $P$ ein
Teiler des Nenners von $\alpha-\sum_{i=1}^nc_i\frac{P_i(t)'}{P_i(t)}$, also
auch des Nenners von $Q(t)'$. Sei nun $Q(t)=\frac{R(t)}{S(t)}$ eine gek\"urzte
Darstellung mit $R(t),S(t)\in K[t]$ und $S(t)$ normiert. Beachte, dass
$\alpha$ ein Element aus $K$ ist, und die Ableitung $F(t)'$ eines Polynoms aus
$K[t]$ wegen $t'\in K$ oder $\frac{t'}{t}\in K$ wieder ein Polynom ist. Wegen
$Q'=\frac{SR'-RS'}{S^2}$ folgt, dass $P$ ein Teiler von $S^2$ ist. Da $P$
irreduzibel ist, ist auch $S$ durch $P$ teilbar. Sei $r$ die h\"ochste Potenz
von $P$, die $S$ teilt. Dann gilt $S=P^rM$ mit $P$ und $M$ teilerfremd. Als
Teiler von $S$ ist $P$ auch teilerfremd zu $R$. Wegen der Teilerfremdheit von
$P$ und $P'$ ist schlie\ss lich $P$ teilerfremd zu $RP'M$.  Wegen
\[
Q'=\frac{PMR'-R(rP'M+PM')}{P^{r+1}M}
\]
ist der Nenner von $Q'$ in der gek\"urzten Darstellung durch $P^{r+1}$, also
insbesondere durch $P^2$ teilbar. Das aber widerspricht Gleichung \eqref{PQ},
da der Hauptnenner von $\alpha-\sum$ nicht durch $P^2$ teilbar ist.

Somit sind alle $P_i$ Elemente aus $K$. Daher ist $Q'=\frac{SR'-RS'}{S^2}\in
K$. \"Ahnlich wie eben sieht man: Ist $F(t)\in K[t]$ ein irreduzibler und
normierter Teiler von $S$, dann ist $F^2$ ein Teiler des Nenners von $Q'$ in
gek\"urzter Darstellung, im Widerspruch zu $Q'\in K$. Daher gilt $S=1$, also
$Q\in K[t]$.

Wir ben\"otigen wieder eine Fallunterscheidung. Ist $t'=\frac{a'}{a}$ mit
$a\in K$, so folgt aus dem Lemma $Q(t)=ct+d$ mit $d\in K$ und $c\in K$ eine
Konstante. Setzen wir $Q'=ct'+d'=c\frac{a'}{a}+d'$ in \eqref{PQ} ein, dann
erhalten wir eine Darstellung der gew\"unschten Form.

Ist hingegen $\frac{t'}{t}=a'$ f\"ur ein $a\in K$, dann ergibt $Q'\in K$
zusammen mit dem Lemma, dass bereits $Q\in K$ gilt, und wir sind fertig.
\end{proof}
Um den Satz nun anzuwenden, ben\"otigen wir ein
\begin{Lemma}
  Sei $g(z)\in\CCC(z)$ eine nicht konstante rationale Funktion. Dann ist
  $e^{g(z)}$ transzendent \"uber $\CCC(z)$.
\end{Lemma}
\begin{proof}[Beweis]
  Sei $f(X)=X^n+a_{n-1}X^{n-1}+\dots+a_1X+a_0$ mit $a_i\in\CCC(z)$ das
  Minimalpolynom von $e^{g(z)}$ \"uber $\CCC(z)$. Differenzieren von
  $0=fe^{g(z)}$ liefert
\[
ng'e^{ng}+\sum_{k=0}^{n-1}(a_k'e^{kg}+a_kkg'e^{kg}).
\]
Wegen der Eindeutigkeit des Minimalpolynoms folgt
\[
ng'a_k=(a_k'+a_kkg')
\]
f\"ur alle $k=0,1,\dots,n-1$. Wegen $a=a_0\ne0$ folgt speziell
\[
ng'=\frac{a'}{a}.
\]
Mittels der Identit\"at der logarithmischen Ableitung erkennt man, dass der
Nenner von $\frac{a'}{a}$ nur einfache Nullstellen hat. Mit einer \"Uberlegung
wie oben folgt aber, dass jede Nullstelle des Nenners von $g$ mindestens
doppelt vorkommt. Daher muss $g$ und somit auch $g'$ ein Polynom sein. Aber
$\frac{a'}{a}$ kann nat\"urlich kein Polynom sein.
\end{proof}
Eine erste Anwendung ist
\begin{Satz}
  Seien $f,g\in\CCC(z)$ rationale Funktionen mit $f\ne0$ und $g'\ne0$. Dann
  ist $f(z)e^{g(z)}$ genau dann elementar integrierbar, wenn es eine rationale
  Funktion $a(z)\in\CCC(z)$ gibt mit $f=a'+ag'$.
\end{Satz}
\begin{proof}[Beweis]
  Hat $f$ die angegebene Form, dann ist $ae^g$ wegen
  $(ae^g)'=a'e^g+ag'e^g=fe^g$ eine elementare Stammfunktion von $fe^g$.
  
  Mit unserem Hauptsatz beweisen wir nun die Umkehrung. Wir setzen
  $t=e^{g(z)}$. Wegen des obigen Lemmas ist $t$ transzendent \"uber $\CCC(z)$.
  Wir betrachten den Differentialk\"orper $K=\CCC(t,z)$. Setze
  $\alpha=f(z)e^{g(z)}=ft$. Falls $\alpha$ elementar integrierbar ist, dann
  hat nach dem Hauptsatz $\alpha$ die Form
  $\alpha=\sum_{i=1}^nc_i\frac{u_i'}{u_i}+v'$ mit $c_i\in\CCC$, $u_i,v\in K$.
  Wir fassen die Elemente $u_i,v$ als rationale Funktionen in $t$ mit
  Koeffizientenk\"orper $\CCC(z)$ auf. Es folgt
\[
ft=\sum_{i=1}^nc_i\frac{u_i(t)'}{u_i(t)}+v(t)'.
\]
Wie beim Beweis des Hauptsatzes folgt, dass wir voraussetzen d\"urfen, dass
die $u_i(t)$ schon alle in $\CCC(z)$ liegen und $v(t)$ ein Polynom in $t$ ist.
Wegen $t'=g't$ und $g'\in\CCC(z)$ hat $v(t)'$ den gleichen Grad wie $v(t)$,
und dieser Grad ist gleich dem Grad von $ft$, also $1$. Es folgt $v(t)=b+at$
mit $a,b\in\CCC(z)$, also $v(t)'=b'+a't+at'=b'+a't+ag't$, und somit
$f=a'+ag'$, und die Behauptung folgt.
\end{proof}
\begin{Korollar}[Liouville]
  Die Funktionen $e^{-x^2}$, $\frac{e^x}{x}$, $\frac{1}{\log(x)}$ und
  $\frac{\sin(x)}{x}$ sind nicht elementar integrierbar.
\end{Korollar}
\begin{proof}[Beweis]
  Der erste Fall f\"uhrt auf die Gleichung
\[
1=a'-2xa
\]
f\"ur $a\in\CCC(x)$. Wir m\"ussen sehen, dass es eine solche rationale
Funktion $a$ nicht gibt. Hat $a$ einen Pol $\alpha\ne0$ der Ordnung $r>0$,
dann ist $\alpha$ ein Pol von $a'$ der Ordnung $r+1$, ein Widerspruch. Mit
einem \"ahnlichen Argument ist $0$ kein Pol von $a$. Somit ist $a$ ein
Polynom, ein Widerspruch aus Gradgr\"unden.

Die Funktion $\frac{e^x}{x}$ f\"uhrt auf die Gleichung
\[
\frac{1}{x}=a'+a.
\]
Wie eben folgt durch Betrachtung der Pole, dass diese Gleichung keine L\"osung
hat. Die Substitution $x=\log(y)$ zeigt, dass auch $\frac{1}{\log(y)}$ nicht
elementar integrierbar ist.

Der letzte Fall sei als \"Ubung \"uberlassen. Mit der Substitution $x=iz$
f\"uhrt man das auf die Integration von $\frac{e^z-e^{-z}}{z}$ zur\"uck, und
argumentiert so \"ahnlich wie im Beweis des obigen Satzes.
\end{proof}

\printindex \bibliographystyle{myalpha}
%\bibliographystyle{alpha}
%\bibliographystyle{myplain2}
%\bibliographystyle{amsplain}
\bibliography{mybib}
\end{document}
